% =============================================================================
% Chapter 3 - Literature Review  (reduced)
% =============================================================================
% Produced from contents/chapter3_shortened.tex per CH3_REDUNDANCY_REVIEW.md.
% Include with % =============================================================================
% Chapter 3 - Literature Review  (reduced)
% =============================================================================
% Produced from contents/chapter3_shortened.tex per CH3_REDUNDANCY_REVIEW.md.
% Include with % =============================================================================
% Chapter 3 - Literature Review  (reduced)
% =============================================================================
% Produced from contents/chapter3_shortened.tex per CH3_REDUNDANCY_REVIEW.md.
% Include with % =============================================================================
% Chapter 3 - Literature Review  (reduced)
% =============================================================================
% Produced from contents/chapter3_shortened.tex per CH3_REDUNDANCY_REVIEW.md.
% Include with \include{contents/chapter3_cut} IN PLACE OF chapter3.tex and
% chapter3_shortened.tex. Never include two of them: they share labels.
%
% WHAT CHANGED
%   * Repeated mechanism/derivation passages removed: they are owned by
%     Chapter 2 (magnification, brightness constancy, strain tensor, SimAM
%     energy, focal loss, F1). This chapter now carries published evidence and
%     transfer limits only, and contains NO display mathematics.
%   * The nine per-section "research gap" endings are replaced by short
%     relevance statements, with one qualified synthesis at the end.
%   * Local decision tables (EVM/flow/strain/TIM/CNN/SimAM/imbalance) removed:
%     they restated Chapter 2 or 4 and were then restated again in the ledger.
%   * 17 tables -> 6. 2 figures -> 1. The dataset-composition figure is gone
%     (its content is now one prose sentence); the fold-composition figure is
%     retained because the pooled-versus-fold argument depends on it.
%
% THE LEDGER'S FIRST COLUMN IS NOW \ref{}, NOT HARD-CODED NUMBERS.
% The twelve divergence rows previously cited 3.1.2 ... 3.9.7 literally, which
% any restructuring silently falsifies. They now resolve automatically.
%
% LABELS: every label referenced from introduction.tex, chapter4_shortened.tex,
% chapter5_shortened.tex or chapter6.tex is preserved (24 of them), plus those
% used by the unshortened chapters, so either manuscript state resolves.
% Verified by scratchpad scan; re-run it after any further edit.
%
% NOTE: the old header warned that Chapter 2 cites Chapter 3 subsections BY
% NUMBER. That is no longer true -- the shortened Chapter 2 has none -- so
% subsection positions here are free. The only by-number dependency was the
% ledger's own first column, now removed.
% =============================================================================

\chapter{Literature Review}
\label{ch:literature}

This chapter reviews the literature bearing on each component of the recognition pipeline evaluated in this thesis: the corpus every experiment runs on, the four techniques placed under ablation, the input transformations they operate on, and the training regime held constant across all of them. Each section states what the relevant studies found, what their evidence does not establish, and what follows for this experiment. Mechanisms already derived in \autoref{ch:background} are not derived again. \autoref{sec:synthesis-and-gap} consolidates the closing statements into one research gap and collects every declared departure from the reviewed methods.


% =============================================================================
\section{The Evaluation Corpus: CASME II and the Spontaneous Benchmark Family}
\label{sec:corpus}
% =============================================================================

The corpus is treated first because on a database of this size several decisions usually presented as modelling choices are in fact forced by properties of the data.

\subsection{Why a dedicated spontaneous corpus was necessary}
\label{sec:why-spontaneous}

Automatic expression recognition matured on large corpora of \textit{posed} expressions, reporting accuracies above 90\,\% on the six basic emotions by the early 2010s~\cite{yan2014}. Micro-expression research could not follow, because the phenomenon is involuntary and cannot be performed on instruction. Two problems motivated CASME II directly.

The first is \textbf{construct validity}: Yan et al.~\cite{yan2014} and Li et al.~\cite{liX2013} both observe that deliberately produced ``micro-expressions'' differ systematically from spontaneous ones, so a system trained on posed data may be a recogniser of fast deliberate movement rather than of micro-expressions. The second is \textbf{acquisition quality}. Defined jointly by very short duration and very low intensity, the phenomenon sits close to the noise floor of ordinary video; Yan et al.~\cite{yan2014} argue that CASME's 60~fps coarsely samples a movement lasting a fraction of a second --- while marking the benefit of higher rates as unresolved --- and that mains-frequency flicker contaminated earlier recordings. Elicitation was a third difficulty: deception paradigms produced samples contaminated by conversational movement, and the workable paradigm, which CASME II adopts, has participants watch high-arousal video while attempting to keep the face neutral~\cite{liX2013}.

\begin{table}[htbp]
	\centering
	\footnotesize
	\setlength{\tabcolsep}{4pt}
	\begin{tabular}{l c c c c c c}
		\hline
		\textbf{Corpus} & \textbf{Subj.} & \textbf{Micro samples} & \textbf{fps} & \textbf{Spont.} & \textbf{FACS} & \textbf{Ethn. / age} \\
		\hline
		SMIC~\cite{liX2013}              & 16 valid          & 164                & 100          & Yes          & No           & 3 / 26.7 \\
		\textbf{CASME II}~\cite{yan2014} & \textbf{26 valid} & \textbf{247 / 255} & \textbf{200} & \textbf{Yes} & \textbf{Yes} & \textbf{1 / 22.03} \\
		SAMM~\cite{davison2018}          & 32                & 159                & 200          & Yes          & Yes          & 13 / 33.24 \\
		CAS(ME)$^2$~\cite{qu2016}        & 22                & 53 (+250 macro)    & 30           & Yes          & Yes          & 1 / 22.59 \\
		\hline
	\end{tabular}
	\caption[The spontaneous micro-expression corpora]{The spontaneous corpora that define the current context, compiled from Yan et al.~\cite{yan2014}, Davison et al.~\cite{davison2018}, Li et al.~\cite{liX2013} and Qu et al.~\cite{qu2016}. CASME II carries two sample counts: 247 in the paper, 255 in the released coding.}
	\label{tab:spontaneous-corpora}
\end{table}

\subsection{CASME II: construction and annotation}
\label{sec:casme2-construction}

CASME II~\cite{yan2014} recruited thirty-five participants, mean age 22.03 (SD 1.60), recording at \textbf{200 fps} and 640~$\times$~480 without inter-frame compression under flicker-free illumination, giving a facial region of roughly 280~$\times$~340 pixels. Each retained sample carries onset and offset indices, FACS action-unit labels from two independent coders with disagreements arbitrated by discussion, and an emotion label assigned from the coded action units, participant self-report and the eliciting episode's content --- Yan et al. arguing against forcing spontaneous samples into the six basic categories, since a single action unit is frequently ambiguous.

Two specifics propagate into this thesis. The original release applies a normalisation from which this work departs: 68 Active Shape Model landmarks are detected on the \textit{first frame only} and the resulting transform applied unchanged to every subsequent frame; this pipeline reads the original recorded frames instead, unregistered and uncropped. And the published paper reports \textbf{247} micro-expressions while the distributed coding file contains \textbf{255}, a post-publication revision rarely acknowledged and one reason sample counts for ``CASME II'' vary. The present work uses the released 255-sample coding.

\subsection{The label taxonomy and the sparsity it creates}
\label{sec:label-taxonomy}

Yan et al.~\cite{yan2014} provide five emotion categories with action-unit criteria and are candid about the result: ``some types of facial expressions are difficult to elicit in laboratory situations, thus the samples in different categories distributed unequally.'' The released coding is finer, exposing seven raw labels. Counting directly from the label table used here gives happiness 32, disgust 63, surprise 25, repression 27, sadness 7, fear 2, and others 99.

\textbf{Seven-class classification is not viable on this corpus.} Fear is represented by two clips in the entire database and sadness by seven. Under any subject-disjoint protocol --- and subject-disjointness has been this benchmark's norm since its introduction (\autoref{sec:original-baseline}) --- a two-clip class cannot be meaningfully trained or tested, and per-class metrics degenerate.

The field's response is to group by affective valence. The convention originates with SMIC, released with exactly three classes~\cite{liX2013}, and was carried into cross-corpus work by MEGC 2019~\cite{see2019}. Discarding ``others'' and grouping the remainder gives the working set used throughout: \textbf{Negative 99, Positive 32, Surprise 25, N = 156}, a ratio of roughly 4 : 1.3 : 1.

\textbf{A declared divergence from the MEGC convention.} MEGC 2019 reports the CASME II portion of its composite database as \textbf{145} samples, footnoting that CASME II's Negative class covers ``Disgust and Repression''~\cite{see2019} and building its subset from the paper's five-category release, in which sadness and fear are absorbed into ``others''. This thesis uses the seven-label coding and assigns sadness and fear to Negative, following the mapping MEGC gives for \textbf{SAMM} rather than for CASME II. The consequence is an \textbf{eleven-clip difference}: N = 156 against 145, with Negative at 99 rather than 88. Nine of the eleven are the sadness and fear clips; the remaining two are unaccounted for, since disgust and repression together number 90 in the released coding against the 88 the challenge reports --- a residual noted rather than resolved. The direction is mild: the nine added clips all fall in the majority class and come from the categories coders found hardest to assign, so if anything they make this task marginally harder than the challenge's.

\subsection{The original baseline, and what it does not establish}
\label{sec:original-baseline}

Yan et al.~\cite{yan2014} baseline the corpus with LBP-TOP~\cite{zhaoG2007} and a support vector machine, sweeping spatial and temporal radii and excluding $R_T = 1$ deliberately, because at 200~fps the change between adjacent frames is negligible --- an early acknowledgement that a high sampling rate is not automatically an advantage. Five-class accuracies under \textbf{leave-one-subject-out} cluster in the high fifties to low sixties, the best configuration reaching \textbf{63.41\,\%}. Three observations matter.

\textbf{Leave-one-subject-out is native to this benchmark}, not a stricter protocol adopted later, so any evaluation reporting a single random split departs from the benchmark's own convention. \textbf{The baseline is reported as accuracy alone}, and on a five-class problem where ``others'' is 102 of 247 samples the majority-constant classifier scores 41\,\% without modelling anything, so a headline near 63.41\,\% is not by itself evidence of a working system. \textbf{And the sources disagree on the protocol behind that figure}: Yan et al.~\cite{yan2014} state leave-one-subject-out in abstract and results prose, the caption of the table reporting it reads ``leave-one-out cross-validation'' without the subject qualifier, and Li et al.~\cite{liX2018} reproduce the figure marked as leave-one-\textit{sample}-out, a materially easier protocol. The most widely quoted number on this benchmark thus rests on a protocol its own sources describe three ways; this thesis follows the originating paper's prose and records the discrepancy rather than resolving it.

Accuracy also moves by several points across a purely descriptive radius sweep, with no change to data, split or classifier. Where effect sizes of that magnitude arise from descriptor settings, similar differences between architectures cannot be interpreted without a variance estimate --- a limitation this thesis inherits and declares.

\subsection{Sibling corpora, and why this study remains single-corpus}
\label{sec:sibling-corpora}

Three further databases define the context, each isolating a specific limitation of CASME II. \textbf{SMIC}~\cite{liX2013} is the origin of the three-class taxonomy and of fixed-length temporal interpolation, both inherited here; it is not used as a corpus because it is neither FACS-coded nor apex-annotated. \textbf{SAMM}~\cite{davison2018} addresses the CASME family's demographic narrowness with 13 ethnicities against CASME II's one, Davison et al. observing that a corpus drawn from one ethnicity and one narrow age band restricts analysis to ``similar looking participants''; it is stronger on that count but smaller in usable samples and spread across seven categories. \textbf{CAS(ME)$^2$}~\cite{qu2016} pairs macro- and micro-expression samples from the same participants, making it the natural resource for domain transfer --- the strategy behind the strongest MEGC 2019 entry~\cite{liu2019,see2019} --- but records at 30~fps, supplies only 53 micro-expression samples, and is evaluated leave-one-\textit{video}-out, which does not enforce subject-disjointness.

CASME II is used alone because it has the highest temporal resolution of any spontaneous corpus of comparable size, a precondition for magnification in a defined band and for dense frame-to-frame flow; because its onset, apex and offset annotation makes temporal normalisation well-defined rather than heuristic; and because it remains the most widely reported single-corpus benchmark. The cost --- that no claim here extends to other populations, capture conditions or corpora --- is declared rather than argued away.

\subsection{How CASME II is evaluated: MEGC 2019 and the metric convention}
\label{sec:megc-metrics}

The most consequential development in how CASME II is \textit{used} is the Composite Database Evaluation protocol of the Second Facial Micro-Expression Grand Challenge~\cite{see2019}. It merges SMIC, CASME II and SAMM into one composite database of \textbf{442 samples} --- Negative 250, Positive 109, Surprise 83 --- of which CASME II contributes 145, and runs leave-one-subject-out across the composite. Plain accuracy is rejected explicitly on grounds of class imbalance, and two balanced metrics mandated: \textbf{Unweighted F1 (UF1)}, the macro-averaged F1 defined in \autoref{sec:confusion-matrix}, and \textbf{Unweighted Average Recall}, the balanced accuracy.

What matters is not the algebra but the \textbf{construction}: UF1 accumulates true positives, false positives and false negatives per class \textbf{over all folds} \textit{before} any per-class F1 is computed. As \autoref{sec:pooled-vs-fold} establishes, that is a different estimator from the average of per-fold macro F1 scores, and the distinction is not cosmetic here. Because clips are distributed very unequally across subjects and many single-subject folds contain only one of the three classes, a per-fold macro F1 is bounded well below one by fold composition alone (\autoref{fig:fold-composition}). The primary metric of this thesis, \textit{pooled macro F1}, is therefore identical in construction to UF1; the per-fold average appearing in the raw logs is not a metric the field uses.

\begin{figure}[htbp]
	\centering
	\includegraphics[width=\textwidth]{fig3_2_fold_composition.png}
	\caption[Class composition of the 25 leave-one-subject-out folds]{Class composition of the 25 leave-one-subject-out folds, and the per-fold macro-F1 ceiling it implies.}
	\label{fig:fold-composition}
\end{figure}

\textbf{One structural caveat governs the challenge's reported figures}, which range up to 0.7885 composite UF1 and 0.8764 on the CASME II subset~\cite{see2019}. Every one comes from a model trained on the \textit{composite} database --- roughly 442 samples across three corpora and 68 subjects --- then scored on the CASME II portion of the pooled predictions. A CASME II column is therefore not a CASME II-only result: it reports how a cross-corpus model performs on CASME II clips, with approximately \textbf{2.8 times} the training data available in any single-corpus fold, where the experiment conducted here trains each fold on roughly 150 clips. The two regimes are not interchangeable, and single-corpus results should not be compared against those figures without stating the difference.

Two further observations from the challenge bear on this design. The top three submissions all chose \textbf{optical flow} as their representation, a convergence that supports the motion-based input used here; and one submission substituted mid-position frames for annotated apex frames with little loss, suggesting that precise apex localisation may be less critical than assumed.

\subsection{Threats to validity carried by the corpus}
\label{sec:corpus-threats}

The following are not criticisms of CASME II, which remains the best available resource for this task, but constraints any study built on it inherits. Fold and duration statistics are computed from the 156-clip label table used here.

\textbf{Scale and imbalance.} 156 usable clips across three classes, distributed 99 : 32 : 25, small enough that parameter count is itself a risk factor. This motivates parameter-free or low-parameter components wherever a choice exists, and it makes accuracy unusable alone: the always-Negative classifier attains \textbf{0.6346 accuracy} at a macro F1 of only \textbf{0.2588}, so that floor is quoted alongside every result.

\textbf{Subject-level non-uniformity.} \textbf{Subject 17 alone contributes 33 clips, 21\,\% of the working set}, while three subjects contribute one each, so fold sizes vary by a factor of 33; \textbf{ten of the 25 folds contain a single class}, eight contain two, and only seven all three. This forces full 25-fold leave-one-subject-out and is why the pooled metric rather than the fold average is primary. It is also why a single held-out split is not a reliable instrument here: which subjects fall in the test set can move the result by more than any architectural change under study. Subject 18's three clips are all ``others'' and are removed by the three-class filter, so the run has \textbf{25 folds, not 26} --- stated explicitly because fold counts differ between papers using this corpus.

\textbf{Demographic homogeneity and laboratory ecology.} All participants are Chinese with mean age 22.03, named by Davison et al.~\cite{davison2018} as a principal motivation for SAMM; the MEGC organisers saw the effect in reverse, reporting that adaptation from CK+ helped on SMIC and SAMM but ``not so for CASME II, which contain predominantly Chinese subjects''~\cite{see2019}. Recording is frontal under fixed lamps with participants instructed not to move or speak, so almost none of the head motion, illumination variation, occlusion or speech-related movement a deployed system would face is present. Together these motivate a \textbf{motion} representation rather than raw pixels --- appearance being exactly what a 26-subject corpus would let a network memorise --- and they are why no cross-dataset or in-the-wild claim is made anywhere in this thesis.

\textbf{Duration, annotation and version ambiguity.} Onset-to-offset lengths range from \textbf{31 to 126 frames}, median 66 --- at 200~fps, 0.155\,s to 0.63\,s --- so every clip is resampled to a fixed length before training. Labels derive from coded action units, self-report and stimulus content: a construct produced by trained coders under a rubric, taken here as given and recorded as a declared limitation rather than as measurement error. And three sample counts circulate under the name ``CASME II'' --- 247, 255 and 145 --- with both five- and seven-class label sets in use and rarely stated, so N is reported explicitly alongside every figure here.

\subsection{Implications for this research}
\label{sec:corpus-implications}

Two further properties determine settings rather than threatening validity: the \textbf{200 fps} capture rate fixes the magnification pass-band at a matched \textbf{5--25 Hz}, and the baseline's sensitivity to a descriptive hyper-parameter sets the standard of evidence, so effect sizes below the resolution of a 156-clip corpus are reported as unresolved rather than as findings. No confidence interval is available, since the single-seed design and unretained per-clip predictions preclude one.

The consequence that shapes interpretation rather than construction is that \textbf{the small-sample regime predicts which components should help}. 156 clips cannot supply the data needed to learn a general spatio-temporal feature extractor from scratch, particularly once much of that extraction has been performed analytically before the network sees the data, whereas a component modelling the temporal arc of an expression exploits exactly the structure the corpus's onset, apex and offset annotation guarantees is present. This is a hypothesis the experiment tests, not a result it assumes. What the corpus literature does not supply is a stable basis for comparing methods: only Li et al.~\cite{liX2018} quantify how much of a reported difference is attributable to the protocol rather than the method, and then only for a single matched row of their own method.


% =============================================================================
\section{Motion Magnification: Eulerian Video Magnification}
\label{sec:evm-review}
% =============================================================================

Magnification is reviewed through the four corpus papers that apply it: Bai et al.~\cite{bai2021}, Li et al.~\cite{liX2018}, and Li, Huang and Zhao~\cite{liY2018,liY2021}. The original amplitude- and phase-based formulations lie outside the corpus and are described only as those papers attribute them. The mechanism itself is given in \autoref{sec:motion-magnification}.

The corpus design of \autoref{sec:why-spontaneous} answers the \textit{short-duration} half of the micro-expression definition but not the \textit{low-intensity} half. A movement spread across dozens of frames at 200~fps may still occupy only a few grey levels across a small patch of skin, and Yan et al.~\cite{yan2014} report discarding samples outright because the movement was too faint to place an onset. Li et al.~\cite{liX2018} state the remedy directly --- amplify the motion before any descriptor is computed, so that the difference between categories is enlarged rather than the classifier being made more sensitive. Being entirely upstream of the model is what makes magnification a clean ablation factor here.

Three parameters are free in any application: the temporal pass-band, which decides \textit{which} motions are amplified; the factor $\alpha$, which decides \textit{how much}; and the pyramid depth. On the first, Bai et al.~\cite{bai2021} report the rule given by the original authors --- the band should match the duration of the phenomenon, for micro-expressions 1/5 to 1/25\,s, equivalently \textbf{5--25 Hz} --- while noting that this full band produced ``noticeable noise'' in their setting, which led them to narrow it to 15--25 Hz. This is the only explicit guidance on band selection anywhere in the corpus.

\subsection{Evidence that magnification improves recognition}
\label{sec:evm-evidence}

The strongest case is Li et al.~\cite{liX2018}, who sweep the magnification factor rather than fixing it, across $\alpha \in \{1, 2, 4, 8, 12, 16, 20, 24, 30\}$ with $\alpha = 1$ denoting no magnification, on both SMIC and CASME II under leave-one-subject-out. Two findings emerge.

\textbf{Magnification helps consistently}, ``for all three kinds of features on all four testing datasets'', and on CASME II substantially (\autoref{tab:evm-effect}); their best full configuration reaches 67.21\,\% on CASME II, competitive with the contemporaneous state of the art.

\begin{table}[htbp]
	\centering
	\small
	\begin{tabular}{l c c c}
		\hline
		\textbf{Descriptor} & \textbf{No magnification} & \textbf{With magnification} & \textbf{$\Delta$} \\
		\hline
		LBP  & 55.87\,\% & 60.73\,\% & \textbf{+4.86} \\
		HOG  & 57.49\,\% & 63.97\,\% & \textbf{+6.48} \\
		HIGO & 57.09\,\% & 67.21\,\% & \textbf{+10.12} \\
		\hline
	\end{tabular}
	\caption[Effect of Eulerian motion magnification on CASME II]{Effect of magnification on CASME II under leave-one-subject-out, after Li et al.~\cite{liX2018} (Table 6). This is the corpus's one controlled magnification-on/off comparison.}
	\label{tab:evm-effect}
\end{table}

\textbf{The relationship between $\alpha$ and accuracy is non-monotonic.} Li et al.~\cite{liX2018} describe the curves as ``rainbow-shaped'', with best performance ``generally achieved when the motion is magnified in the range of $[8, 16]$'', and give a mechanistic account of both failure modes: too little magnification ``might not be enough to reveal the ME motion progress'', while too much ``degrades the performance because too many artifacts are induced''. Magnification is therefore a parameter with an interior optimum, not a monotone improvement to be maximised.

Bai et al.~\cite{bai2021} provide a second demonstration with a fundamentally different downstream model --- a pre-trained VGGFace2 encoder feeding a bidirectional LSTM --- reaching \textbf{75.76\,\% with phase-based magnification against a 60.60\,\% baseline} on SMIC, with the phase-based variant outperforming the amplitude-based one. The comparison is a distinct one and worth retaining, but the baseline is drawn from an earlier publication rather than from a magnification-off arm of their own, so it is not a matched control in the sense of \autoref{tab:evm-effect}. Because the pipeline shares nothing with hand-crafted descriptors, the two results together indicate the benefit is not an artefact of one feature family --- at least when the downstream representation is appearance-based.

Two further papers apply magnification in a different regime. Li, Huang and Zhao~\cite{liY2018,liY2021} magnify a single spotted apex frame at \textbf{$\alpha = 30$}, at the extreme end of Li et al.'s sweep, and augment by retaining the apex plus two frames either side. That $\alpha$ is asserted rather than derived, and the regimes differ enough --- a narrow window consumed as largely static appearance, against full magnified sequences described by spatio-temporal descriptors --- that the discrepancy need not be a contradiction. The safe reading is the weaker one: $\alpha$ is a parameter of the \textit{combination} of magnification and downstream representation, and cannot be transferred between pipelines without re-validation.

\subsection{Known limitations}
\label{sec:evm-limitations}

Four failure modes are documented, and the corpus is unusually candid about them.

\textbf{Magnification amplifies everything in band, not only expressions.} This is the most important and the most easily overlooked. Li et al.~\cite{liX2018} decline to use magnification for micro-expression \textit{spotting} at all, ``because it magnifies unwanted motions (e.g. head movements) at the same time''. A rigid head translation whose temporal frequency falls inside the pass-band is amplified exactly as a brow raise is. Prior alignment and cropping mitigate this; they do not eliminate it.

\textbf{Artefacts grow with $\alpha$}, reported by all four papers. \textbf{There is no principled procedure for choosing $\alpha$}: one value in the corpus is obtained by exhaustive sweep~\cite{liX2018}, and the apex-frame papers assert theirs. And \textbf{evaluation is confined to appearance-based downstream representations} --- LBP, HOG and HIGO in Li et al.~\cite{liX2018}, VGGFace2 features in Bai et al.~\cite{bai2021}, VGG-Face in both Li, Huang and Zhao papers. Every one is computed on magnified pixel intensities. \textbf{No study in the corpus measures what magnification contributes when the representation handed to the classifier is itself an explicit motion field.}

\subsection{Implications for this research}
\label{sec:evm-implications}
\label{sec:evm-motion-hypothesis}

The review fixes two settings and leaves one hypothesis. The temporal band is set to the \textbf{unnarrowed 5--25 Hz} duration rule, so Bai et al.'s narrowing is not followed and the pass-band admits the low-frequency range they rejected as noisy; this is a declared divergence. The factor is set to \textbf{$\alpha = 10$}, inside the reported $[8,16]$ optimum, and the apex papers' $\alpha = 30$ is explicitly not adopted because this thesis consumes full sequences. Amplitude-based magnification is used over a four-level Laplacian pyramid, and the band-pass is realised by an \textbf{FFT coefficient mask} rather than the Butterworth filter the reviewed description uses (\autoref{sec:temporal-filtering}) --- a further declared difference. Stage 1 runs twice, producing magnified and non-magnified tensor sets, so the switch selects a directory and leaves the network untouched.

Behind those settings lies a hypothesis the corpus does not test, stated because it frames the interpretation. Magnification amplifies the \textit{amplitude} of intensity variation in a band, and an appearance descriptor on magnified frames sees a directly stronger signal. An optical flow field is a different object: flow estimates \textit{displacement}, and amplifying the underlying intensity change by $\alpha$ does not simply scale the recovered displacement by $\alpha$. It changes the conditioning of the estimation problem, potentially making faint motion recoverable where it was not, while introducing displacement artefacts of the kind all four papers warn about; and any residual global scaling is further attenuated by the per-channel normalisation applied before training. The plausible outcome is that magnification in front of a flow-and-strain representation retains the part of its benefit consisting of \textit{making otherwise-unrecoverable motion recoverable} while losing the part consisting of \textit{increasing contrast for an appearance descriptor} --- so its measured contribution should be smaller on average than the +5 to +10 points of \autoref{tab:evm-effect}. Six matched pairs differing in nothing but the tensor directory are what allow that to be tested.

\textbf{A defect the review made detectable.} In this project's earliest evaluations every magnification-enabled configuration produced results \textit{bit-identical} to its disabled counterpart. Given that all four reviewed papers report clear, repeatable effects, an effect of exactly zero across six independent pairs is not a plausible null but a symptom: both arms were reading the same tensor directory, so the hypothesis had never been exercised. The defect was repaired before the reported run, and wiring the existing verification helper in as a start-up assertion is declared further work.

What the literature leaves open is threefold: every evaluation places magnification in front of an appearance representation, conflating two mechanisms that no experiment separates; $\alpha$ is transferred between studies without re-validation; and magnification is always evaluated as a two-cell addition to an otherwise fixed pipeline, so its interaction with other components is never measured. This thesis places it in front of an explicit motion representation, selects $\alpha$ from the full-sequence optimum, and measures it as one factor in a factorial ablation whose spread across architectures is itself the finding.


% =============================================================================
\section{Motion Representation: Dense Optical Flow}
\label{sec:flow-review}
% =============================================================================

Dense optical flow supplies the horizontal and vertical displacement channels. The relevant literature concerns three choices: the estimator, the frame pairs it processes, and whether flow serves as a recognition feature or an alignment tool. The estimation assumptions themselves are given in \autoref{sec:optical-flow}.

Xu et al.~\cite{xu2017} motivate facial-dynamics descriptors by the difficulty of interpreting micro-expressions through texture alone: flow describes apparent movement directly, whereas image intensity also records appearance and illumination. This does not guarantee independence from identity or recording conditions, since facial geometry, alignment and brightness changes still affect estimation. The benchmark evidence supports the choice --- the MEGC 2019 organisers report that their three highest-ranked submissions used optical flow~\cite{see2019}, OFF-ApexNet~\cite{liong2019a} uses horizontal and vertical flow as network inputs, STSTNet~\cite{liong2019b} adds strain, and Zhao et al.~\cite{zhaoS2021} retain a sequence of flow fields --- but this convergence is not a controlled demonstration that motion always outperforms appearance.

\textbf{The estimator is not settled by the corpus.} Zhao et al.~\cite{zhaoS2021} adopt Farneb\"ack through OpenCV, citing established availability; OFF-ApexNet~\cite{liong2019a} uses TV-L1, citing preservation of motion discontinuities; Xu et al.~\cite{xu2017} use the method of Sun, Roth and Black. STSTNet~\cite{liong2019b} specifies onset--apex flow without naming its estimator, so its similarity to OFF-ApexNet is insufficient to attribute TV-L1 to it. The advantages claimed for TV-L1 should not be read as a measured recognition gain over Farneb\"ack on CASME II: no controlled three-estimator comparison exists in the corpus. This thesis fixes Farneb\"ack as a preprocessing choice and leaves the comparison open.

\textbf{Onset--apex flow or a sequence is the consequential choice.} OFF-ApexNet~\cite{liong2019a} and STSTNet~\cite{liong2019b} represent each clip by flow from onset to apex, resized to $28\times28$ --- an economical summary of the displacement from neutral to peak, suitable for a small training set, which omits the intermediate trajectory and the relaxation after the apex. Zhao et al.~\cite{zhaoS2021} instead build an eleven-frame key-frame sequence and compute ten adjacent-pair fields, retaining changes in motion across the expression. The distinction is between one displacement summary and a sequence whose order a temporal model can exploit; it is not merely a difference in input size. This study follows the sequence approach but samples existing frames uniformly rather than constructing key frames, which preserves a temporal input for the transformer ablation. It does not measure the recognition cost of replacing a sequence with an onset--apex pair: that comparison is not an experimental arm.

Flow is also used for alignment rather than recognition. Xu et al.~\cite{xu2017} coarsely align by landmarks and then refine in-sequence using pixel-level movement, motivated specifically by micro-expressions, where a small registration error can be comparable to the displacement being recognised. This pipeline performs no flow-based fine alignment, so corpus preprocessing must not be treated as evidence that this error source has been eliminated. It remains the most promising unexploited improvement to preprocessing.

\textbf{Limitations of the representation.} Flow remains sensitive to brightness-constancy violations, alignment error and estimator settings; apex-based systems additionally depend on selecting an informative frame, and OFF-ApexNet~\cite{liong2019a} attributes weaker SMIC performance to imprecise apex selection. Using an onset--offset sequence avoids dependence on one selected apex but still relies on clip boundaries and sampling: repeated sampled frames yield near-zero-flow pairs, and large sampling intervals make the small-motion approximation less appropriate.

\subsection{Implications for this research}
\label{sec:flow-implications}

The central interpretive point is that optical flow is an \textit{analytical spatio-temporal feature extractor}: it estimates motion from paired frames before the network is trained. A learned spatial stem therefore acts on a representation that already makes displacement explicit, and with 156 clips, learning additional spatial filters may offer limited benefit or duplicate structure preprocessing has already exposed. This is a hypothesis about the stem's marginal value, not proof that learned spatial processing is unnecessary.

A temporal encoder addresses a different task: modelling relations among successive displacement fields, which pairwise flow estimation does not itself learn. Keeping those two roles distinct is necessary when attributing performance to the stem as against the transformer, and it is what the matched comparisons of this thesis are designed to separate. The design does not rank estimators, isolate alignment quality, or compare onset--apex against sequence representations; those remain limitations rather than gaps claimed as closed.


% =============================================================================
\section{Deformation Representation: The Optical Strain Tensor}
\label{sec:strain-review}
% =============================================================================

Optical strain supplies the third input channel. Its literature establishes a deformation descriptor derived from flow, several ways to use it, and limits relevant to interpreting the ablation. The tensor and its components are defined in \autoref{sec:optical-strain}; this section concerns the published evidence and one consequential divergence.

Flow describes displacement; strain describes its local spatial variation. A spatially uniform translation has zero displacement gradient, whereas neighbouring skin regions moving differently produce deformation. Shreve et al.~\cite{shreve2011} use this distinction to detect facial activity and describe robustness to moderate head translation, which does not imply immunity to arbitrary head movement or registration error. Their illumination- and make-up-robustness statements are attributed to earlier studies rather than measured anew, and should be read as reported motivation rather than as evidence about the present pipeline.

\subsection{The implemented magnitude and its divergence}
\label{sec:strain-divergence}

Liong et al.~\cite{liong2014a,liong2014b,liong2016} reduce the tensor to a scalar by the four-component Frobenius norm, which counts the two equal off-diagonal shear entries separately. \textbf{The implementation here counts shear once rather than twice}, giving the three-term magnitude of \autoref{sec:optical-strain}. Spatial derivatives come from central differences at one-pixel spacing, following Liong et al.~\cite{liong2014a}, but the scalar reduction differs from their convention.

STSTNet~\cite{liong2019b} cannot establish a precedent for the three-term form. Its printed magnitude equation repeats $\partial u$ in the mixed derivative where the symmetric tensor calls for $\partial v/\partial x$ alongside $\partial u/\partial y$, and it places the factor one-half \emph{outside} the squared derivative sum. Correcting the repeated derivative while retaining that placement yields the \textit{four}-term convention; moving the factor inside the square would instead yield the implemented form, but that is a further change to the printed expression. The intended implementation cannot be settled from the published equation alone, so this thesis declares a departure from the unambiguous four-term convention rather than claiming to reproduce STSTNet's magnitude.

The four-term form weights pure shear by $\sqrt{2}$ relative to the implemented one. Independent min--max normalisation removes a uniform scale factor but does not generally remove this difference, because normal and shear contributions vary across pixels and time. Neither convention was compared experimentally here, so their practical equivalence is not assumed.

\subsection{Recognition evidence and its protocol limits}
\label{sec:strain-evidence}

Strain has been used three ways: as a pooled descriptor (optical strain features, OSF), as a spatial weighting on LBP-TOP histograms (optical strain weighted features, OSW), and as a direct input channel. Liong et al.~\cite{liong2016} combine the first two, pooling spatially then temporally, and report the results of \autoref{tab:strain-performance}; the early work also applies Wiener and Gaussian filters explicitly to suppress noise~\cite{liong2014a}. STSTNet~\cite{liong2019b} supplies strain directly alongside onset--apex flow, which is the precedent for the three-channel input used here, independently of the magnitude ambiguity --- though this study retains all three channels across successive frame pairs rather than reducing each clip to one map.

\begin{table}[htbp]
	\centering
	\small
	\begin{tabular}{@{}lcc@{}}
		\hline
		\textbf{Method} & \textbf{SMIC accuracy (\%)} & \textbf{CASME II accuracy (\%)} \\
		\hline
		Baseline (texture) & 45.73 & 61.94 \\
		OSF                & 41.46 & 51.01 \\
		OSF + OSW          & 52.44 & \textbf{63.16} \\
		OFF + OFW (flow)   & 40.24 & \textbf{55.87} \\
		\hline
	\end{tabular}
	\caption[Strain-based descriptors against a matched flow control]{Strain-based descriptors against a matched flow-based control, after Liong et al.~\cite{liong2016} (Table 8). SMIC uses leave-one-subject-out; CASME II uses leave-one-\textit{video}-out, which is not subject-disjoint. The summary combines block settings, so equal entries elsewhere in the source are not matched-setting comparisons.}
	\label{tab:strain-performance}
\end{table}

The CASME II evidence depends on how strain is used and against what. OSF alone falls below the texture baseline; the combined descriptor improves on it by 0.41 to 1.22 points across block settings~\cite{liong2016}. Crucially, Liong et al. also test flow features and weights within the same framework, reaching 55.87\,\% against 63.16\,\% for the strain pair --- a \textbf{matched flow-versus-strain control}, so it would be incorrect to claim the literature never isolates strain against flow. What it concerns is hand-crafted descriptors under a protocol that is not subject-disjoint, and these accuracy gains are not comparable with pooled macro F1. It does not establish the gain from adding a strain channel to a learned sequence model.

\subsection{Implications for this research}
\label{sec:strain-implications}

Strain extends the analytical-feature argument of \autoref{sec:flow-implications}: the input contains a \textit{hand-specified differential operator}, exposing local deformation without learning that operator from 156 clips. A spatial stem may still learn useful combinations of those features, but it operates after both motion estimation and deformation extraction. This motivates testing its incremental contribution; it does not establish the cause of any observed architectural effect.

Three limits bound what this thesis can say about strain. Spatial differentiation can amplify flow-estimation noise, and \textbf{no dedicated strain post-filter is applied}, departing from the explicit filtering of Liong et al.~\cite{liong2014a}; Farneb\"ack's internal smoothing is not a separately evaluated strain filter. The shear weighting is the untested choice of \autoref{sec:strain-divergence}. Above all, \textbf{strain is present in all twelve evaluated configurations}, so its marginal contribution is not measured at all: a strain-on/strain-off comparison is the natural next ablation, and neither published gains from a different descriptor nor the performance of the complete system substitutes for it. A deterministic transformation of $u$ and $v$ adds no independent observation; its potential value lies in making deformation directly accessible to the classifier, and whether that helps must be measured rather than inferred from its construction.


% =============================================================================
\section{Temporal Normalisation: Fixed-Length Sampling}
\label{sec:tim-review}
% =============================================================================

Flow and strain fix \textit{what} is computed from each frame pair; temporal normalisation fixes \textit{which} pairs, and how many. The Temporal Interpolation Model (TIM) is credited in this corpus to Zhou et al. by Li et al.~\cite{liX2013,liX2018} and to Pfister et al. by Lu et al.~\cite{lu2015}, and is described here only through the reviewed papers that implement it.

Normalisation is necessary because any batched network requires a fixed input shape and the corpus supplies no constant, spanning a four-fold range of sequence length; and because duration is itself a nuisance variable --- Lu et al.~\cite{lu2015} normalise to remove ``temporal fluctuations of micro-expressions, which may introduce noise irrelevant to classification''. They also give the corpus's most compact description of TIM: ``a manifold-based interpolation method, which builds a low-dimensional embedding of an image sequence, and then interpolates a curve in the low-dimensional space. The interpolated frames are mapped back to the original high-dimensional space.'' The property that matters downstream is that TIM \textbf{synthesises frames that were never recorded}, having been proposed ``to generate sufficient frames from an image sequence'' --- an upsampler rather than only a length-standardiser.

\textbf{Ten frames is a convention, not a finding, and the studies that test it disagree.} TIM appears throughout the corpus as an unremarked preprocessing step, almost always at ten frames~\cite{liX2013,liX2018,lu2015}. Li et al.~\cite{liX2018} sweep interpolation length from 10 to 80 on the SMIC variants with the descriptor held fixed, and report that ``the best performance for all three datasets is achieved with interpolation to 10 frames'', because ``in longer sequences, changes along the time-dimension are diluted''. Ben et al.~\cite{ben2021} reach the opposite conclusion: holding alignment, descriptor and classifier fixed, they interpolate to 10 through 110 frames by both Newton interpolation and TIM on MMEW, finding Newton peaking at 33.3\,\% at 60 frames and TIM at \textbf{38.9\,\%} at 30 or 60, and adopt TIM at 60. Accuracy is non-monotonic in length in both. Xu et al.~\cite{xu2017} add the only multi-corpus sweep --- 10, 15 and 20 frames across CASME, CASME II and both SMIC subsets --- and depart from convention by using \textbf{linear} interpolation ``instead of Temporal Interpolation Model because in a high frame-rate scenario, linear interpolation suffices to characterize the motion pattern with less error''.

Two caveats limit what transfers. Ben et al.'s experiment uses MMEW with LBP-TOP, not CASME II and not a learned model, and runs largely in the \textit{upsampling} regime, so an optimum measured while synthesising frames need not hold when discarding them. And \textbf{frame synthesis is in tension with motion analysis}: where the downstream representation is appearance-based, synthesised frames are benign, but where it is motion, synthesised frames yield synthesised displacements. Xu et al.~\cite{xu2017} are the only reviewed paper to raise this, and they assert it rather than measure it. No reviewed study compares the two interpolators against a motion representation.

\subsection{Implications for this research}
\label{sec:tim-implications}

\textbf{A declaration about what this pipeline implements.} The module performing this step is named for the Temporal Interpolation Model, but what it computes is \textbf{uniform index sampling} of existing frames (\autoref{sec:three-channel-tensor}). No embedding is built, no curve is interpolated, no frame is synthesised. That is a materially different operation from TIM, and it is named accurately here rather than in the pipeline's own terms. It is defensible: because the downstream representation is flow between adjacent sampled frames, interpolation would manufacture displacement fields for motion that never occurred, whereas subsampling discards temporal resolution without fabricating it. What is given up is TIM's ability to \textit{lengthen} short clips.

The length is fixed at \textbf{33 frames sampled, giving 32 flow pairs}, bounded by the annotated onset and offset so the window spans the expression rather than surrounding neutral frames. That is close to the lower of Ben et al.'s two reported optima, but the correspondence should not be overstated: their 30 frames are \textit{interpolated} frames under a hand-crafted descriptor on a different corpus, which is not direct validation of 32 sampled flow fields under a learned model here. The length is not swept, and a sweep is declared as further work.

Three quantitative consequences follow, and they constrain the magnification conclusion rather than merely describing the sampler.

\textbf{Effective temporal resolution falls by a clip-dependent factor.} Sampling 33 frames from a clip of $N$ frames recorded at 200~fps gives an effective rate of $6600/N$ fps: \textbf{100 fps} at the median $N = 66$, which is SMIC's recording rate and half what CASME II was built to provide, and roughly \textbf{52 fps} for the longest clip, below the original CASME database's 60~fps~\cite{yan2014}. Eighty-one of the 156 clips are downsampled by a factor of two or more, which gives back part of what the 200~fps capture rate bought. No reviewed paper reports this cost.

\textbf{Magnification is applied after subsampling, which shifts the realised pass-band.} Magnification runs on the 33 sampled frames at the nominal 200~fps, but those frames span the clip's full duration, so the true inter-frame interval is $N/(200 \times 32)$ seconds rather than $1/200$. At the median, a band specified as 5--25 Hz corresponds to roughly 2.5--12.5 Hz physically, and the discrepancy grows with clip length. The realised band is therefore neither the 5--25 Hz that \autoref{sec:evm-implications} derives nor constant across clips. This is a declared divergence; magnifying before subsampling is declared as further work.

\textbf{Three clips are shorter than the sampling window.} Three of the 156 contain 31 frames against a target of 33, so sampling repeats indices and yields near-zero flow at those positions.

The unresolved question this leaves is what temporal normalisation costs a representation made of motion rather than appearance. Three studies sweep interpolation length and give three answers, only one of them touching CASME II and none using a learned temporal model. This thesis takes the position that follows from its own representation and quantifies what that costs in effective frame rate, but does not sweep the length, compare interpolators, or separate the effect of sequence length from that of the sampling regime.


% =============================================================================
\section{The Learned Spatial Backbone: Shallow Convolutional Networks}
\label{sec:cnn-review}
% =============================================================================

By 2019 the field had converged on optical flow as the input; the open question was what should consume it. A learned backbone promises what hand-crafted descriptors cannot~\cite{liong2019b}, but runs into the small-data constraint of \autoref{sec:corpus-threats}.

\textbf{The term ``3D CNN'' covers three materially different architectures in this literature.} STSTNet~\cite{liong2019b}, the canonical shallow model, calls itself a ``Shallow Triple Stream Three-dimensional CNN'', but its input is a single $28 \times 28 \times 3$ tensor whose third dimension holds horizontal flow, vertical flow and strain --- \textbf{not time}. Its three parallel single-layer streams, with 3, 5 and 8 kernels, supplement ``the small scale input data \ldots to avoid the problem of underfitting''; the clip was collapsed to one frame pair upstream. STRCN~\cite{xia2020a} by contrast extends convolutional connectivity genuinely ``across temporal domain'' with recurrent convolutional layers. Zhao et al.~\cite{zhaoS2021} apply a C3D-style network to a ten-frame flow sequence, so the temporal axis is present but short. Results reported across the corpus under the same label are therefore not like with like.

\textbf{The central controlled finding is that shrinking both model and input helps.} Xia et al.~\cite{xia2020b} treat depth and input resolution as controlled variables, building four models from a single convolutional layer to one convolutional plus three recurrent convolutional layers, and evaluating each across input resolutions from 20 to 300 pixels on optical flow. Deeper is worse: models 3 and 4 ``have worse recognition performances compared to shallower \ldots models'', even though all four are far shallower than ResNet-101. Higher resolution is worse with a threshold: performance ``will be degraded with larger input resolutions'' and for the deeper models ``begins to decrease dramatically when the resolution is bigger than \textbf{100 $\times$ 100}'' --- though shallow models are ``basically robust to the change of input resolution''. That qualification matters for this thesis, and so does the setting: Xia et al.'s study and STSTNet are both developed for the MEGC composite regime, which supplies substantially more training data per fold than a single corpus (\autoref{sec:megc-metrics}), so conclusions about affordable capacity understate the constraint here. Xia et al.~\cite{xia2020b} also take a third route to capacity control that this thesis follows --- developing ``three parameter-free modules \ldots integrate with RCN without increasing any learnable parameters'', so representational power grows while overfitting risk does not.

\textbf{No study isolates the backbone itself.} The corpus is not short of matched comparisons: it sets descriptor-only methods against learned ones on the same folds, and Xia et al.~\cite{xia2020b} vary individual architectural modules against a common baseline with every other parameter held constant. But those comparisons all \textit{add} to a backbone that stays in place. None removes the learned spatial stem entirely and passes the motion representation straight to a classifier or temporal model, so what the stem contributes against what it costs is nowhere measured.

\subsection{Implications for this research}
\label{sec:cnn-implications}

The component ablated here is a \textbf{three-stream shallow convolutional backbone}, one unshared two-layer stream per input channel, carrying \textbf{14,544 parameters} --- shallow by the corpus's standards, though roughly nine times STSTNet's. Where STSTNet's streams differ in filter count over the same input, these differ in input channel. Removing it leaves not an empty slot but spatial pooling followed by a learned linear projection, so the no-backbone route still contains learned input processing.

\textbf{A declaration about what this component is.} Every kernel in the backbone is spatial only: no filter spans two time steps, so despite the \texttt{Conv3d} implementation it is a \textbf{purely spatial} feature extractor (\autoref{sec:conv3d}). The ablation therefore does not test spatio-temporal convolution; it tests a learned spatial stem, and whatever it measures is evidence about \textit{this} stem --- spatial-only, at 224~$\times$~224, on flow and strain input --- and \textbf{not} evidence that convolution over time is unhelpful for micro-expression recognition, which this study never evaluated. In the field's own terms the choice is defensible: STSTNet, one of the strongest shallow baselines on the CASME II subset, likewise performs no temporal convolution.

\textbf{Input resolution is the clearest departure from Xia et al.'s guidance.} The backbone is fed at \textbf{224 $\times$ 224}, more than twice their threshold. This is not a departure from the corpus as a whole --- Bai et al.~\cite{bai2021} resize to 224 $\times$ 224 without alignment or cropping, as here; Li, Huang and Zhao~\cite{liY2021} resize faces to the same; and Xia et al.'s own earlier work~\cite{xia2020a} runs one stream at 300 $\times$ 245 --- but it departs from the specific threshold above, which is a composite-database result that bites hardest on models deeper than this one. Since this stem is shallow, the reviewed evidence predicts a smaller resolution effect than the deep-model figures suggest. The choice is also the expensive one: a 224 $\times$ 224 map carries sixty-four times the spatial elements of STSTNet's input, and the backbone convolves it at every one of the 32 time steps. Any effect measured for the stem is therefore \textbf{a resolution result as much as an architecture result}, and the reviewed evidence does not determine which mechanism dominates; re-running the backbone arm at a lower resolution is the first item of declared further work because it is cheap and would settle the question, not because the corpus implies a particular outcome.


% =============================================================================
\section{Parameter-Free Attention: SimAM}
\label{sec:simam-review}
% =============================================================================

SimAM re-weights convolutional feature maps without adding learned parameters. Yang et al.~\cite{yang2021} supply the mechanism and its image-recognition evidence; Xia et al.~\cite{xia2020b} motivate parameter-free modules in composite-database micro-expression recognition, supporting the design principle rather than validating SimAM on this task. The energy formulation and gate are given in \autoref{sec:simam}.

Yang et al.~\cite{yang2021} motivate attention through spatial suppression: neurons distinguishable from their surroundings are assigned greater importance, expressed as the linear separability of one activation from the others in its channel. This is a prior about activation statistics, not a detector of action units or minority classes --- a distinctive response may encode useful deformation, noise or an artefact --- so its suitability for flow-derived features requires evaluation rather than following from the neuroscience motivation. Two implementation details are owned here rather than by \autoref{ch:background}: the shared variance printed with the paper's equation~(5) divides by $M$ whereas its Figure~3 reference code divides by $M-1$, and this thesis follows the code; and the module's ``three-dimensional weights'' describe $C \times H \times W$ with $M = HW$ per frame, not a temporal volume, so extending the context across frames is an additional design decision rather than part of the original formulation.

Yang et al. evaluate SimAM with several network families on CIFAR and ImageNet, reporting competitive performance against learned attention modules, and their timing results show that unchanged parameter totals do not imply unchanged inference speed. For CIFAR they sweep $\lambda$ from $10^{-1}$ to $10^{-6}$ and select $10^{-4}$; they repeat the selection for ImageNet and choose $0.1$. Thus $10^{-4}$ is a dataset-selected operating point, not a universal setting justified by the formula.

\textbf{The limits of transfer are specific.} The source experiments concern natural images, different training-set sizes and learned feature distributions. Here the input is flow and strain, so distinctiveness need not bear the same relationship to useful evidence. The image experiments also do not decide whether video statistics should be computed per frame or over a clip. And the absence of learned attention weights does not establish that runtime, memory or overfitting concerns disappear. These are empirical limits on transferring the result, not reasons to assume success or failure.

\subsection{Implications for this research}
\label{sec:simam-implications}

SimAM is applied separately to each convolutional stream before concatenation, introducing no trainable weights or biases anywhere. Two choices diverge from the source. Mean and variance span $(T,H,W)$ within each channel with divisor $THW - 1$, rather than one $(H,W)$ plane, so an activation is judged distinctive relative to its whole clip rather than its own frame; this could emphasise a brief contraction, but does not guarantee low weights at onset or high weights at the apex, since the score depends on squared deviation and not on expression labels. And $\lambda = 10^{-4}$ is retained without re-tuning on CASME II. Neither context choice is compared experimentally.

SimAM is also enabled only alongside the convolutional stem in this model, which describes this implementation rather than a mathematical prohibition. Four of the sixteen candidate configurations are therefore excluded, leaving twelve evaluated and \textbf{four SimAM matched pairs}, and the measured effect is conditional on the stem being present. Retention should be judged against both recognition performance and computational cost: parameter count alone cannot settle it, because the module is cheap rather than free.

\subsection{Relevance and scope}
\label{sec:simam-gap}

The study evaluates SimAM in a motion-input, small-corpus setting through four matched comparisons under complete leave-one-subject-out, testing the implemented module as a whole rather than establishing the benefit of each design choice. Three follow-ups remain: whether the distinctiveness prior suits flow-derived features, re-tuning $\lambda$, and comparing per-frame with whole-clip statistics. Each requires new evaluations despite the module introducing no learned weights.


% =============================================================================
\section{The Temporal Encoder: Transformers}
\label{sec:transformer-review}
% =============================================================================

The temporal encoder is motivated by Zhang et al.'s Short and Long Range Relation Based Spatio-Temporal Transformer (SLSTT)~\cite{zhang2022} and its Vision Transformer foundation~\cite{dosovitskiy2021}. The component evaluated here differs substantially from published SLSTT, particularly in the axis over which self-attention operates, and distinguishing the two is the purpose of this section.

A micro-expression evolves through onset, peak and relaxation, so retaining a motion sequence permits a model to relate separated stages rather than compressing them immediately into one descriptor; Zhang et al.~\cite{zhang2022} motivate their design through the need to capture both local and global patterns. This thesis compares a temporal transformer against mean pooling over the same sequence: mean pooling discards order, whereas positional information and self-attention allow relations among frame-pair features to be modelled. That motivates the comparison without establishing in advance which generalises better. ViT itself~\cite{dosovitskiy2021} splits an image into patches, projects them into tokens, adds learned positional embeddings and relates them by attention; Dosovitskiy et al. also show that without the locality and translation-equivariance biases of convolution, image transformers can perform poorly when trained with insufficient data, while large-scale pre-training improves them. That is a caution for a 156-clip study, not a universal requirement that every transformer be pre-trained.

\textbf{Published SLSTT combines three stages.} First, \textbf{onset-referenced optical flow} describes sampled frames relative to the onset rather than to their immediate predecessors, which Zhang et al. argue gives a more structured trajectory towards and away from the apex; their experiments use it throughout, without a matched short-term-flow ablation establishing its independent gain. Second, a \textbf{spatial ViT encoder} processes each flow field as patches, using ImageNet-pre-trained ViT-B/16 with twelve layers at $384 \times 384$. Third, an \textbf{LSTM temporal aggregator} combines the per-frame features, with a mean aggregator supplying the comparison variant. Self-attention is spatial in this architecture; recurrence handles time.

\subsection{The aggregation comparison, and what it can support}
\label{sec:slstt-results}

The comparison relevant here is internal to SLSTT and holds the spatial transformer fixed in both arms: CASME II UF1 rises from \textbf{0.844} with the mean aggregator to \textbf{0.901} with the LSTM, a difference of 0.057~\cite{zhang2022}. That supports evaluating aggregation choices within that framework. It is \textbf{not} an estimate of what an LSTM would add to this thesis, where mean pooling follows \emph{temporal} self-attention rather than independent spatial encoding, and it isolates an aggregator rather than the presence of a transformer, since both variants contain one.

Protocol labels need care. Zhang et al. describe three-class leave-one-subject-out evaluations on the databases separately and together under a Composite Database Evaluation heading, so that heading alone does not establish that every CASME II entry uses pooled cross-database training; they also report a distinct Sole Database Evaluation using original labels, including five CASME II classes~\cite{zhang2022}. Neither setting should be equated automatically with this thesis's 156-clip three-class pool. Published SLSTT also retains a pre-trained spatial encoder, so its results do not isolate the benefit of a transformer trained from scratch on this study's motion representation.

\subsection{Implementation and six declared divergences}
\label{sec:transformer-implications}

The evaluated component receives 32 tokens of width 96, from either the convolutional stem or spatial pooling followed by a learned linear projection. The encoder uses four pre-norm layers, eight heads, feed-forward width 256 and dropout 0.1; fixed sinusoidal positions precede attention and mean pooling follows it. With the encoder disabled the input tokens are averaged directly.

\begin{table}[htbp]
	\centering
	\small
	\begin{tabular}{@{}p{0.23\linewidth} p{0.34\linewidth} p{0.35\linewidth}@{}}
		\hline
		\textbf{Property} & \textbf{Published SLSTT~\cite{zhang2022}} & \textbf{This thesis} \\
		\hline
		Attention axis   & Spatial patches within a frame              & Temporal tokens within a clip \\
		Aggregation      & LSTM over spatially encoded frames          & Mean after temporal self-attention \\
		Initialisation   & ImageNet-pre-trained ViT-B/16               & Trained from scratch \\
		Capacity         & Twelve layers, width 768; $384\times384$ in & Four layers, width 96; 32 tokens \\
		Positions        & Learned spatial embeddings                  & Fixed sinusoidal temporal encoding \\
		Input flow       & Relative to onset                           & Relative to the preceding sampled frame \\
		\hline
	\end{tabular}
	\caption[Divergences from published SLSTT]{Six material divergences from published SLSTT. The component evaluated here is a small temporal transformer encoder, not a reproduction of the published architecture.}
	\label{tab:slstt-divergences}
\end{table}

Six matched pairs evaluate the encoder with other component choices fixed within each pair. They test the whole encoder --- positional encoding, feed-forward layers and normalisation included --- against direct mean pooling, and do not isolate self-attention from those accompanying operations or compare it against another learned temporal model. Three follow-ups follow from \autoref{tab:slstt-divergences}: testing an LSTM aggregator, comparing onset-referenced with consecutive-frame flow, and replacing fixed positions with learned embeddings. All three remain untested, and changing the input representation requires re-training before any benefit can be claimed. What this study can establish is a matched comparison of the implemented encoder against an order-blind baseline within this pipeline; it neither validates published SLSTT nor establishes a general ranking of attention, recurrence and convolution.


% =============================================================================
\section{Training Under Severe Class Imbalance}
\label{sec:imbalance-review}
% =============================================================================

The 99 : 32 : 25 distribution of \autoref{sec:label-taxonomy} is a property of what can be elicited in a laboratory, not a sampling accident, so it cannot be corrected by collecting more data from the same paradigm. The corpus corrects for it at three points --- the data, the loss and the metric --- and this thesis uses all three. The mechanisms are given in \autoref{sec:focal-loss} and \autoref{sec:balanced-sampling}.

\textbf{At the loss.} Zhao et al.~\cite{zhaoS2021} bring focal loss into micro-expression recognition ``to alleviate the inefficient model training caused by the class imbalance in the MEs Datasets'', adopting the formulation developed for one-stage object detection. The division of labour inside it is what matters: $\gamma$ re-weights by \textit{difficulty} and is treated as a constant, ``set as 2 in practice'', while $\alpha$ re-weights by \textit{class frequency} and is ``treated as a hyper-parameter to set by cross validation''. Conflating the two is the error discussed below. Zhao et al. compare focal loss against cross-entropy across a sweep of $\alpha$, but hold $\gamma$ fixed, so the focusing exponent remains unvalidated on this task.

\textbf{At the data.} Xia et al.~\cite{xia2020a} pair ``temporal data augmentation strategies as well as a balanced loss \ldots jointly'' to overcome ``limited and imbalanced training samples''. Their own ablation removes each correction in turn and reports the balanced loss improving performance only ``slightly'' --- at most about two points, and in one cell costing accuracy. This is the corpus's only controlled look at composing two corrections, and no reviewed paper reports the pairing failing outright.

\textbf{At the metric.} \autoref{sec:megc-metrics} defines the pooled macro-averaged measure. Liong et al.~\cite{liong2016} add the argument from the subject side, reporting F1, precision and recall because these ``provide a more meaningful perspective than accuracy rates when the datasets used are naturally imbalanced since each subject has a different number of video samples.'' Both arguments compound here, where classes are skewed \textbf{and} per-subject clip counts range from 1 to 33.

One further gap: \textbf{label smoothing is named in the corpus but never examined on this task.} Dosovitskiy et al.~\cite{dosovitskiy2021} list it among the regularisers they tune, but no reviewed micro-expression paper uses it or reports what it does under class skew.

\subsection{Implications for this research}
\label{sec:imbalance-double-correction}

Focal loss is adopted with \textbf{$\gamma = 2.0$}, matching Zhao et al.'s reported setting. The class-frequency term is implemented as an optional per-class $\alpha$ vector but \textbf{switched off at run time}, because frequency correction is assigned to a weighted sampler active in every configuration. Label smoothing of 0.05 is applied inside the focal loss, declared as a choice this corpus supports neither way. The primary metric is pooled macro F1, with accuracy quoted only alongside the always-Negative floor.

\textbf{What this project contributes is a negative result about composing corrections.} The configuration is resolved at run time by a single rule: class weighting in the loss is enabled only when the balanced sampler is \textit{not} active. That rule exists because of two failures in opposite directions, both recoverable from this project's own run history. In an earlier code era the loss carried inverse-frequency class weights and no sampler existed; under that single correction the proposed configuration abandoned the \textbf{Positive} class outright, a failure initially misdiagnosed as over-capacity overfitting. A later era added the sampler while leaving the loss weights unconditional, so both corrections ran at once; the collapse then inverted, and configurations abandoned the \textbf{Negative} majority class instead, one predicting a single class for every clip. With the guard in place, no configuration in the reported study abandons any class.

This is worth reporting because no reviewed paper warns that combining corrections can invert the intended effect rather than merely blunt it. It must be stated carefully, however: the two collapses occurred in code eras differing from the reported run in more than the correction rule, and the final sweep contains no matched arm with the guard disabled. \textbf{The guard's causal benefit is therefore not isolated}, and the evidence does not establish that correcting at exactly one point is universally best --- only that correcting twice, in this pipeline, produced a class collapse that correcting once did not.

A second declaration follows. Because every configuration shares the same sampler, loss and $\gamma$, this thesis measures \textbf{no} effect attributable to the imbalance treatment: it is a constant, not a factor. The claim supported is that the four architectural components were compared under a training regime that does not abandon minority classes, not that this regime is optimal.

\subsection{Relevance and scope}
\label{sec:imbalance-gap}

The corrections themselves are well established. What the corpus does not settle is how two corrections compose, what $\gamma$ should be on this task, and what over-correction looks like when it fails. This thesis contributes to the first and third by documenting two opposite class-collapse failures and the run-time rule that resolved them. On the second it offers no improvement, since $\gamma = 2.0$ and label smoothing of 0.05 are held constant across all twelve configurations; sweeping either is declared as further work.


% =============================================================================
\section{Synthesis and Research Gap}
\label{sec:synthesis-and-gap}
% =============================================================================

This section consolidates the preceding relevance statements into one research gap, and collects in one place every declared divergence between the system built here and the literature it draws on.

Read as a whole, the reviewed literature settles a good deal. The corpora are small, demographically narrow, laboratory-bound and irreducibly imbalanced by what can be elicited, and CASME II remains the primary single-corpus benchmark under its native leave-one-subject-out protocol. Motion rather than appearance is the appropriate input; a strain field adds a further cue, though on CASME II only weakly and under a protocol that is not subject-disjoint; and both the magnification factor and the sequence length have interior optima located only by sweep. At this data size, capability bought without parameters outperforms capability bought with them, and self-attention is the strongest-performing mechanism reported. Accuracy alone is inadequate under skew and must give way to macro-averaged measures accumulated across folds before averaging.

\subsection{Four limitations that recur across every topic}
\label{sec:four-limitations}

The per-section statements above are not independent observations. Four structural problems generate most of them.

\textbf{(a) Components are isolated one at a time, inside one proposed design, and never jointly.} Matched-pair measurement is not absent from the corpus: \autoref{tab:evm-effect} isolates magnification, \autoref{tab:strain-performance} isolates strain against a flow control, Xia et al.~\cite{xia2020b} compare a recurrent convolutional network against three variants each adding one module with every other parameter held constant, and SLSTT's Mean-versus-LSTM comparison isolates its aggregator (\autoref{sec:slstt-results}). What no reviewed study does is vary \textit{several} components together, so that their interactions can be read. Every ablation is internal to one proposed design and moves one factor from that design's own defaults; across papers, systems are compared whole, so a difference in the reported number cannot be attributed to any one stage. There is therefore no basis for deciding which learned stages are load-bearing, nor whether one component's value depends on another's presence.

\textbf{(b) The evaluation protocol is unstable and under-reported.} Sample counts, label-set size and the choice among leave-one-subject-out, leave-one-video-out and leave-one-sample-out vary inconsistently, and the field's most-quoted baseline is described three ways by its own sources (\autoref{sec:original-baseline}). Given this corpus's fold composition these are not bookkeeping differences: they change what the numbers mean.

\textbf{(c) The strongest results depend on regimes a single-corpus study cannot access.} The best CASME II figures come from composite-database training (\autoref{sec:megc-metrics}) and, for SLSTT, ImageNet pre-training --- neither a property of the architecture being credited. Almost no evidence exists about what these components do when trained from scratch on one small corpus, which is the situation most practitioners are in.

\textbf{(d) Preprocessing parameters have interior optima and are transferred without re-validation.} The magnification factor, the sequence length and the input resolution all rise then fall, and none has a principled selection rule. Values are carried between pipelines whose downstream models differ entirely: $\alpha = 30$ moves from a single-frame appearance encoder to other settings, ten-frame interpolation propagates from one baseline paper to the whole field, and $\lambda = 10^{-4}$ transfers from CIFAR though its own authors re-tuned it for ImageNet.

\subsection{The research gap}
\label{sec:research-gap}

Taking (a) to (d) together:

\textbf{The reviewed literature has established a standard pipeline for micro-expression recognition --- magnify, compute motion, normalise duration, extract spatially, model temporally --- and contains matched ablations of individual stages, but no study varies several of those stages jointly so that their relative contributions and interactions can be read. Because components are varied one at a time inside a single proposed system, under protocols that differ in unstated ways, and predominantly in training regimes richer than a single small corpus provides, it is not currently possible to say from this literature which parts of that pipeline earn their place, which are inherited convention, and which depend on the presence of another.}

This matters practically as well as scientifically. A practitioner assembling this pipeline today has no published basis for allocating effort: for judging whether the most computationally expensive stage is the most valuable, whether a component with a negligible average effect may still matter in a specific failure, or whether a ranking read off an incomplete evaluation of a corpus this size is stable at all.

\subsection{How this thesis addresses the gap, and what it does not}
\label{sec:addressing-gap}

\textbf{What it does.} This thesis treats the pipeline as a factorial experiment rather than a proposal: four components varied independently across every architecturally valid cell of a 2\textsuperscript{4} matrix, twelve configurations, each evaluated under \textbf{complete 25-fold leave-one-subject-out on CASME II} with every non-varied factor held identical. That supplies the joint measurement (a) identifies as missing, reporting the protocol in full addresses (b) for this study at least, and training from scratch is the regime (c) identifies as unevidenced. Each component's marginal contribution is read from pairs of runs differing in one factor alone, and the distribution of those four effects is the substantive answer to the gap. Two further contributions follow from the review rather than the experiment: a \textbf{repaired measurement}, where an early data-routing defect left the magnification switch inert and the review's evidence of a clear, repeatable magnification effect is what made bit-identical results detectable as a defect rather than a null (\autoref{sec:evm-implications}); and \textbf{two opposite class-collapse failures} that no reviewed paper reports (\autoref{sec:imbalance-double-correction}).

\textbf{What it does not do.} Three of the four limitations are only partially addressed. On (b), this study reports its own protocol completely but cannot repair the field's inconsistency, and its own N = 156 differs from MEGC's 145. On (c), it establishes what these components do from scratch on one corpus, which is a different question from --- not an answer to --- how they behave under composite training. On (d), the thesis inherits parameter values rather than sweeping them: $\alpha = 10$, T = 32, $\lambda = 10^{-4}$ and $\gamma = 2.0$ are fixed throughout, so the ablation measures architecture at one point in preprocessing space rather than across it.

\subsection{A consolidated ledger of declared divergences}
\label{sec:divergence-ledger}

Reviewing the literature closely enough to write this chapter surfaced twelve points at which the implemented system departs from the work it draws on. Several were discovered only through the review. They are collected in \autoref{tab:divergence-ledger} so that no result rests on an unstated assumption.

\begin{table}[htbp]
	\centering
	\footnotesize
	\begin{tabular}{@{}p{0.85cm} p{6.0cm} p{6.6cm}@{}}
		\hline
		\textbf{\S} & \textbf{Divergence} & \textbf{Consequence} \\
		\hline
		\ref{sec:label-taxonomy} & Working set is N = 156, including sadness and fear in Negative; MEGC's CASME II subset is N = 145 & Comparisons against challenge figures are not exactly like-for-like \\
		\ref{sec:casme2-construction} & The corpus's own registered release is not used; the pipeline reads the original recorded frames, unregistered and uncropped & No landmarking or alignment of any kind precedes the network; tolerance to pose and scale must be learned from 156 clips \\
		\ref{sec:evm-implications} & Temporal band is the unnarrowed 5--25 Hz duration rule, where Bai et al. narrowed it to 15--25 Hz as noisy & The pass-band admits the low-frequency range they rejected \\
		\ref{sec:evm-implications} & Band-pass is realised by an FFT coefficient mask, not the Butterworth filter of the reviewed description & A different filter shape, with different transition behaviour at the band edges \\
		\ref{sec:strain-divergence} & Strain magnitude counts shear once, not twice as in the published Frobenius norm & Shear weighted lower by $\sqrt{2}$ relative to normal strain; untested \\
		\ref{sec:strain-implications} & No dedicated strain filtering, where the literature applies Wiener and Gaussian filters & Spatial differentiation can amplify flow-estimation noise, and nothing downstream suppresses it \\
		\ref{sec:tim-implications} & Temporal normalisation is uniform index sampling, \textbf{not} manifold-based TIM; no frames are synthesised & Avoids fabricating motion; forfeits the ability to lengthen short clips \\
		\ref{sec:tim-implications} & Magnification is applied \textbf{after} subsampling, at a nominal 200 fps & Realised pass-band is clip-dependent and below the intended 5--25 Hz \\
		\ref{sec:cnn-implications} & Backbone kernels are spatial only --- \textbf{no temporal mixing} & The ablation tests a learned \textit{spatial} stem, not spatio-temporal convolution \\
		\ref{sec:cnn-implications} & Backbone input is 224 $\times$ 224, against Xia et al.'s guidance of $\leq$ 100 $\times$ 100 --- though Bai et al. and Li, Huang and Zhao both use 224 $\times$ 224 & Whatever the stem is measured to contribute may be a resolution effect as much as an architecture effect \\
		\ref{sec:simam-implications} & SimAM statistics pooled over the whole clip rather than per frame as in the source; $\lambda$ not re-tuned & Neurons judged against the clip, not the frame; defensible but untested \\
		\ref{sec:transformer-implications} & Six divergences from published SLSTT (temporal not spatial attention, mean not LSTM aggregation, from scratch, smaller, sinusoidal encoding, short-term flow) & The temporal encoder is \textbf{not} a reproduction of SLSTT \\
		\ref{sec:imbalance-double-correction} & Label smoothing of 0.05, named in the corpus only by Dosovitskiy et al. and never examined on this task & A choice the micro-expression corpus supports neither way, held constant and unmeasured \\
		\hline
	\end{tabular}
	\caption[Every declared departure from the reviewed literature]{Every declared departure from the reviewed literature, with its consequence. The first column resolves to the subsection that owns each divergence.}
	\label{tab:divergence-ledger}
\end{table}

Three of these --- the strain formula, the magnification ordering and the input resolution --- were identified by the review itself rather than being known design decisions, and each is declared as proposed further work. None invalidates the ablation, since every configuration shares the same choices. What they constrain is the \textit{wording} of the conclusions: what is measured is the contribution of a spatial stem at 224 $\times$ 224 and of temporal self-attention in this pipeline, not a verdict on convolution over time or on SLSTT. The limits of what can be claimed are therefore fixed before any number is reported --- a small, single-seed, demographically narrow corpus measured against richer published regimes, with these divergences standing throughout.
 IN PLACE OF chapter3.tex and
% chapter3_shortened.tex. Never include two of them: they share labels.
%
% WHAT CHANGED
%   * Repeated mechanism/derivation passages removed: they are owned by
%     Chapter 2 (magnification, brightness constancy, strain tensor, SimAM
%     energy, focal loss, F1). This chapter now carries published evidence and
%     transfer limits only, and contains NO display mathematics.
%   * The nine per-section "research gap" endings are replaced by short
%     relevance statements, with one qualified synthesis at the end.
%   * Local decision tables (EVM/flow/strain/TIM/CNN/SimAM/imbalance) removed:
%     they restated Chapter 2 or 4 and were then restated again in the ledger.
%   * 17 tables -> 6. 2 figures -> 1. The dataset-composition figure is gone
%     (its content is now one prose sentence); the fold-composition figure is
%     retained because the pooled-versus-fold argument depends on it.
%
% THE LEDGER'S FIRST COLUMN IS NOW \ref{}, NOT HARD-CODED NUMBERS.
% The twelve divergence rows previously cited 3.1.2 ... 3.9.7 literally, which
% any restructuring silently falsifies. They now resolve automatically.
%
% LABELS: every label referenced from introduction.tex, chapter4_shortened.tex,
% chapter5_shortened.tex or chapter6.tex is preserved (24 of them), plus those
% used by the unshortened chapters, so either manuscript state resolves.
% Verified by scratchpad scan; re-run it after any further edit.
%
% NOTE: the old header warned that Chapter 2 cites Chapter 3 subsections BY
% NUMBER. That is no longer true -- the shortened Chapter 2 has none -- so
% subsection positions here are free. The only by-number dependency was the
% ledger's own first column, now removed.
% =============================================================================

\chapter{Literature Review}
\label{ch:literature}

This chapter reviews the literature bearing on each component of the recognition pipeline evaluated in this thesis: the corpus every experiment runs on, the four techniques placed under ablation, the input transformations they operate on, and the training regime held constant across all of them. Each section states what the relevant studies found, what their evidence does not establish, and what follows for this experiment. Mechanisms already derived in \autoref{ch:background} are not derived again. \autoref{sec:synthesis-and-gap} consolidates the closing statements into one research gap and collects every declared departure from the reviewed methods.


% =============================================================================
\section{The Evaluation Corpus: CASME II and the Spontaneous Benchmark Family}
\label{sec:corpus}
% =============================================================================

The corpus is treated first because on a database of this size several decisions usually presented as modelling choices are in fact forced by properties of the data.

\subsection{Why a dedicated spontaneous corpus was necessary}
\label{sec:why-spontaneous}

Automatic expression recognition matured on large corpora of \textit{posed} expressions, reporting accuracies above 90\,\% on the six basic emotions by the early 2010s~\cite{yan2014}. Micro-expression research could not follow, because the phenomenon is involuntary and cannot be performed on instruction. Two problems motivated CASME II directly.

The first is \textbf{construct validity}: Yan et al.~\cite{yan2014} and Li et al.~\cite{liX2013} both observe that deliberately produced ``micro-expressions'' differ systematically from spontaneous ones, so a system trained on posed data may be a recogniser of fast deliberate movement rather than of micro-expressions. The second is \textbf{acquisition quality}. Defined jointly by very short duration and very low intensity, the phenomenon sits close to the noise floor of ordinary video; Yan et al.~\cite{yan2014} argue that CASME's 60~fps coarsely samples a movement lasting a fraction of a second --- while marking the benefit of higher rates as unresolved --- and that mains-frequency flicker contaminated earlier recordings. Elicitation was a third difficulty: deception paradigms produced samples contaminated by conversational movement, and the workable paradigm, which CASME II adopts, has participants watch high-arousal video while attempting to keep the face neutral~\cite{liX2013}.

\begin{table}[htbp]
	\centering
	\footnotesize
	\setlength{\tabcolsep}{4pt}
	\begin{tabular}{l c c c c c c}
		\hline
		\textbf{Corpus} & \textbf{Subj.} & \textbf{Micro samples} & \textbf{fps} & \textbf{Spont.} & \textbf{FACS} & \textbf{Ethn. / age} \\
		\hline
		SMIC~\cite{liX2013}              & 16 valid          & 164                & 100          & Yes          & No           & 3 / 26.7 \\
		\textbf{CASME II}~\cite{yan2014} & \textbf{26 valid} & \textbf{247 / 255} & \textbf{200} & \textbf{Yes} & \textbf{Yes} & \textbf{1 / 22.03} \\
		SAMM~\cite{davison2018}          & 32                & 159                & 200          & Yes          & Yes          & 13 / 33.24 \\
		CAS(ME)$^2$~\cite{qu2016}        & 22                & 53 (+250 macro)    & 30           & Yes          & Yes          & 1 / 22.59 \\
		\hline
	\end{tabular}
	\caption[The spontaneous micro-expression corpora]{The spontaneous corpora that define the current context, compiled from Yan et al.~\cite{yan2014}, Davison et al.~\cite{davison2018}, Li et al.~\cite{liX2013} and Qu et al.~\cite{qu2016}. CASME II carries two sample counts: 247 in the paper, 255 in the released coding.}
	\label{tab:spontaneous-corpora}
\end{table}

\subsection{CASME II: construction and annotation}
\label{sec:casme2-construction}

CASME II~\cite{yan2014} recruited thirty-five participants, mean age 22.03 (SD 1.60), recording at \textbf{200 fps} and 640~$\times$~480 without inter-frame compression under flicker-free illumination, giving a facial region of roughly 280~$\times$~340 pixels. Each retained sample carries onset and offset indices, FACS action-unit labels from two independent coders with disagreements arbitrated by discussion, and an emotion label assigned from the coded action units, participant self-report and the eliciting episode's content --- Yan et al. arguing against forcing spontaneous samples into the six basic categories, since a single action unit is frequently ambiguous.

Two specifics propagate into this thesis. The original release applies a normalisation from which this work departs: 68 Active Shape Model landmarks are detected on the \textit{first frame only} and the resulting transform applied unchanged to every subsequent frame; this pipeline reads the original recorded frames instead, unregistered and uncropped. And the published paper reports \textbf{247} micro-expressions while the distributed coding file contains \textbf{255}, a post-publication revision rarely acknowledged and one reason sample counts for ``CASME II'' vary. The present work uses the released 255-sample coding.

\subsection{The label taxonomy and the sparsity it creates}
\label{sec:label-taxonomy}

Yan et al.~\cite{yan2014} provide five emotion categories with action-unit criteria and are candid about the result: ``some types of facial expressions are difficult to elicit in laboratory situations, thus the samples in different categories distributed unequally.'' The released coding is finer, exposing seven raw labels. Counting directly from the label table used here gives happiness 32, disgust 63, surprise 25, repression 27, sadness 7, fear 2, and others 99.

\textbf{Seven-class classification is not viable on this corpus.} Fear is represented by two clips in the entire database and sadness by seven. Under any subject-disjoint protocol --- and subject-disjointness has been this benchmark's norm since its introduction (\autoref{sec:original-baseline}) --- a two-clip class cannot be meaningfully trained or tested, and per-class metrics degenerate.

The field's response is to group by affective valence. The convention originates with SMIC, released with exactly three classes~\cite{liX2013}, and was carried into cross-corpus work by MEGC 2019~\cite{see2019}. Discarding ``others'' and grouping the remainder gives the working set used throughout: \textbf{Negative 99, Positive 32, Surprise 25, N = 156}, a ratio of roughly 4 : 1.3 : 1.

\textbf{A declared divergence from the MEGC convention.} MEGC 2019 reports the CASME II portion of its composite database as \textbf{145} samples, footnoting that CASME II's Negative class covers ``Disgust and Repression''~\cite{see2019} and building its subset from the paper's five-category release, in which sadness and fear are absorbed into ``others''. This thesis uses the seven-label coding and assigns sadness and fear to Negative, following the mapping MEGC gives for \textbf{SAMM} rather than for CASME II. The consequence is an \textbf{eleven-clip difference}: N = 156 against 145, with Negative at 99 rather than 88. Nine of the eleven are the sadness and fear clips; the remaining two are unaccounted for, since disgust and repression together number 90 in the released coding against the 88 the challenge reports --- a residual noted rather than resolved. The direction is mild: the nine added clips all fall in the majority class and come from the categories coders found hardest to assign, so if anything they make this task marginally harder than the challenge's.

\subsection{The original baseline, and what it does not establish}
\label{sec:original-baseline}

Yan et al.~\cite{yan2014} baseline the corpus with LBP-TOP~\cite{zhaoG2007} and a support vector machine, sweeping spatial and temporal radii and excluding $R_T = 1$ deliberately, because at 200~fps the change between adjacent frames is negligible --- an early acknowledgement that a high sampling rate is not automatically an advantage. Five-class accuracies under \textbf{leave-one-subject-out} cluster in the high fifties to low sixties, the best configuration reaching \textbf{63.41\,\%}. Three observations matter.

\textbf{Leave-one-subject-out is native to this benchmark}, not a stricter protocol adopted later, so any evaluation reporting a single random split departs from the benchmark's own convention. \textbf{The baseline is reported as accuracy alone}, and on a five-class problem where ``others'' is 102 of 247 samples the majority-constant classifier scores 41\,\% without modelling anything, so a headline near 63.41\,\% is not by itself evidence of a working system. \textbf{And the sources disagree on the protocol behind that figure}: Yan et al.~\cite{yan2014} state leave-one-subject-out in abstract and results prose, the caption of the table reporting it reads ``leave-one-out cross-validation'' without the subject qualifier, and Li et al.~\cite{liX2018} reproduce the figure marked as leave-one-\textit{sample}-out, a materially easier protocol. The most widely quoted number on this benchmark thus rests on a protocol its own sources describe three ways; this thesis follows the originating paper's prose and records the discrepancy rather than resolving it.

Accuracy also moves by several points across a purely descriptive radius sweep, with no change to data, split or classifier. Where effect sizes of that magnitude arise from descriptor settings, similar differences between architectures cannot be interpreted without a variance estimate --- a limitation this thesis inherits and declares.

\subsection{Sibling corpora, and why this study remains single-corpus}
\label{sec:sibling-corpora}

Three further databases define the context, each isolating a specific limitation of CASME II. \textbf{SMIC}~\cite{liX2013} is the origin of the three-class taxonomy and of fixed-length temporal interpolation, both inherited here; it is not used as a corpus because it is neither FACS-coded nor apex-annotated. \textbf{SAMM}~\cite{davison2018} addresses the CASME family's demographic narrowness with 13 ethnicities against CASME II's one, Davison et al. observing that a corpus drawn from one ethnicity and one narrow age band restricts analysis to ``similar looking participants''; it is stronger on that count but smaller in usable samples and spread across seven categories. \textbf{CAS(ME)$^2$}~\cite{qu2016} pairs macro- and micro-expression samples from the same participants, making it the natural resource for domain transfer --- the strategy behind the strongest MEGC 2019 entry~\cite{liu2019,see2019} --- but records at 30~fps, supplies only 53 micro-expression samples, and is evaluated leave-one-\textit{video}-out, which does not enforce subject-disjointness.

CASME II is used alone because it has the highest temporal resolution of any spontaneous corpus of comparable size, a precondition for magnification in a defined band and for dense frame-to-frame flow; because its onset, apex and offset annotation makes temporal normalisation well-defined rather than heuristic; and because it remains the most widely reported single-corpus benchmark. The cost --- that no claim here extends to other populations, capture conditions or corpora --- is declared rather than argued away.

\subsection{How CASME II is evaluated: MEGC 2019 and the metric convention}
\label{sec:megc-metrics}

The most consequential development in how CASME II is \textit{used} is the Composite Database Evaluation protocol of the Second Facial Micro-Expression Grand Challenge~\cite{see2019}. It merges SMIC, CASME II and SAMM into one composite database of \textbf{442 samples} --- Negative 250, Positive 109, Surprise 83 --- of which CASME II contributes 145, and runs leave-one-subject-out across the composite. Plain accuracy is rejected explicitly on grounds of class imbalance, and two balanced metrics mandated: \textbf{Unweighted F1 (UF1)}, the macro-averaged F1 defined in \autoref{sec:confusion-matrix}, and \textbf{Unweighted Average Recall}, the balanced accuracy.

What matters is not the algebra but the \textbf{construction}: UF1 accumulates true positives, false positives and false negatives per class \textbf{over all folds} \textit{before} any per-class F1 is computed. As \autoref{sec:pooled-vs-fold} establishes, that is a different estimator from the average of per-fold macro F1 scores, and the distinction is not cosmetic here. Because clips are distributed very unequally across subjects and many single-subject folds contain only one of the three classes, a per-fold macro F1 is bounded well below one by fold composition alone (\autoref{fig:fold-composition}). The primary metric of this thesis, \textit{pooled macro F1}, is therefore identical in construction to UF1; the per-fold average appearing in the raw logs is not a metric the field uses.

\begin{figure}[htbp]
	\centering
	\includegraphics[width=\textwidth]{fig3_2_fold_composition.png}
	\caption[Class composition of the 25 leave-one-subject-out folds]{Class composition of the 25 leave-one-subject-out folds, and the per-fold macro-F1 ceiling it implies.}
	\label{fig:fold-composition}
\end{figure}

\textbf{One structural caveat governs the challenge's reported figures}, which range up to 0.7885 composite UF1 and 0.8764 on the CASME II subset~\cite{see2019}. Every one comes from a model trained on the \textit{composite} database --- roughly 442 samples across three corpora and 68 subjects --- then scored on the CASME II portion of the pooled predictions. A CASME II column is therefore not a CASME II-only result: it reports how a cross-corpus model performs on CASME II clips, with approximately \textbf{2.8 times} the training data available in any single-corpus fold, where the experiment conducted here trains each fold on roughly 150 clips. The two regimes are not interchangeable, and single-corpus results should not be compared against those figures without stating the difference.

Two further observations from the challenge bear on this design. The top three submissions all chose \textbf{optical flow} as their representation, a convergence that supports the motion-based input used here; and one submission substituted mid-position frames for annotated apex frames with little loss, suggesting that precise apex localisation may be less critical than assumed.

\subsection{Threats to validity carried by the corpus}
\label{sec:corpus-threats}

The following are not criticisms of CASME II, which remains the best available resource for this task, but constraints any study built on it inherits. Fold and duration statistics are computed from the 156-clip label table used here.

\textbf{Scale and imbalance.} 156 usable clips across three classes, distributed 99 : 32 : 25, small enough that parameter count is itself a risk factor. This motivates parameter-free or low-parameter components wherever a choice exists, and it makes accuracy unusable alone: the always-Negative classifier attains \textbf{0.6346 accuracy} at a macro F1 of only \textbf{0.2588}, so that floor is quoted alongside every result.

\textbf{Subject-level non-uniformity.} \textbf{Subject 17 alone contributes 33 clips, 21\,\% of the working set}, while three subjects contribute one each, so fold sizes vary by a factor of 33; \textbf{ten of the 25 folds contain a single class}, eight contain two, and only seven all three. This forces full 25-fold leave-one-subject-out and is why the pooled metric rather than the fold average is primary. It is also why a single held-out split is not a reliable instrument here: which subjects fall in the test set can move the result by more than any architectural change under study. Subject 18's three clips are all ``others'' and are removed by the three-class filter, so the run has \textbf{25 folds, not 26} --- stated explicitly because fold counts differ between papers using this corpus.

\textbf{Demographic homogeneity and laboratory ecology.} All participants are Chinese with mean age 22.03, named by Davison et al.~\cite{davison2018} as a principal motivation for SAMM; the MEGC organisers saw the effect in reverse, reporting that adaptation from CK+ helped on SMIC and SAMM but ``not so for CASME II, which contain predominantly Chinese subjects''~\cite{see2019}. Recording is frontal under fixed lamps with participants instructed not to move or speak, so almost none of the head motion, illumination variation, occlusion or speech-related movement a deployed system would face is present. Together these motivate a \textbf{motion} representation rather than raw pixels --- appearance being exactly what a 26-subject corpus would let a network memorise --- and they are why no cross-dataset or in-the-wild claim is made anywhere in this thesis.

\textbf{Duration, annotation and version ambiguity.} Onset-to-offset lengths range from \textbf{31 to 126 frames}, median 66 --- at 200~fps, 0.155\,s to 0.63\,s --- so every clip is resampled to a fixed length before training. Labels derive from coded action units, self-report and stimulus content: a construct produced by trained coders under a rubric, taken here as given and recorded as a declared limitation rather than as measurement error. And three sample counts circulate under the name ``CASME II'' --- 247, 255 and 145 --- with both five- and seven-class label sets in use and rarely stated, so N is reported explicitly alongside every figure here.

\subsection{Implications for this research}
\label{sec:corpus-implications}

Two further properties determine settings rather than threatening validity: the \textbf{200 fps} capture rate fixes the magnification pass-band at a matched \textbf{5--25 Hz}, and the baseline's sensitivity to a descriptive hyper-parameter sets the standard of evidence, so effect sizes below the resolution of a 156-clip corpus are reported as unresolved rather than as findings. No confidence interval is available, since the single-seed design and unretained per-clip predictions preclude one.

The consequence that shapes interpretation rather than construction is that \textbf{the small-sample regime predicts which components should help}. 156 clips cannot supply the data needed to learn a general spatio-temporal feature extractor from scratch, particularly once much of that extraction has been performed analytically before the network sees the data, whereas a component modelling the temporal arc of an expression exploits exactly the structure the corpus's onset, apex and offset annotation guarantees is present. This is a hypothesis the experiment tests, not a result it assumes. What the corpus literature does not supply is a stable basis for comparing methods: only Li et al.~\cite{liX2018} quantify how much of a reported difference is attributable to the protocol rather than the method, and then only for a single matched row of their own method.


% =============================================================================
\section{Motion Magnification: Eulerian Video Magnification}
\label{sec:evm-review}
% =============================================================================

Magnification is reviewed through the four corpus papers that apply it: Bai et al.~\cite{bai2021}, Li et al.~\cite{liX2018}, and Li, Huang and Zhao~\cite{liY2018,liY2021}. The original amplitude- and phase-based formulations lie outside the corpus and are described only as those papers attribute them. The mechanism itself is given in \autoref{sec:motion-magnification}.

The corpus design of \autoref{sec:why-spontaneous} answers the \textit{short-duration} half of the micro-expression definition but not the \textit{low-intensity} half. A movement spread across dozens of frames at 200~fps may still occupy only a few grey levels across a small patch of skin, and Yan et al.~\cite{yan2014} report discarding samples outright because the movement was too faint to place an onset. Li et al.~\cite{liX2018} state the remedy directly --- amplify the motion before any descriptor is computed, so that the difference between categories is enlarged rather than the classifier being made more sensitive. Being entirely upstream of the model is what makes magnification a clean ablation factor here.

Three parameters are free in any application: the temporal pass-band, which decides \textit{which} motions are amplified; the factor $\alpha$, which decides \textit{how much}; and the pyramid depth. On the first, Bai et al.~\cite{bai2021} report the rule given by the original authors --- the band should match the duration of the phenomenon, for micro-expressions 1/5 to 1/25\,s, equivalently \textbf{5--25 Hz} --- while noting that this full band produced ``noticeable noise'' in their setting, which led them to narrow it to 15--25 Hz. This is the only explicit guidance on band selection anywhere in the corpus.

\subsection{Evidence that magnification improves recognition}
\label{sec:evm-evidence}

The strongest case is Li et al.~\cite{liX2018}, who sweep the magnification factor rather than fixing it, across $\alpha \in \{1, 2, 4, 8, 12, 16, 20, 24, 30\}$ with $\alpha = 1$ denoting no magnification, on both SMIC and CASME II under leave-one-subject-out. Two findings emerge.

\textbf{Magnification helps consistently}, ``for all three kinds of features on all four testing datasets'', and on CASME II substantially (\autoref{tab:evm-effect}); their best full configuration reaches 67.21\,\% on CASME II, competitive with the contemporaneous state of the art.

\begin{table}[htbp]
	\centering
	\small
	\begin{tabular}{l c c c}
		\hline
		\textbf{Descriptor} & \textbf{No magnification} & \textbf{With magnification} & \textbf{$\Delta$} \\
		\hline
		LBP  & 55.87\,\% & 60.73\,\% & \textbf{+4.86} \\
		HOG  & 57.49\,\% & 63.97\,\% & \textbf{+6.48} \\
		HIGO & 57.09\,\% & 67.21\,\% & \textbf{+10.12} \\
		\hline
	\end{tabular}
	\caption[Effect of Eulerian motion magnification on CASME II]{Effect of magnification on CASME II under leave-one-subject-out, after Li et al.~\cite{liX2018} (Table 6). This is the corpus's one controlled magnification-on/off comparison.}
	\label{tab:evm-effect}
\end{table}

\textbf{The relationship between $\alpha$ and accuracy is non-monotonic.} Li et al.~\cite{liX2018} describe the curves as ``rainbow-shaped'', with best performance ``generally achieved when the motion is magnified in the range of $[8, 16]$'', and give a mechanistic account of both failure modes: too little magnification ``might not be enough to reveal the ME motion progress'', while too much ``degrades the performance because too many artifacts are induced''. Magnification is therefore a parameter with an interior optimum, not a monotone improvement to be maximised.

Bai et al.~\cite{bai2021} provide a second demonstration with a fundamentally different downstream model --- a pre-trained VGGFace2 encoder feeding a bidirectional LSTM --- reaching \textbf{75.76\,\% with phase-based magnification against a 60.60\,\% baseline} on SMIC, with the phase-based variant outperforming the amplitude-based one. The comparison is a distinct one and worth retaining, but the baseline is drawn from an earlier publication rather than from a magnification-off arm of their own, so it is not a matched control in the sense of \autoref{tab:evm-effect}. Because the pipeline shares nothing with hand-crafted descriptors, the two results together indicate the benefit is not an artefact of one feature family --- at least when the downstream representation is appearance-based.

Two further papers apply magnification in a different regime. Li, Huang and Zhao~\cite{liY2018,liY2021} magnify a single spotted apex frame at \textbf{$\alpha = 30$}, at the extreme end of Li et al.'s sweep, and augment by retaining the apex plus two frames either side. That $\alpha$ is asserted rather than derived, and the regimes differ enough --- a narrow window consumed as largely static appearance, against full magnified sequences described by spatio-temporal descriptors --- that the discrepancy need not be a contradiction. The safe reading is the weaker one: $\alpha$ is a parameter of the \textit{combination} of magnification and downstream representation, and cannot be transferred between pipelines without re-validation.

\subsection{Known limitations}
\label{sec:evm-limitations}

Four failure modes are documented, and the corpus is unusually candid about them.

\textbf{Magnification amplifies everything in band, not only expressions.} This is the most important and the most easily overlooked. Li et al.~\cite{liX2018} decline to use magnification for micro-expression \textit{spotting} at all, ``because it magnifies unwanted motions (e.g. head movements) at the same time''. A rigid head translation whose temporal frequency falls inside the pass-band is amplified exactly as a brow raise is. Prior alignment and cropping mitigate this; they do not eliminate it.

\textbf{Artefacts grow with $\alpha$}, reported by all four papers. \textbf{There is no principled procedure for choosing $\alpha$}: one value in the corpus is obtained by exhaustive sweep~\cite{liX2018}, and the apex-frame papers assert theirs. And \textbf{evaluation is confined to appearance-based downstream representations} --- LBP, HOG and HIGO in Li et al.~\cite{liX2018}, VGGFace2 features in Bai et al.~\cite{bai2021}, VGG-Face in both Li, Huang and Zhao papers. Every one is computed on magnified pixel intensities. \textbf{No study in the corpus measures what magnification contributes when the representation handed to the classifier is itself an explicit motion field.}

\subsection{Implications for this research}
\label{sec:evm-implications}
\label{sec:evm-motion-hypothesis}

The review fixes two settings and leaves one hypothesis. The temporal band is set to the \textbf{unnarrowed 5--25 Hz} duration rule, so Bai et al.'s narrowing is not followed and the pass-band admits the low-frequency range they rejected as noisy; this is a declared divergence. The factor is set to \textbf{$\alpha = 10$}, inside the reported $[8,16]$ optimum, and the apex papers' $\alpha = 30$ is explicitly not adopted because this thesis consumes full sequences. Amplitude-based magnification is used over a four-level Laplacian pyramid, and the band-pass is realised by an \textbf{FFT coefficient mask} rather than the Butterworth filter the reviewed description uses (\autoref{sec:temporal-filtering}) --- a further declared difference. Stage 1 runs twice, producing magnified and non-magnified tensor sets, so the switch selects a directory and leaves the network untouched.

Behind those settings lies a hypothesis the corpus does not test, stated because it frames the interpretation. Magnification amplifies the \textit{amplitude} of intensity variation in a band, and an appearance descriptor on magnified frames sees a directly stronger signal. An optical flow field is a different object: flow estimates \textit{displacement}, and amplifying the underlying intensity change by $\alpha$ does not simply scale the recovered displacement by $\alpha$. It changes the conditioning of the estimation problem, potentially making faint motion recoverable where it was not, while introducing displacement artefacts of the kind all four papers warn about; and any residual global scaling is further attenuated by the per-channel normalisation applied before training. The plausible outcome is that magnification in front of a flow-and-strain representation retains the part of its benefit consisting of \textit{making otherwise-unrecoverable motion recoverable} while losing the part consisting of \textit{increasing contrast for an appearance descriptor} --- so its measured contribution should be smaller on average than the +5 to +10 points of \autoref{tab:evm-effect}. Six matched pairs differing in nothing but the tensor directory are what allow that to be tested.

\textbf{A defect the review made detectable.} In this project's earliest evaluations every magnification-enabled configuration produced results \textit{bit-identical} to its disabled counterpart. Given that all four reviewed papers report clear, repeatable effects, an effect of exactly zero across six independent pairs is not a plausible null but a symptom: both arms were reading the same tensor directory, so the hypothesis had never been exercised. The defect was repaired before the reported run, and wiring the existing verification helper in as a start-up assertion is declared further work.

What the literature leaves open is threefold: every evaluation places magnification in front of an appearance representation, conflating two mechanisms that no experiment separates; $\alpha$ is transferred between studies without re-validation; and magnification is always evaluated as a two-cell addition to an otherwise fixed pipeline, so its interaction with other components is never measured. This thesis places it in front of an explicit motion representation, selects $\alpha$ from the full-sequence optimum, and measures it as one factor in a factorial ablation whose spread across architectures is itself the finding.


% =============================================================================
\section{Motion Representation: Dense Optical Flow}
\label{sec:flow-review}
% =============================================================================

Dense optical flow supplies the horizontal and vertical displacement channels. The relevant literature concerns three choices: the estimator, the frame pairs it processes, and whether flow serves as a recognition feature or an alignment tool. The estimation assumptions themselves are given in \autoref{sec:optical-flow}.

Xu et al.~\cite{xu2017} motivate facial-dynamics descriptors by the difficulty of interpreting micro-expressions through texture alone: flow describes apparent movement directly, whereas image intensity also records appearance and illumination. This does not guarantee independence from identity or recording conditions, since facial geometry, alignment and brightness changes still affect estimation. The benchmark evidence supports the choice --- the MEGC 2019 organisers report that their three highest-ranked submissions used optical flow~\cite{see2019}, OFF-ApexNet~\cite{liong2019a} uses horizontal and vertical flow as network inputs, STSTNet~\cite{liong2019b} adds strain, and Zhao et al.~\cite{zhaoS2021} retain a sequence of flow fields --- but this convergence is not a controlled demonstration that motion always outperforms appearance.

\textbf{The estimator is not settled by the corpus.} Zhao et al.~\cite{zhaoS2021} adopt Farneb\"ack through OpenCV, citing established availability; OFF-ApexNet~\cite{liong2019a} uses TV-L1, citing preservation of motion discontinuities; Xu et al.~\cite{xu2017} use the method of Sun, Roth and Black. STSTNet~\cite{liong2019b} specifies onset--apex flow without naming its estimator, so its similarity to OFF-ApexNet is insufficient to attribute TV-L1 to it. The advantages claimed for TV-L1 should not be read as a measured recognition gain over Farneb\"ack on CASME II: no controlled three-estimator comparison exists in the corpus. This thesis fixes Farneb\"ack as a preprocessing choice and leaves the comparison open.

\textbf{Onset--apex flow or a sequence is the consequential choice.} OFF-ApexNet~\cite{liong2019a} and STSTNet~\cite{liong2019b} represent each clip by flow from onset to apex, resized to $28\times28$ --- an economical summary of the displacement from neutral to peak, suitable for a small training set, which omits the intermediate trajectory and the relaxation after the apex. Zhao et al.~\cite{zhaoS2021} instead build an eleven-frame key-frame sequence and compute ten adjacent-pair fields, retaining changes in motion across the expression. The distinction is between one displacement summary and a sequence whose order a temporal model can exploit; it is not merely a difference in input size. This study follows the sequence approach but samples existing frames uniformly rather than constructing key frames, which preserves a temporal input for the transformer ablation. It does not measure the recognition cost of replacing a sequence with an onset--apex pair: that comparison is not an experimental arm.

Flow is also used for alignment rather than recognition. Xu et al.~\cite{xu2017} coarsely align by landmarks and then refine in-sequence using pixel-level movement, motivated specifically by micro-expressions, where a small registration error can be comparable to the displacement being recognised. This pipeline performs no flow-based fine alignment, so corpus preprocessing must not be treated as evidence that this error source has been eliminated. It remains the most promising unexploited improvement to preprocessing.

\textbf{Limitations of the representation.} Flow remains sensitive to brightness-constancy violations, alignment error and estimator settings; apex-based systems additionally depend on selecting an informative frame, and OFF-ApexNet~\cite{liong2019a} attributes weaker SMIC performance to imprecise apex selection. Using an onset--offset sequence avoids dependence on one selected apex but still relies on clip boundaries and sampling: repeated sampled frames yield near-zero-flow pairs, and large sampling intervals make the small-motion approximation less appropriate.

\subsection{Implications for this research}
\label{sec:flow-implications}

The central interpretive point is that optical flow is an \textit{analytical spatio-temporal feature extractor}: it estimates motion from paired frames before the network is trained. A learned spatial stem therefore acts on a representation that already makes displacement explicit, and with 156 clips, learning additional spatial filters may offer limited benefit or duplicate structure preprocessing has already exposed. This is a hypothesis about the stem's marginal value, not proof that learned spatial processing is unnecessary.

A temporal encoder addresses a different task: modelling relations among successive displacement fields, which pairwise flow estimation does not itself learn. Keeping those two roles distinct is necessary when attributing performance to the stem as against the transformer, and it is what the matched comparisons of this thesis are designed to separate. The design does not rank estimators, isolate alignment quality, or compare onset--apex against sequence representations; those remain limitations rather than gaps claimed as closed.


% =============================================================================
\section{Deformation Representation: The Optical Strain Tensor}
\label{sec:strain-review}
% =============================================================================

Optical strain supplies the third input channel. Its literature establishes a deformation descriptor derived from flow, several ways to use it, and limits relevant to interpreting the ablation. The tensor and its components are defined in \autoref{sec:optical-strain}; this section concerns the published evidence and one consequential divergence.

Flow describes displacement; strain describes its local spatial variation. A spatially uniform translation has zero displacement gradient, whereas neighbouring skin regions moving differently produce deformation. Shreve et al.~\cite{shreve2011} use this distinction to detect facial activity and describe robustness to moderate head translation, which does not imply immunity to arbitrary head movement or registration error. Their illumination- and make-up-robustness statements are attributed to earlier studies rather than measured anew, and should be read as reported motivation rather than as evidence about the present pipeline.

\subsection{The implemented magnitude and its divergence}
\label{sec:strain-divergence}

Liong et al.~\cite{liong2014a,liong2014b,liong2016} reduce the tensor to a scalar by the four-component Frobenius norm, which counts the two equal off-diagonal shear entries separately. \textbf{The implementation here counts shear once rather than twice}, giving the three-term magnitude of \autoref{sec:optical-strain}. Spatial derivatives come from central differences at one-pixel spacing, following Liong et al.~\cite{liong2014a}, but the scalar reduction differs from their convention.

STSTNet~\cite{liong2019b} cannot establish a precedent for the three-term form. Its printed magnitude equation repeats $\partial u$ in the mixed derivative where the symmetric tensor calls for $\partial v/\partial x$ alongside $\partial u/\partial y$, and it places the factor one-half \emph{outside} the squared derivative sum. Correcting the repeated derivative while retaining that placement yields the \textit{four}-term convention; moving the factor inside the square would instead yield the implemented form, but that is a further change to the printed expression. The intended implementation cannot be settled from the published equation alone, so this thesis declares a departure from the unambiguous four-term convention rather than claiming to reproduce STSTNet's magnitude.

The four-term form weights pure shear by $\sqrt{2}$ relative to the implemented one. Independent min--max normalisation removes a uniform scale factor but does not generally remove this difference, because normal and shear contributions vary across pixels and time. Neither convention was compared experimentally here, so their practical equivalence is not assumed.

\subsection{Recognition evidence and its protocol limits}
\label{sec:strain-evidence}

Strain has been used three ways: as a pooled descriptor (optical strain features, OSF), as a spatial weighting on LBP-TOP histograms (optical strain weighted features, OSW), and as a direct input channel. Liong et al.~\cite{liong2016} combine the first two, pooling spatially then temporally, and report the results of \autoref{tab:strain-performance}; the early work also applies Wiener and Gaussian filters explicitly to suppress noise~\cite{liong2014a}. STSTNet~\cite{liong2019b} supplies strain directly alongside onset--apex flow, which is the precedent for the three-channel input used here, independently of the magnitude ambiguity --- though this study retains all three channels across successive frame pairs rather than reducing each clip to one map.

\begin{table}[htbp]
	\centering
	\small
	\begin{tabular}{@{}lcc@{}}
		\hline
		\textbf{Method} & \textbf{SMIC accuracy (\%)} & \textbf{CASME II accuracy (\%)} \\
		\hline
		Baseline (texture) & 45.73 & 61.94 \\
		OSF                & 41.46 & 51.01 \\
		OSF + OSW          & 52.44 & \textbf{63.16} \\
		OFF + OFW (flow)   & 40.24 & \textbf{55.87} \\
		\hline
	\end{tabular}
	\caption[Strain-based descriptors against a matched flow control]{Strain-based descriptors against a matched flow-based control, after Liong et al.~\cite{liong2016} (Table 8). SMIC uses leave-one-subject-out; CASME II uses leave-one-\textit{video}-out, which is not subject-disjoint. The summary combines block settings, so equal entries elsewhere in the source are not matched-setting comparisons.}
	\label{tab:strain-performance}
\end{table}

The CASME II evidence depends on how strain is used and against what. OSF alone falls below the texture baseline; the combined descriptor improves on it by 0.41 to 1.22 points across block settings~\cite{liong2016}. Crucially, Liong et al. also test flow features and weights within the same framework, reaching 55.87\,\% against 63.16\,\% for the strain pair --- a \textbf{matched flow-versus-strain control}, so it would be incorrect to claim the literature never isolates strain against flow. What it concerns is hand-crafted descriptors under a protocol that is not subject-disjoint, and these accuracy gains are not comparable with pooled macro F1. It does not establish the gain from adding a strain channel to a learned sequence model.

\subsection{Implications for this research}
\label{sec:strain-implications}

Strain extends the analytical-feature argument of \autoref{sec:flow-implications}: the input contains a \textit{hand-specified differential operator}, exposing local deformation without learning that operator from 156 clips. A spatial stem may still learn useful combinations of those features, but it operates after both motion estimation and deformation extraction. This motivates testing its incremental contribution; it does not establish the cause of any observed architectural effect.

Three limits bound what this thesis can say about strain. Spatial differentiation can amplify flow-estimation noise, and \textbf{no dedicated strain post-filter is applied}, departing from the explicit filtering of Liong et al.~\cite{liong2014a}; Farneb\"ack's internal smoothing is not a separately evaluated strain filter. The shear weighting is the untested choice of \autoref{sec:strain-divergence}. Above all, \textbf{strain is present in all twelve evaluated configurations}, so its marginal contribution is not measured at all: a strain-on/strain-off comparison is the natural next ablation, and neither published gains from a different descriptor nor the performance of the complete system substitutes for it. A deterministic transformation of $u$ and $v$ adds no independent observation; its potential value lies in making deformation directly accessible to the classifier, and whether that helps must be measured rather than inferred from its construction.


% =============================================================================
\section{Temporal Normalisation: Fixed-Length Sampling}
\label{sec:tim-review}
% =============================================================================

Flow and strain fix \textit{what} is computed from each frame pair; temporal normalisation fixes \textit{which} pairs, and how many. The Temporal Interpolation Model (TIM) is credited in this corpus to Zhou et al. by Li et al.~\cite{liX2013,liX2018} and to Pfister et al. by Lu et al.~\cite{lu2015}, and is described here only through the reviewed papers that implement it.

Normalisation is necessary because any batched network requires a fixed input shape and the corpus supplies no constant, spanning a four-fold range of sequence length; and because duration is itself a nuisance variable --- Lu et al.~\cite{lu2015} normalise to remove ``temporal fluctuations of micro-expressions, which may introduce noise irrelevant to classification''. They also give the corpus's most compact description of TIM: ``a manifold-based interpolation method, which builds a low-dimensional embedding of an image sequence, and then interpolates a curve in the low-dimensional space. The interpolated frames are mapped back to the original high-dimensional space.'' The property that matters downstream is that TIM \textbf{synthesises frames that were never recorded}, having been proposed ``to generate sufficient frames from an image sequence'' --- an upsampler rather than only a length-standardiser.

\textbf{Ten frames is a convention, not a finding, and the studies that test it disagree.} TIM appears throughout the corpus as an unremarked preprocessing step, almost always at ten frames~\cite{liX2013,liX2018,lu2015}. Li et al.~\cite{liX2018} sweep interpolation length from 10 to 80 on the SMIC variants with the descriptor held fixed, and report that ``the best performance for all three datasets is achieved with interpolation to 10 frames'', because ``in longer sequences, changes along the time-dimension are diluted''. Ben et al.~\cite{ben2021} reach the opposite conclusion: holding alignment, descriptor and classifier fixed, they interpolate to 10 through 110 frames by both Newton interpolation and TIM on MMEW, finding Newton peaking at 33.3\,\% at 60 frames and TIM at \textbf{38.9\,\%} at 30 or 60, and adopt TIM at 60. Accuracy is non-monotonic in length in both. Xu et al.~\cite{xu2017} add the only multi-corpus sweep --- 10, 15 and 20 frames across CASME, CASME II and both SMIC subsets --- and depart from convention by using \textbf{linear} interpolation ``instead of Temporal Interpolation Model because in a high frame-rate scenario, linear interpolation suffices to characterize the motion pattern with less error''.

Two caveats limit what transfers. Ben et al.'s experiment uses MMEW with LBP-TOP, not CASME II and not a learned model, and runs largely in the \textit{upsampling} regime, so an optimum measured while synthesising frames need not hold when discarding them. And \textbf{frame synthesis is in tension with motion analysis}: where the downstream representation is appearance-based, synthesised frames are benign, but where it is motion, synthesised frames yield synthesised displacements. Xu et al.~\cite{xu2017} are the only reviewed paper to raise this, and they assert it rather than measure it. No reviewed study compares the two interpolators against a motion representation.

\subsection{Implications for this research}
\label{sec:tim-implications}

\textbf{A declaration about what this pipeline implements.} The module performing this step is named for the Temporal Interpolation Model, but what it computes is \textbf{uniform index sampling} of existing frames (\autoref{sec:three-channel-tensor}). No embedding is built, no curve is interpolated, no frame is synthesised. That is a materially different operation from TIM, and it is named accurately here rather than in the pipeline's own terms. It is defensible: because the downstream representation is flow between adjacent sampled frames, interpolation would manufacture displacement fields for motion that never occurred, whereas subsampling discards temporal resolution without fabricating it. What is given up is TIM's ability to \textit{lengthen} short clips.

The length is fixed at \textbf{33 frames sampled, giving 32 flow pairs}, bounded by the annotated onset and offset so the window spans the expression rather than surrounding neutral frames. That is close to the lower of Ben et al.'s two reported optima, but the correspondence should not be overstated: their 30 frames are \textit{interpolated} frames under a hand-crafted descriptor on a different corpus, which is not direct validation of 32 sampled flow fields under a learned model here. The length is not swept, and a sweep is declared as further work.

Three quantitative consequences follow, and they constrain the magnification conclusion rather than merely describing the sampler.

\textbf{Effective temporal resolution falls by a clip-dependent factor.} Sampling 33 frames from a clip of $N$ frames recorded at 200~fps gives an effective rate of $6600/N$ fps: \textbf{100 fps} at the median $N = 66$, which is SMIC's recording rate and half what CASME II was built to provide, and roughly \textbf{52 fps} for the longest clip, below the original CASME database's 60~fps~\cite{yan2014}. Eighty-one of the 156 clips are downsampled by a factor of two or more, which gives back part of what the 200~fps capture rate bought. No reviewed paper reports this cost.

\textbf{Magnification is applied after subsampling, which shifts the realised pass-band.} Magnification runs on the 33 sampled frames at the nominal 200~fps, but those frames span the clip's full duration, so the true inter-frame interval is $N/(200 \times 32)$ seconds rather than $1/200$. At the median, a band specified as 5--25 Hz corresponds to roughly 2.5--12.5 Hz physically, and the discrepancy grows with clip length. The realised band is therefore neither the 5--25 Hz that \autoref{sec:evm-implications} derives nor constant across clips. This is a declared divergence; magnifying before subsampling is declared as further work.

\textbf{Three clips are shorter than the sampling window.} Three of the 156 contain 31 frames against a target of 33, so sampling repeats indices and yields near-zero flow at those positions.

The unresolved question this leaves is what temporal normalisation costs a representation made of motion rather than appearance. Three studies sweep interpolation length and give three answers, only one of them touching CASME II and none using a learned temporal model. This thesis takes the position that follows from its own representation and quantifies what that costs in effective frame rate, but does not sweep the length, compare interpolators, or separate the effect of sequence length from that of the sampling regime.


% =============================================================================
\section{The Learned Spatial Backbone: Shallow Convolutional Networks}
\label{sec:cnn-review}
% =============================================================================

By 2019 the field had converged on optical flow as the input; the open question was what should consume it. A learned backbone promises what hand-crafted descriptors cannot~\cite{liong2019b}, but runs into the small-data constraint of \autoref{sec:corpus-threats}.

\textbf{The term ``3D CNN'' covers three materially different architectures in this literature.} STSTNet~\cite{liong2019b}, the canonical shallow model, calls itself a ``Shallow Triple Stream Three-dimensional CNN'', but its input is a single $28 \times 28 \times 3$ tensor whose third dimension holds horizontal flow, vertical flow and strain --- \textbf{not time}. Its three parallel single-layer streams, with 3, 5 and 8 kernels, supplement ``the small scale input data \ldots to avoid the problem of underfitting''; the clip was collapsed to one frame pair upstream. STRCN~\cite{xia2020a} by contrast extends convolutional connectivity genuinely ``across temporal domain'' with recurrent convolutional layers. Zhao et al.~\cite{zhaoS2021} apply a C3D-style network to a ten-frame flow sequence, so the temporal axis is present but short. Results reported across the corpus under the same label are therefore not like with like.

\textbf{The central controlled finding is that shrinking both model and input helps.} Xia et al.~\cite{xia2020b} treat depth and input resolution as controlled variables, building four models from a single convolutional layer to one convolutional plus three recurrent convolutional layers, and evaluating each across input resolutions from 20 to 300 pixels on optical flow. Deeper is worse: models 3 and 4 ``have worse recognition performances compared to shallower \ldots models'', even though all four are far shallower than ResNet-101. Higher resolution is worse with a threshold: performance ``will be degraded with larger input resolutions'' and for the deeper models ``begins to decrease dramatically when the resolution is bigger than \textbf{100 $\times$ 100}'' --- though shallow models are ``basically robust to the change of input resolution''. That qualification matters for this thesis, and so does the setting: Xia et al.'s study and STSTNet are both developed for the MEGC composite regime, which supplies substantially more training data per fold than a single corpus (\autoref{sec:megc-metrics}), so conclusions about affordable capacity understate the constraint here. Xia et al.~\cite{xia2020b} also take a third route to capacity control that this thesis follows --- developing ``three parameter-free modules \ldots integrate with RCN without increasing any learnable parameters'', so representational power grows while overfitting risk does not.

\textbf{No study isolates the backbone itself.} The corpus is not short of matched comparisons: it sets descriptor-only methods against learned ones on the same folds, and Xia et al.~\cite{xia2020b} vary individual architectural modules against a common baseline with every other parameter held constant. But those comparisons all \textit{add} to a backbone that stays in place. None removes the learned spatial stem entirely and passes the motion representation straight to a classifier or temporal model, so what the stem contributes against what it costs is nowhere measured.

\subsection{Implications for this research}
\label{sec:cnn-implications}

The component ablated here is a \textbf{three-stream shallow convolutional backbone}, one unshared two-layer stream per input channel, carrying \textbf{14,544 parameters} --- shallow by the corpus's standards, though roughly nine times STSTNet's. Where STSTNet's streams differ in filter count over the same input, these differ in input channel. Removing it leaves not an empty slot but spatial pooling followed by a learned linear projection, so the no-backbone route still contains learned input processing.

\textbf{A declaration about what this component is.} Every kernel in the backbone is spatial only: no filter spans two time steps, so despite the \texttt{Conv3d} implementation it is a \textbf{purely spatial} feature extractor (\autoref{sec:conv3d}). The ablation therefore does not test spatio-temporal convolution; it tests a learned spatial stem, and whatever it measures is evidence about \textit{this} stem --- spatial-only, at 224~$\times$~224, on flow and strain input --- and \textbf{not} evidence that convolution over time is unhelpful for micro-expression recognition, which this study never evaluated. In the field's own terms the choice is defensible: STSTNet, one of the strongest shallow baselines on the CASME II subset, likewise performs no temporal convolution.

\textbf{Input resolution is the clearest departure from Xia et al.'s guidance.} The backbone is fed at \textbf{224 $\times$ 224}, more than twice their threshold. This is not a departure from the corpus as a whole --- Bai et al.~\cite{bai2021} resize to 224 $\times$ 224 without alignment or cropping, as here; Li, Huang and Zhao~\cite{liY2021} resize faces to the same; and Xia et al.'s own earlier work~\cite{xia2020a} runs one stream at 300 $\times$ 245 --- but it departs from the specific threshold above, which is a composite-database result that bites hardest on models deeper than this one. Since this stem is shallow, the reviewed evidence predicts a smaller resolution effect than the deep-model figures suggest. The choice is also the expensive one: a 224 $\times$ 224 map carries sixty-four times the spatial elements of STSTNet's input, and the backbone convolves it at every one of the 32 time steps. Any effect measured for the stem is therefore \textbf{a resolution result as much as an architecture result}, and the reviewed evidence does not determine which mechanism dominates; re-running the backbone arm at a lower resolution is the first item of declared further work because it is cheap and would settle the question, not because the corpus implies a particular outcome.


% =============================================================================
\section{Parameter-Free Attention: SimAM}
\label{sec:simam-review}
% =============================================================================

SimAM re-weights convolutional feature maps without adding learned parameters. Yang et al.~\cite{yang2021} supply the mechanism and its image-recognition evidence; Xia et al.~\cite{xia2020b} motivate parameter-free modules in composite-database micro-expression recognition, supporting the design principle rather than validating SimAM on this task. The energy formulation and gate are given in \autoref{sec:simam}.

Yang et al.~\cite{yang2021} motivate attention through spatial suppression: neurons distinguishable from their surroundings are assigned greater importance, expressed as the linear separability of one activation from the others in its channel. This is a prior about activation statistics, not a detector of action units or minority classes --- a distinctive response may encode useful deformation, noise or an artefact --- so its suitability for flow-derived features requires evaluation rather than following from the neuroscience motivation. Two implementation details are owned here rather than by \autoref{ch:background}: the shared variance printed with the paper's equation~(5) divides by $M$ whereas its Figure~3 reference code divides by $M-1$, and this thesis follows the code; and the module's ``three-dimensional weights'' describe $C \times H \times W$ with $M = HW$ per frame, not a temporal volume, so extending the context across frames is an additional design decision rather than part of the original formulation.

Yang et al. evaluate SimAM with several network families on CIFAR and ImageNet, reporting competitive performance against learned attention modules, and their timing results show that unchanged parameter totals do not imply unchanged inference speed. For CIFAR they sweep $\lambda$ from $10^{-1}$ to $10^{-6}$ and select $10^{-4}$; they repeat the selection for ImageNet and choose $0.1$. Thus $10^{-4}$ is a dataset-selected operating point, not a universal setting justified by the formula.

\textbf{The limits of transfer are specific.} The source experiments concern natural images, different training-set sizes and learned feature distributions. Here the input is flow and strain, so distinctiveness need not bear the same relationship to useful evidence. The image experiments also do not decide whether video statistics should be computed per frame or over a clip. And the absence of learned attention weights does not establish that runtime, memory or overfitting concerns disappear. These are empirical limits on transferring the result, not reasons to assume success or failure.

\subsection{Implications for this research}
\label{sec:simam-implications}

SimAM is applied separately to each convolutional stream before concatenation, introducing no trainable weights or biases anywhere. Two choices diverge from the source. Mean and variance span $(T,H,W)$ within each channel with divisor $THW - 1$, rather than one $(H,W)$ plane, so an activation is judged distinctive relative to its whole clip rather than its own frame; this could emphasise a brief contraction, but does not guarantee low weights at onset or high weights at the apex, since the score depends on squared deviation and not on expression labels. And $\lambda = 10^{-4}$ is retained without re-tuning on CASME II. Neither context choice is compared experimentally.

SimAM is also enabled only alongside the convolutional stem in this model, which describes this implementation rather than a mathematical prohibition. Four of the sixteen candidate configurations are therefore excluded, leaving twelve evaluated and \textbf{four SimAM matched pairs}, and the measured effect is conditional on the stem being present. Retention should be judged against both recognition performance and computational cost: parameter count alone cannot settle it, because the module is cheap rather than free.

\subsection{Relevance and scope}
\label{sec:simam-gap}

The study evaluates SimAM in a motion-input, small-corpus setting through four matched comparisons under complete leave-one-subject-out, testing the implemented module as a whole rather than establishing the benefit of each design choice. Three follow-ups remain: whether the distinctiveness prior suits flow-derived features, re-tuning $\lambda$, and comparing per-frame with whole-clip statistics. Each requires new evaluations despite the module introducing no learned weights.


% =============================================================================
\section{The Temporal Encoder: Transformers}
\label{sec:transformer-review}
% =============================================================================

The temporal encoder is motivated by Zhang et al.'s Short and Long Range Relation Based Spatio-Temporal Transformer (SLSTT)~\cite{zhang2022} and its Vision Transformer foundation~\cite{dosovitskiy2021}. The component evaluated here differs substantially from published SLSTT, particularly in the axis over which self-attention operates, and distinguishing the two is the purpose of this section.

A micro-expression evolves through onset, peak and relaxation, so retaining a motion sequence permits a model to relate separated stages rather than compressing them immediately into one descriptor; Zhang et al.~\cite{zhang2022} motivate their design through the need to capture both local and global patterns. This thesis compares a temporal transformer against mean pooling over the same sequence: mean pooling discards order, whereas positional information and self-attention allow relations among frame-pair features to be modelled. That motivates the comparison without establishing in advance which generalises better. ViT itself~\cite{dosovitskiy2021} splits an image into patches, projects them into tokens, adds learned positional embeddings and relates them by attention; Dosovitskiy et al. also show that without the locality and translation-equivariance biases of convolution, image transformers can perform poorly when trained with insufficient data, while large-scale pre-training improves them. That is a caution for a 156-clip study, not a universal requirement that every transformer be pre-trained.

\textbf{Published SLSTT combines three stages.} First, \textbf{onset-referenced optical flow} describes sampled frames relative to the onset rather than to their immediate predecessors, which Zhang et al. argue gives a more structured trajectory towards and away from the apex; their experiments use it throughout, without a matched short-term-flow ablation establishing its independent gain. Second, a \textbf{spatial ViT encoder} processes each flow field as patches, using ImageNet-pre-trained ViT-B/16 with twelve layers at $384 \times 384$. Third, an \textbf{LSTM temporal aggregator} combines the per-frame features, with a mean aggregator supplying the comparison variant. Self-attention is spatial in this architecture; recurrence handles time.

\subsection{The aggregation comparison, and what it can support}
\label{sec:slstt-results}

The comparison relevant here is internal to SLSTT and holds the spatial transformer fixed in both arms: CASME II UF1 rises from \textbf{0.844} with the mean aggregator to \textbf{0.901} with the LSTM, a difference of 0.057~\cite{zhang2022}. That supports evaluating aggregation choices within that framework. It is \textbf{not} an estimate of what an LSTM would add to this thesis, where mean pooling follows \emph{temporal} self-attention rather than independent spatial encoding, and it isolates an aggregator rather than the presence of a transformer, since both variants contain one.

Protocol labels need care. Zhang et al. describe three-class leave-one-subject-out evaluations on the databases separately and together under a Composite Database Evaluation heading, so that heading alone does not establish that every CASME II entry uses pooled cross-database training; they also report a distinct Sole Database Evaluation using original labels, including five CASME II classes~\cite{zhang2022}. Neither setting should be equated automatically with this thesis's 156-clip three-class pool. Published SLSTT also retains a pre-trained spatial encoder, so its results do not isolate the benefit of a transformer trained from scratch on this study's motion representation.

\subsection{Implementation and six declared divergences}
\label{sec:transformer-implications}

The evaluated component receives 32 tokens of width 96, from either the convolutional stem or spatial pooling followed by a learned linear projection. The encoder uses four pre-norm layers, eight heads, feed-forward width 256 and dropout 0.1; fixed sinusoidal positions precede attention and mean pooling follows it. With the encoder disabled the input tokens are averaged directly.

\begin{table}[htbp]
	\centering
	\small
	\begin{tabular}{@{}p{0.23\linewidth} p{0.34\linewidth} p{0.35\linewidth}@{}}
		\hline
		\textbf{Property} & \textbf{Published SLSTT~\cite{zhang2022}} & \textbf{This thesis} \\
		\hline
		Attention axis   & Spatial patches within a frame              & Temporal tokens within a clip \\
		Aggregation      & LSTM over spatially encoded frames          & Mean after temporal self-attention \\
		Initialisation   & ImageNet-pre-trained ViT-B/16               & Trained from scratch \\
		Capacity         & Twelve layers, width 768; $384\times384$ in & Four layers, width 96; 32 tokens \\
		Positions        & Learned spatial embeddings                  & Fixed sinusoidal temporal encoding \\
		Input flow       & Relative to onset                           & Relative to the preceding sampled frame \\
		\hline
	\end{tabular}
	\caption[Divergences from published SLSTT]{Six material divergences from published SLSTT. The component evaluated here is a small temporal transformer encoder, not a reproduction of the published architecture.}
	\label{tab:slstt-divergences}
\end{table}

Six matched pairs evaluate the encoder with other component choices fixed within each pair. They test the whole encoder --- positional encoding, feed-forward layers and normalisation included --- against direct mean pooling, and do not isolate self-attention from those accompanying operations or compare it against another learned temporal model. Three follow-ups follow from \autoref{tab:slstt-divergences}: testing an LSTM aggregator, comparing onset-referenced with consecutive-frame flow, and replacing fixed positions with learned embeddings. All three remain untested, and changing the input representation requires re-training before any benefit can be claimed. What this study can establish is a matched comparison of the implemented encoder against an order-blind baseline within this pipeline; it neither validates published SLSTT nor establishes a general ranking of attention, recurrence and convolution.


% =============================================================================
\section{Training Under Severe Class Imbalance}
\label{sec:imbalance-review}
% =============================================================================

The 99 : 32 : 25 distribution of \autoref{sec:label-taxonomy} is a property of what can be elicited in a laboratory, not a sampling accident, so it cannot be corrected by collecting more data from the same paradigm. The corpus corrects for it at three points --- the data, the loss and the metric --- and this thesis uses all three. The mechanisms are given in \autoref{sec:focal-loss} and \autoref{sec:balanced-sampling}.

\textbf{At the loss.} Zhao et al.~\cite{zhaoS2021} bring focal loss into micro-expression recognition ``to alleviate the inefficient model training caused by the class imbalance in the MEs Datasets'', adopting the formulation developed for one-stage object detection. The division of labour inside it is what matters: $\gamma$ re-weights by \textit{difficulty} and is treated as a constant, ``set as 2 in practice'', while $\alpha$ re-weights by \textit{class frequency} and is ``treated as a hyper-parameter to set by cross validation''. Conflating the two is the error discussed below. Zhao et al. compare focal loss against cross-entropy across a sweep of $\alpha$, but hold $\gamma$ fixed, so the focusing exponent remains unvalidated on this task.

\textbf{At the data.} Xia et al.~\cite{xia2020a} pair ``temporal data augmentation strategies as well as a balanced loss \ldots jointly'' to overcome ``limited and imbalanced training samples''. Their own ablation removes each correction in turn and reports the balanced loss improving performance only ``slightly'' --- at most about two points, and in one cell costing accuracy. This is the corpus's only controlled look at composing two corrections, and no reviewed paper reports the pairing failing outright.

\textbf{At the metric.} \autoref{sec:megc-metrics} defines the pooled macro-averaged measure. Liong et al.~\cite{liong2016} add the argument from the subject side, reporting F1, precision and recall because these ``provide a more meaningful perspective than accuracy rates when the datasets used are naturally imbalanced since each subject has a different number of video samples.'' Both arguments compound here, where classes are skewed \textbf{and} per-subject clip counts range from 1 to 33.

One further gap: \textbf{label smoothing is named in the corpus but never examined on this task.} Dosovitskiy et al.~\cite{dosovitskiy2021} list it among the regularisers they tune, but no reviewed micro-expression paper uses it or reports what it does under class skew.

\subsection{Implications for this research}
\label{sec:imbalance-double-correction}

Focal loss is adopted with \textbf{$\gamma = 2.0$}, matching Zhao et al.'s reported setting. The class-frequency term is implemented as an optional per-class $\alpha$ vector but \textbf{switched off at run time}, because frequency correction is assigned to a weighted sampler active in every configuration. Label smoothing of 0.05 is applied inside the focal loss, declared as a choice this corpus supports neither way. The primary metric is pooled macro F1, with accuracy quoted only alongside the always-Negative floor.

\textbf{What this project contributes is a negative result about composing corrections.} The configuration is resolved at run time by a single rule: class weighting in the loss is enabled only when the balanced sampler is \textit{not} active. That rule exists because of two failures in opposite directions, both recoverable from this project's own run history. In an earlier code era the loss carried inverse-frequency class weights and no sampler existed; under that single correction the proposed configuration abandoned the \textbf{Positive} class outright, a failure initially misdiagnosed as over-capacity overfitting. A later era added the sampler while leaving the loss weights unconditional, so both corrections ran at once; the collapse then inverted, and configurations abandoned the \textbf{Negative} majority class instead, one predicting a single class for every clip. With the guard in place, no configuration in the reported study abandons any class.

This is worth reporting because no reviewed paper warns that combining corrections can invert the intended effect rather than merely blunt it. It must be stated carefully, however: the two collapses occurred in code eras differing from the reported run in more than the correction rule, and the final sweep contains no matched arm with the guard disabled. \textbf{The guard's causal benefit is therefore not isolated}, and the evidence does not establish that correcting at exactly one point is universally best --- only that correcting twice, in this pipeline, produced a class collapse that correcting once did not.

A second declaration follows. Because every configuration shares the same sampler, loss and $\gamma$, this thesis measures \textbf{no} effect attributable to the imbalance treatment: it is a constant, not a factor. The claim supported is that the four architectural components were compared under a training regime that does not abandon minority classes, not that this regime is optimal.

\subsection{Relevance and scope}
\label{sec:imbalance-gap}

The corrections themselves are well established. What the corpus does not settle is how two corrections compose, what $\gamma$ should be on this task, and what over-correction looks like when it fails. This thesis contributes to the first and third by documenting two opposite class-collapse failures and the run-time rule that resolved them. On the second it offers no improvement, since $\gamma = 2.0$ and label smoothing of 0.05 are held constant across all twelve configurations; sweeping either is declared as further work.


% =============================================================================
\section{Synthesis and Research Gap}
\label{sec:synthesis-and-gap}
% =============================================================================

This section consolidates the preceding relevance statements into one research gap, and collects in one place every declared divergence between the system built here and the literature it draws on.

Read as a whole, the reviewed literature settles a good deal. The corpora are small, demographically narrow, laboratory-bound and irreducibly imbalanced by what can be elicited, and CASME II remains the primary single-corpus benchmark under its native leave-one-subject-out protocol. Motion rather than appearance is the appropriate input; a strain field adds a further cue, though on CASME II only weakly and under a protocol that is not subject-disjoint; and both the magnification factor and the sequence length have interior optima located only by sweep. At this data size, capability bought without parameters outperforms capability bought with them, and self-attention is the strongest-performing mechanism reported. Accuracy alone is inadequate under skew and must give way to macro-averaged measures accumulated across folds before averaging.

\subsection{Four limitations that recur across every topic}
\label{sec:four-limitations}

The per-section statements above are not independent observations. Four structural problems generate most of them.

\textbf{(a) Components are isolated one at a time, inside one proposed design, and never jointly.} Matched-pair measurement is not absent from the corpus: \autoref{tab:evm-effect} isolates magnification, \autoref{tab:strain-performance} isolates strain against a flow control, Xia et al.~\cite{xia2020b} compare a recurrent convolutional network against three variants each adding one module with every other parameter held constant, and SLSTT's Mean-versus-LSTM comparison isolates its aggregator (\autoref{sec:slstt-results}). What no reviewed study does is vary \textit{several} components together, so that their interactions can be read. Every ablation is internal to one proposed design and moves one factor from that design's own defaults; across papers, systems are compared whole, so a difference in the reported number cannot be attributed to any one stage. There is therefore no basis for deciding which learned stages are load-bearing, nor whether one component's value depends on another's presence.

\textbf{(b) The evaluation protocol is unstable and under-reported.} Sample counts, label-set size and the choice among leave-one-subject-out, leave-one-video-out and leave-one-sample-out vary inconsistently, and the field's most-quoted baseline is described three ways by its own sources (\autoref{sec:original-baseline}). Given this corpus's fold composition these are not bookkeeping differences: they change what the numbers mean.

\textbf{(c) The strongest results depend on regimes a single-corpus study cannot access.} The best CASME II figures come from composite-database training (\autoref{sec:megc-metrics}) and, for SLSTT, ImageNet pre-training --- neither a property of the architecture being credited. Almost no evidence exists about what these components do when trained from scratch on one small corpus, which is the situation most practitioners are in.

\textbf{(d) Preprocessing parameters have interior optima and are transferred without re-validation.} The magnification factor, the sequence length and the input resolution all rise then fall, and none has a principled selection rule. Values are carried between pipelines whose downstream models differ entirely: $\alpha = 30$ moves from a single-frame appearance encoder to other settings, ten-frame interpolation propagates from one baseline paper to the whole field, and $\lambda = 10^{-4}$ transfers from CIFAR though its own authors re-tuned it for ImageNet.

\subsection{The research gap}
\label{sec:research-gap}

Taking (a) to (d) together:

\textbf{The reviewed literature has established a standard pipeline for micro-expression recognition --- magnify, compute motion, normalise duration, extract spatially, model temporally --- and contains matched ablations of individual stages, but no study varies several of those stages jointly so that their relative contributions and interactions can be read. Because components are varied one at a time inside a single proposed system, under protocols that differ in unstated ways, and predominantly in training regimes richer than a single small corpus provides, it is not currently possible to say from this literature which parts of that pipeline earn their place, which are inherited convention, and which depend on the presence of another.}

This matters practically as well as scientifically. A practitioner assembling this pipeline today has no published basis for allocating effort: for judging whether the most computationally expensive stage is the most valuable, whether a component with a negligible average effect may still matter in a specific failure, or whether a ranking read off an incomplete evaluation of a corpus this size is stable at all.

\subsection{How this thesis addresses the gap, and what it does not}
\label{sec:addressing-gap}

\textbf{What it does.} This thesis treats the pipeline as a factorial experiment rather than a proposal: four components varied independently across every architecturally valid cell of a 2\textsuperscript{4} matrix, twelve configurations, each evaluated under \textbf{complete 25-fold leave-one-subject-out on CASME II} with every non-varied factor held identical. That supplies the joint measurement (a) identifies as missing, reporting the protocol in full addresses (b) for this study at least, and training from scratch is the regime (c) identifies as unevidenced. Each component's marginal contribution is read from pairs of runs differing in one factor alone, and the distribution of those four effects is the substantive answer to the gap. Two further contributions follow from the review rather than the experiment: a \textbf{repaired measurement}, where an early data-routing defect left the magnification switch inert and the review's evidence of a clear, repeatable magnification effect is what made bit-identical results detectable as a defect rather than a null (\autoref{sec:evm-implications}); and \textbf{two opposite class-collapse failures} that no reviewed paper reports (\autoref{sec:imbalance-double-correction}).

\textbf{What it does not do.} Three of the four limitations are only partially addressed. On (b), this study reports its own protocol completely but cannot repair the field's inconsistency, and its own N = 156 differs from MEGC's 145. On (c), it establishes what these components do from scratch on one corpus, which is a different question from --- not an answer to --- how they behave under composite training. On (d), the thesis inherits parameter values rather than sweeping them: $\alpha = 10$, T = 32, $\lambda = 10^{-4}$ and $\gamma = 2.0$ are fixed throughout, so the ablation measures architecture at one point in preprocessing space rather than across it.

\subsection{A consolidated ledger of declared divergences}
\label{sec:divergence-ledger}

Reviewing the literature closely enough to write this chapter surfaced twelve points at which the implemented system departs from the work it draws on. Several were discovered only through the review. They are collected in \autoref{tab:divergence-ledger} so that no result rests on an unstated assumption.

\begin{table}[htbp]
	\centering
	\footnotesize
	\begin{tabular}{@{}p{0.85cm} p{6.0cm} p{6.6cm}@{}}
		\hline
		\textbf{\S} & \textbf{Divergence} & \textbf{Consequence} \\
		\hline
		\ref{sec:label-taxonomy} & Working set is N = 156, including sadness and fear in Negative; MEGC's CASME II subset is N = 145 & Comparisons against challenge figures are not exactly like-for-like \\
		\ref{sec:casme2-construction} & The corpus's own registered release is not used; the pipeline reads the original recorded frames, unregistered and uncropped & No landmarking or alignment of any kind precedes the network; tolerance to pose and scale must be learned from 156 clips \\
		\ref{sec:evm-implications} & Temporal band is the unnarrowed 5--25 Hz duration rule, where Bai et al. narrowed it to 15--25 Hz as noisy & The pass-band admits the low-frequency range they rejected \\
		\ref{sec:evm-implications} & Band-pass is realised by an FFT coefficient mask, not the Butterworth filter of the reviewed description & A different filter shape, with different transition behaviour at the band edges \\
		\ref{sec:strain-divergence} & Strain magnitude counts shear once, not twice as in the published Frobenius norm & Shear weighted lower by $\sqrt{2}$ relative to normal strain; untested \\
		\ref{sec:strain-implications} & No dedicated strain filtering, where the literature applies Wiener and Gaussian filters & Spatial differentiation can amplify flow-estimation noise, and nothing downstream suppresses it \\
		\ref{sec:tim-implications} & Temporal normalisation is uniform index sampling, \textbf{not} manifold-based TIM; no frames are synthesised & Avoids fabricating motion; forfeits the ability to lengthen short clips \\
		\ref{sec:tim-implications} & Magnification is applied \textbf{after} subsampling, at a nominal 200 fps & Realised pass-band is clip-dependent and below the intended 5--25 Hz \\
		\ref{sec:cnn-implications} & Backbone kernels are spatial only --- \textbf{no temporal mixing} & The ablation tests a learned \textit{spatial} stem, not spatio-temporal convolution \\
		\ref{sec:cnn-implications} & Backbone input is 224 $\times$ 224, against Xia et al.'s guidance of $\leq$ 100 $\times$ 100 --- though Bai et al. and Li, Huang and Zhao both use 224 $\times$ 224 & Whatever the stem is measured to contribute may be a resolution effect as much as an architecture effect \\
		\ref{sec:simam-implications} & SimAM statistics pooled over the whole clip rather than per frame as in the source; $\lambda$ not re-tuned & Neurons judged against the clip, not the frame; defensible but untested \\
		\ref{sec:transformer-implications} & Six divergences from published SLSTT (temporal not spatial attention, mean not LSTM aggregation, from scratch, smaller, sinusoidal encoding, short-term flow) & The temporal encoder is \textbf{not} a reproduction of SLSTT \\
		\ref{sec:imbalance-double-correction} & Label smoothing of 0.05, named in the corpus only by Dosovitskiy et al. and never examined on this task & A choice the micro-expression corpus supports neither way, held constant and unmeasured \\
		\hline
	\end{tabular}
	\caption[Every declared departure from the reviewed literature]{Every declared departure from the reviewed literature, with its consequence. The first column resolves to the subsection that owns each divergence.}
	\label{tab:divergence-ledger}
\end{table}

Three of these --- the strain formula, the magnification ordering and the input resolution --- were identified by the review itself rather than being known design decisions, and each is declared as proposed further work. None invalidates the ablation, since every configuration shares the same choices. What they constrain is the \textit{wording} of the conclusions: what is measured is the contribution of a spatial stem at 224 $\times$ 224 and of temporal self-attention in this pipeline, not a verdict on convolution over time or on SLSTT. The limits of what can be claimed are therefore fixed before any number is reported --- a small, single-seed, demographically narrow corpus measured against richer published regimes, with these divergences standing throughout.
 IN PLACE OF chapter3.tex and
% chapter3_shortened.tex. Never include two of them: they share labels.
%
% WHAT CHANGED
%   * Repeated mechanism/derivation passages removed: they are owned by
%     Chapter 2 (magnification, brightness constancy, strain tensor, SimAM
%     energy, focal loss, F1). This chapter now carries published evidence and
%     transfer limits only, and contains NO display mathematics.
%   * The nine per-section "research gap" endings are replaced by short
%     relevance statements, with one qualified synthesis at the end.
%   * Local decision tables (EVM/flow/strain/TIM/CNN/SimAM/imbalance) removed:
%     they restated Chapter 2 or 4 and were then restated again in the ledger.
%   * 17 tables -> 6. 2 figures -> 1. The dataset-composition figure is gone
%     (its content is now one prose sentence); the fold-composition figure is
%     retained because the pooled-versus-fold argument depends on it.
%
% THE LEDGER'S FIRST COLUMN IS NOW \ref{}, NOT HARD-CODED NUMBERS.
% The twelve divergence rows previously cited 3.1.2 ... 3.9.7 literally, which
% any restructuring silently falsifies. They now resolve automatically.
%
% LABELS: every label referenced from introduction.tex, chapter4_shortened.tex,
% chapter5_shortened.tex or chapter6.tex is preserved (24 of them), plus those
% used by the unshortened chapters, so either manuscript state resolves.
% Verified by scratchpad scan; re-run it after any further edit.
%
% NOTE: the old header warned that Chapter 2 cites Chapter 3 subsections BY
% NUMBER. That is no longer true -- the shortened Chapter 2 has none -- so
% subsection positions here are free. The only by-number dependency was the
% ledger's own first column, now removed.
% =============================================================================

\chapter{Literature Review}
\label{ch:literature}

This chapter reviews the literature bearing on each component of the recognition pipeline evaluated in this thesis: the corpus every experiment runs on, the four techniques placed under ablation, the input transformations they operate on, and the training regime held constant across all of them. Each section states what the relevant studies found, what their evidence does not establish, and what follows for this experiment. Mechanisms already derived in \autoref{ch:background} are not derived again. \autoref{sec:synthesis-and-gap} consolidates the closing statements into one research gap and collects every declared departure from the reviewed methods.


% =============================================================================
\section{The Evaluation Corpus: CASME II and the Spontaneous Benchmark Family}
\label{sec:corpus}
% =============================================================================

The corpus is treated first because on a database of this size several decisions usually presented as modelling choices are in fact forced by properties of the data.

\subsection{Why a dedicated spontaneous corpus was necessary}
\label{sec:why-spontaneous}

Automatic expression recognition matured on large corpora of \textit{posed} expressions, reporting accuracies above 90\,\% on the six basic emotions by the early 2010s~\cite{yan2014}. Micro-expression research could not follow, because the phenomenon is involuntary and cannot be performed on instruction. Two problems motivated CASME II directly.

The first is \textbf{construct validity}: Yan et al.~\cite{yan2014} and Li et al.~\cite{liX2013} both observe that deliberately produced ``micro-expressions'' differ systematically from spontaneous ones, so a system trained on posed data may be a recogniser of fast deliberate movement rather than of micro-expressions. The second is \textbf{acquisition quality}. Defined jointly by very short duration and very low intensity, the phenomenon sits close to the noise floor of ordinary video; Yan et al.~\cite{yan2014} argue that CASME's 60~fps coarsely samples a movement lasting a fraction of a second --- while marking the benefit of higher rates as unresolved --- and that mains-frequency flicker contaminated earlier recordings. Elicitation was a third difficulty: deception paradigms produced samples contaminated by conversational movement, and the workable paradigm, which CASME II adopts, has participants watch high-arousal video while attempting to keep the face neutral~\cite{liX2013}.

\begin{table}[htbp]
	\centering
	\footnotesize
	\setlength{\tabcolsep}{4pt}
	\begin{tabular}{l c c c c c c}
		\hline
		\textbf{Corpus} & \textbf{Subj.} & \textbf{Micro samples} & \textbf{fps} & \textbf{Spont.} & \textbf{FACS} & \textbf{Ethn. / age} \\
		\hline
		SMIC~\cite{liX2013}              & 16 valid          & 164                & 100          & Yes          & No           & 3 / 26.7 \\
		\textbf{CASME II}~\cite{yan2014} & \textbf{26 valid} & \textbf{247 / 255} & \textbf{200} & \textbf{Yes} & \textbf{Yes} & \textbf{1 / 22.03} \\
		SAMM~\cite{davison2018}          & 32                & 159                & 200          & Yes          & Yes          & 13 / 33.24 \\
		CAS(ME)$^2$~\cite{qu2016}        & 22                & 53 (+250 macro)    & 30           & Yes          & Yes          & 1 / 22.59 \\
		\hline
	\end{tabular}
	\caption[The spontaneous micro-expression corpora]{The spontaneous corpora that define the current context, compiled from Yan et al.~\cite{yan2014}, Davison et al.~\cite{davison2018}, Li et al.~\cite{liX2013} and Qu et al.~\cite{qu2016}. CASME II carries two sample counts: 247 in the paper, 255 in the released coding.}
	\label{tab:spontaneous-corpora}
\end{table}

\subsection{CASME II: construction and annotation}
\label{sec:casme2-construction}

CASME II~\cite{yan2014} recruited thirty-five participants, mean age 22.03 (SD 1.60), recording at \textbf{200 fps} and 640~$\times$~480 without inter-frame compression under flicker-free illumination, giving a facial region of roughly 280~$\times$~340 pixels. Each retained sample carries onset and offset indices, FACS action-unit labels from two independent coders with disagreements arbitrated by discussion, and an emotion label assigned from the coded action units, participant self-report and the eliciting episode's content --- Yan et al. arguing against forcing spontaneous samples into the six basic categories, since a single action unit is frequently ambiguous.

Two specifics propagate into this thesis. The original release applies a normalisation from which this work departs: 68 Active Shape Model landmarks are detected on the \textit{first frame only} and the resulting transform applied unchanged to every subsequent frame; this pipeline reads the original recorded frames instead, unregistered and uncropped. And the published paper reports \textbf{247} micro-expressions while the distributed coding file contains \textbf{255}, a post-publication revision rarely acknowledged and one reason sample counts for ``CASME II'' vary. The present work uses the released 255-sample coding.

\subsection{The label taxonomy and the sparsity it creates}
\label{sec:label-taxonomy}

Yan et al.~\cite{yan2014} provide five emotion categories with action-unit criteria and are candid about the result: ``some types of facial expressions are difficult to elicit in laboratory situations, thus the samples in different categories distributed unequally.'' The released coding is finer, exposing seven raw labels. Counting directly from the label table used here gives happiness 32, disgust 63, surprise 25, repression 27, sadness 7, fear 2, and others 99.

\textbf{Seven-class classification is not viable on this corpus.} Fear is represented by two clips in the entire database and sadness by seven. Under any subject-disjoint protocol --- and subject-disjointness has been this benchmark's norm since its introduction (\autoref{sec:original-baseline}) --- a two-clip class cannot be meaningfully trained or tested, and per-class metrics degenerate.

The field's response is to group by affective valence. The convention originates with SMIC, released with exactly three classes~\cite{liX2013}, and was carried into cross-corpus work by MEGC 2019~\cite{see2019}. Discarding ``others'' and grouping the remainder gives the working set used throughout: \textbf{Negative 99, Positive 32, Surprise 25, N = 156}, a ratio of roughly 4 : 1.3 : 1.

\textbf{A declared divergence from the MEGC convention.} MEGC 2019 reports the CASME II portion of its composite database as \textbf{145} samples, footnoting that CASME II's Negative class covers ``Disgust and Repression''~\cite{see2019} and building its subset from the paper's five-category release, in which sadness and fear are absorbed into ``others''. This thesis uses the seven-label coding and assigns sadness and fear to Negative, following the mapping MEGC gives for \textbf{SAMM} rather than for CASME II. The consequence is an \textbf{eleven-clip difference}: N = 156 against 145, with Negative at 99 rather than 88. Nine of the eleven are the sadness and fear clips; the remaining two are unaccounted for, since disgust and repression together number 90 in the released coding against the 88 the challenge reports --- a residual noted rather than resolved. The direction is mild: the nine added clips all fall in the majority class and come from the categories coders found hardest to assign, so if anything they make this task marginally harder than the challenge's.

\subsection{The original baseline, and what it does not establish}
\label{sec:original-baseline}

Yan et al.~\cite{yan2014} baseline the corpus with LBP-TOP~\cite{zhaoG2007} and a support vector machine, sweeping spatial and temporal radii and excluding $R_T = 1$ deliberately, because at 200~fps the change between adjacent frames is negligible --- an early acknowledgement that a high sampling rate is not automatically an advantage. Five-class accuracies under \textbf{leave-one-subject-out} cluster in the high fifties to low sixties, the best configuration reaching \textbf{63.41\,\%}. Three observations matter.

\textbf{Leave-one-subject-out is native to this benchmark}, not a stricter protocol adopted later, so any evaluation reporting a single random split departs from the benchmark's own convention. \textbf{The baseline is reported as accuracy alone}, and on a five-class problem where ``others'' is 102 of 247 samples the majority-constant classifier scores 41\,\% without modelling anything, so a headline near 63.41\,\% is not by itself evidence of a working system. \textbf{And the sources disagree on the protocol behind that figure}: Yan et al.~\cite{yan2014} state leave-one-subject-out in abstract and results prose, the caption of the table reporting it reads ``leave-one-out cross-validation'' without the subject qualifier, and Li et al.~\cite{liX2018} reproduce the figure marked as leave-one-\textit{sample}-out, a materially easier protocol. The most widely quoted number on this benchmark thus rests on a protocol its own sources describe three ways; this thesis follows the originating paper's prose and records the discrepancy rather than resolving it.

Accuracy also moves by several points across a purely descriptive radius sweep, with no change to data, split or classifier. Where effect sizes of that magnitude arise from descriptor settings, similar differences between architectures cannot be interpreted without a variance estimate --- a limitation this thesis inherits and declares.

\subsection{Sibling corpora, and why this study remains single-corpus}
\label{sec:sibling-corpora}

Three further databases define the context, each isolating a specific limitation of CASME II. \textbf{SMIC}~\cite{liX2013} is the origin of the three-class taxonomy and of fixed-length temporal interpolation, both inherited here; it is not used as a corpus because it is neither FACS-coded nor apex-annotated. \textbf{SAMM}~\cite{davison2018} addresses the CASME family's demographic narrowness with 13 ethnicities against CASME II's one, Davison et al. observing that a corpus drawn from one ethnicity and one narrow age band restricts analysis to ``similar looking participants''; it is stronger on that count but smaller in usable samples and spread across seven categories. \textbf{CAS(ME)$^2$}~\cite{qu2016} pairs macro- and micro-expression samples from the same participants, making it the natural resource for domain transfer --- the strategy behind the strongest MEGC 2019 entry~\cite{liu2019,see2019} --- but records at 30~fps, supplies only 53 micro-expression samples, and is evaluated leave-one-\textit{video}-out, which does not enforce subject-disjointness.

CASME II is used alone because it has the highest temporal resolution of any spontaneous corpus of comparable size, a precondition for magnification in a defined band and for dense frame-to-frame flow; because its onset, apex and offset annotation makes temporal normalisation well-defined rather than heuristic; and because it remains the most widely reported single-corpus benchmark. The cost --- that no claim here extends to other populations, capture conditions or corpora --- is declared rather than argued away.

\subsection{How CASME II is evaluated: MEGC 2019 and the metric convention}
\label{sec:megc-metrics}

The most consequential development in how CASME II is \textit{used} is the Composite Database Evaluation protocol of the Second Facial Micro-Expression Grand Challenge~\cite{see2019}. It merges SMIC, CASME II and SAMM into one composite database of \textbf{442 samples} --- Negative 250, Positive 109, Surprise 83 --- of which CASME II contributes 145, and runs leave-one-subject-out across the composite. Plain accuracy is rejected explicitly on grounds of class imbalance, and two balanced metrics mandated: \textbf{Unweighted F1 (UF1)}, the macro-averaged F1 defined in \autoref{sec:confusion-matrix}, and \textbf{Unweighted Average Recall}, the balanced accuracy.

What matters is not the algebra but the \textbf{construction}: UF1 accumulates true positives, false positives and false negatives per class \textbf{over all folds} \textit{before} any per-class F1 is computed. As \autoref{sec:pooled-vs-fold} establishes, that is a different estimator from the average of per-fold macro F1 scores, and the distinction is not cosmetic here. Because clips are distributed very unequally across subjects and many single-subject folds contain only one of the three classes, a per-fold macro F1 is bounded well below one by fold composition alone (\autoref{fig:fold-composition}). The primary metric of this thesis, \textit{pooled macro F1}, is therefore identical in construction to UF1; the per-fold average appearing in the raw logs is not a metric the field uses.

\begin{figure}[htbp]
	\centering
	\includegraphics[width=\textwidth]{fig3_2_fold_composition.png}
	\caption[Class composition of the 25 leave-one-subject-out folds]{Class composition of the 25 leave-one-subject-out folds, and the per-fold macro-F1 ceiling it implies.}
	\label{fig:fold-composition}
\end{figure}

\textbf{One structural caveat governs the challenge's reported figures}, which range up to 0.7885 composite UF1 and 0.8764 on the CASME II subset~\cite{see2019}. Every one comes from a model trained on the \textit{composite} database --- roughly 442 samples across three corpora and 68 subjects --- then scored on the CASME II portion of the pooled predictions. A CASME II column is therefore not a CASME II-only result: it reports how a cross-corpus model performs on CASME II clips, with approximately \textbf{2.8 times} the training data available in any single-corpus fold, where the experiment conducted here trains each fold on roughly 150 clips. The two regimes are not interchangeable, and single-corpus results should not be compared against those figures without stating the difference.

Two further observations from the challenge bear on this design. The top three submissions all chose \textbf{optical flow} as their representation, a convergence that supports the motion-based input used here; and one submission substituted mid-position frames for annotated apex frames with little loss, suggesting that precise apex localisation may be less critical than assumed.

\subsection{Threats to validity carried by the corpus}
\label{sec:corpus-threats}

The following are not criticisms of CASME II, which remains the best available resource for this task, but constraints any study built on it inherits. Fold and duration statistics are computed from the 156-clip label table used here.

\textbf{Scale and imbalance.} 156 usable clips across three classes, distributed 99 : 32 : 25, small enough that parameter count is itself a risk factor. This motivates parameter-free or low-parameter components wherever a choice exists, and it makes accuracy unusable alone: the always-Negative classifier attains \textbf{0.6346 accuracy} at a macro F1 of only \textbf{0.2588}, so that floor is quoted alongside every result.

\textbf{Subject-level non-uniformity.} \textbf{Subject 17 alone contributes 33 clips, 21\,\% of the working set}, while three subjects contribute one each, so fold sizes vary by a factor of 33; \textbf{ten of the 25 folds contain a single class}, eight contain two, and only seven all three. This forces full 25-fold leave-one-subject-out and is why the pooled metric rather than the fold average is primary. It is also why a single held-out split is not a reliable instrument here: which subjects fall in the test set can move the result by more than any architectural change under study. Subject 18's three clips are all ``others'' and are removed by the three-class filter, so the run has \textbf{25 folds, not 26} --- stated explicitly because fold counts differ between papers using this corpus.

\textbf{Demographic homogeneity and laboratory ecology.} All participants are Chinese with mean age 22.03, named by Davison et al.~\cite{davison2018} as a principal motivation for SAMM; the MEGC organisers saw the effect in reverse, reporting that adaptation from CK+ helped on SMIC and SAMM but ``not so for CASME II, which contain predominantly Chinese subjects''~\cite{see2019}. Recording is frontal under fixed lamps with participants instructed not to move or speak, so almost none of the head motion, illumination variation, occlusion or speech-related movement a deployed system would face is present. Together these motivate a \textbf{motion} representation rather than raw pixels --- appearance being exactly what a 26-subject corpus would let a network memorise --- and they are why no cross-dataset or in-the-wild claim is made anywhere in this thesis.

\textbf{Duration, annotation and version ambiguity.} Onset-to-offset lengths range from \textbf{31 to 126 frames}, median 66 --- at 200~fps, 0.155\,s to 0.63\,s --- so every clip is resampled to a fixed length before training. Labels derive from coded action units, self-report and stimulus content: a construct produced by trained coders under a rubric, taken here as given and recorded as a declared limitation rather than as measurement error. And three sample counts circulate under the name ``CASME II'' --- 247, 255 and 145 --- with both five- and seven-class label sets in use and rarely stated, so N is reported explicitly alongside every figure here.

\subsection{Implications for this research}
\label{sec:corpus-implications}

Two further properties determine settings rather than threatening validity: the \textbf{200 fps} capture rate fixes the magnification pass-band at a matched \textbf{5--25 Hz}, and the baseline's sensitivity to a descriptive hyper-parameter sets the standard of evidence, so effect sizes below the resolution of a 156-clip corpus are reported as unresolved rather than as findings. No confidence interval is available, since the single-seed design and unretained per-clip predictions preclude one.

The consequence that shapes interpretation rather than construction is that \textbf{the small-sample regime predicts which components should help}. 156 clips cannot supply the data needed to learn a general spatio-temporal feature extractor from scratch, particularly once much of that extraction has been performed analytically before the network sees the data, whereas a component modelling the temporal arc of an expression exploits exactly the structure the corpus's onset, apex and offset annotation guarantees is present. This is a hypothesis the experiment tests, not a result it assumes. What the corpus literature does not supply is a stable basis for comparing methods: only Li et al.~\cite{liX2018} quantify how much of a reported difference is attributable to the protocol rather than the method, and then only for a single matched row of their own method.


% =============================================================================
\section{Motion Magnification: Eulerian Video Magnification}
\label{sec:evm-review}
% =============================================================================

Magnification is reviewed through the four corpus papers that apply it: Bai et al.~\cite{bai2021}, Li et al.~\cite{liX2018}, and Li, Huang and Zhao~\cite{liY2018,liY2021}. The original amplitude- and phase-based formulations lie outside the corpus and are described only as those papers attribute them. The mechanism itself is given in \autoref{sec:motion-magnification}.

The corpus design of \autoref{sec:why-spontaneous} answers the \textit{short-duration} half of the micro-expression definition but not the \textit{low-intensity} half. A movement spread across dozens of frames at 200~fps may still occupy only a few grey levels across a small patch of skin, and Yan et al.~\cite{yan2014} report discarding samples outright because the movement was too faint to place an onset. Li et al.~\cite{liX2018} state the remedy directly --- amplify the motion before any descriptor is computed, so that the difference between categories is enlarged rather than the classifier being made more sensitive. Being entirely upstream of the model is what makes magnification a clean ablation factor here.

Three parameters are free in any application: the temporal pass-band, which decides \textit{which} motions are amplified; the factor $\alpha$, which decides \textit{how much}; and the pyramid depth. On the first, Bai et al.~\cite{bai2021} report the rule given by the original authors --- the band should match the duration of the phenomenon, for micro-expressions 1/5 to 1/25\,s, equivalently \textbf{5--25 Hz} --- while noting that this full band produced ``noticeable noise'' in their setting, which led them to narrow it to 15--25 Hz. This is the only explicit guidance on band selection anywhere in the corpus.

\subsection{Evidence that magnification improves recognition}
\label{sec:evm-evidence}

The strongest case is Li et al.~\cite{liX2018}, who sweep the magnification factor rather than fixing it, across $\alpha \in \{1, 2, 4, 8, 12, 16, 20, 24, 30\}$ with $\alpha = 1$ denoting no magnification, on both SMIC and CASME II under leave-one-subject-out. Two findings emerge.

\textbf{Magnification helps consistently}, ``for all three kinds of features on all four testing datasets'', and on CASME II substantially (\autoref{tab:evm-effect}); their best full configuration reaches 67.21\,\% on CASME II, competitive with the contemporaneous state of the art.

\begin{table}[htbp]
	\centering
	\small
	\begin{tabular}{l c c c}
		\hline
		\textbf{Descriptor} & \textbf{No magnification} & \textbf{With magnification} & \textbf{$\Delta$} \\
		\hline
		LBP  & 55.87\,\% & 60.73\,\% & \textbf{+4.86} \\
		HOG  & 57.49\,\% & 63.97\,\% & \textbf{+6.48} \\
		HIGO & 57.09\,\% & 67.21\,\% & \textbf{+10.12} \\
		\hline
	\end{tabular}
	\caption[Effect of Eulerian motion magnification on CASME II]{Effect of magnification on CASME II under leave-one-subject-out, after Li et al.~\cite{liX2018} (Table 6). This is the corpus's one controlled magnification-on/off comparison.}
	\label{tab:evm-effect}
\end{table}

\textbf{The relationship between $\alpha$ and accuracy is non-monotonic.} Li et al.~\cite{liX2018} describe the curves as ``rainbow-shaped'', with best performance ``generally achieved when the motion is magnified in the range of $[8, 16]$'', and give a mechanistic account of both failure modes: too little magnification ``might not be enough to reveal the ME motion progress'', while too much ``degrades the performance because too many artifacts are induced''. Magnification is therefore a parameter with an interior optimum, not a monotone improvement to be maximised.

Bai et al.~\cite{bai2021} provide a second demonstration with a fundamentally different downstream model --- a pre-trained VGGFace2 encoder feeding a bidirectional LSTM --- reaching \textbf{75.76\,\% with phase-based magnification against a 60.60\,\% baseline} on SMIC, with the phase-based variant outperforming the amplitude-based one. The comparison is a distinct one and worth retaining, but the baseline is drawn from an earlier publication rather than from a magnification-off arm of their own, so it is not a matched control in the sense of \autoref{tab:evm-effect}. Because the pipeline shares nothing with hand-crafted descriptors, the two results together indicate the benefit is not an artefact of one feature family --- at least when the downstream representation is appearance-based.

Two further papers apply magnification in a different regime. Li, Huang and Zhao~\cite{liY2018,liY2021} magnify a single spotted apex frame at \textbf{$\alpha = 30$}, at the extreme end of Li et al.'s sweep, and augment by retaining the apex plus two frames either side. That $\alpha$ is asserted rather than derived, and the regimes differ enough --- a narrow window consumed as largely static appearance, against full magnified sequences described by spatio-temporal descriptors --- that the discrepancy need not be a contradiction. The safe reading is the weaker one: $\alpha$ is a parameter of the \textit{combination} of magnification and downstream representation, and cannot be transferred between pipelines without re-validation.

\subsection{Known limitations}
\label{sec:evm-limitations}

Four failure modes are documented, and the corpus is unusually candid about them.

\textbf{Magnification amplifies everything in band, not only expressions.} This is the most important and the most easily overlooked. Li et al.~\cite{liX2018} decline to use magnification for micro-expression \textit{spotting} at all, ``because it magnifies unwanted motions (e.g. head movements) at the same time''. A rigid head translation whose temporal frequency falls inside the pass-band is amplified exactly as a brow raise is. Prior alignment and cropping mitigate this; they do not eliminate it.

\textbf{Artefacts grow with $\alpha$}, reported by all four papers. \textbf{There is no principled procedure for choosing $\alpha$}: one value in the corpus is obtained by exhaustive sweep~\cite{liX2018}, and the apex-frame papers assert theirs. And \textbf{evaluation is confined to appearance-based downstream representations} --- LBP, HOG and HIGO in Li et al.~\cite{liX2018}, VGGFace2 features in Bai et al.~\cite{bai2021}, VGG-Face in both Li, Huang and Zhao papers. Every one is computed on magnified pixel intensities. \textbf{No study in the corpus measures what magnification contributes when the representation handed to the classifier is itself an explicit motion field.}

\subsection{Implications for this research}
\label{sec:evm-implications}
\label{sec:evm-motion-hypothesis}

The review fixes two settings and leaves one hypothesis. The temporal band is set to the \textbf{unnarrowed 5--25 Hz} duration rule, so Bai et al.'s narrowing is not followed and the pass-band admits the low-frequency range they rejected as noisy; this is a declared divergence. The factor is set to \textbf{$\alpha = 10$}, inside the reported $[8,16]$ optimum, and the apex papers' $\alpha = 30$ is explicitly not adopted because this thesis consumes full sequences. Amplitude-based magnification is used over a four-level Laplacian pyramid, and the band-pass is realised by an \textbf{FFT coefficient mask} rather than the Butterworth filter the reviewed description uses (\autoref{sec:temporal-filtering}) --- a further declared difference. Stage 1 runs twice, producing magnified and non-magnified tensor sets, so the switch selects a directory and leaves the network untouched.

Behind those settings lies a hypothesis the corpus does not test, stated because it frames the interpretation. Magnification amplifies the \textit{amplitude} of intensity variation in a band, and an appearance descriptor on magnified frames sees a directly stronger signal. An optical flow field is a different object: flow estimates \textit{displacement}, and amplifying the underlying intensity change by $\alpha$ does not simply scale the recovered displacement by $\alpha$. It changes the conditioning of the estimation problem, potentially making faint motion recoverable where it was not, while introducing displacement artefacts of the kind all four papers warn about; and any residual global scaling is further attenuated by the per-channel normalisation applied before training. The plausible outcome is that magnification in front of a flow-and-strain representation retains the part of its benefit consisting of \textit{making otherwise-unrecoverable motion recoverable} while losing the part consisting of \textit{increasing contrast for an appearance descriptor} --- so its measured contribution should be smaller on average than the +5 to +10 points of \autoref{tab:evm-effect}. Six matched pairs differing in nothing but the tensor directory are what allow that to be tested.

\textbf{A defect the review made detectable.} In this project's earliest evaluations every magnification-enabled configuration produced results \textit{bit-identical} to its disabled counterpart. Given that all four reviewed papers report clear, repeatable effects, an effect of exactly zero across six independent pairs is not a plausible null but a symptom: both arms were reading the same tensor directory, so the hypothesis had never been exercised. The defect was repaired before the reported run, and wiring the existing verification helper in as a start-up assertion is declared further work.

What the literature leaves open is threefold: every evaluation places magnification in front of an appearance representation, conflating two mechanisms that no experiment separates; $\alpha$ is transferred between studies without re-validation; and magnification is always evaluated as a two-cell addition to an otherwise fixed pipeline, so its interaction with other components is never measured. This thesis places it in front of an explicit motion representation, selects $\alpha$ from the full-sequence optimum, and measures it as one factor in a factorial ablation whose spread across architectures is itself the finding.


% =============================================================================
\section{Motion Representation: Dense Optical Flow}
\label{sec:flow-review}
% =============================================================================

Dense optical flow supplies the horizontal and vertical displacement channels. The relevant literature concerns three choices: the estimator, the frame pairs it processes, and whether flow serves as a recognition feature or an alignment tool. The estimation assumptions themselves are given in \autoref{sec:optical-flow}.

Xu et al.~\cite{xu2017} motivate facial-dynamics descriptors by the difficulty of interpreting micro-expressions through texture alone: flow describes apparent movement directly, whereas image intensity also records appearance and illumination. This does not guarantee independence from identity or recording conditions, since facial geometry, alignment and brightness changes still affect estimation. The benchmark evidence supports the choice --- the MEGC 2019 organisers report that their three highest-ranked submissions used optical flow~\cite{see2019}, OFF-ApexNet~\cite{liong2019a} uses horizontal and vertical flow as network inputs, STSTNet~\cite{liong2019b} adds strain, and Zhao et al.~\cite{zhaoS2021} retain a sequence of flow fields --- but this convergence is not a controlled demonstration that motion always outperforms appearance.

\textbf{The estimator is not settled by the corpus.} Zhao et al.~\cite{zhaoS2021} adopt Farneb\"ack through OpenCV, citing established availability; OFF-ApexNet~\cite{liong2019a} uses TV-L1, citing preservation of motion discontinuities; Xu et al.~\cite{xu2017} use the method of Sun, Roth and Black. STSTNet~\cite{liong2019b} specifies onset--apex flow without naming its estimator, so its similarity to OFF-ApexNet is insufficient to attribute TV-L1 to it. The advantages claimed for TV-L1 should not be read as a measured recognition gain over Farneb\"ack on CASME II: no controlled three-estimator comparison exists in the corpus. This thesis fixes Farneb\"ack as a preprocessing choice and leaves the comparison open.

\textbf{Onset--apex flow or a sequence is the consequential choice.} OFF-ApexNet~\cite{liong2019a} and STSTNet~\cite{liong2019b} represent each clip by flow from onset to apex, resized to $28\times28$ --- an economical summary of the displacement from neutral to peak, suitable for a small training set, which omits the intermediate trajectory and the relaxation after the apex. Zhao et al.~\cite{zhaoS2021} instead build an eleven-frame key-frame sequence and compute ten adjacent-pair fields, retaining changes in motion across the expression. The distinction is between one displacement summary and a sequence whose order a temporal model can exploit; it is not merely a difference in input size. This study follows the sequence approach but samples existing frames uniformly rather than constructing key frames, which preserves a temporal input for the transformer ablation. It does not measure the recognition cost of replacing a sequence with an onset--apex pair: that comparison is not an experimental arm.

Flow is also used for alignment rather than recognition. Xu et al.~\cite{xu2017} coarsely align by landmarks and then refine in-sequence using pixel-level movement, motivated specifically by micro-expressions, where a small registration error can be comparable to the displacement being recognised. This pipeline performs no flow-based fine alignment, so corpus preprocessing must not be treated as evidence that this error source has been eliminated. It remains the most promising unexploited improvement to preprocessing.

\textbf{Limitations of the representation.} Flow remains sensitive to brightness-constancy violations, alignment error and estimator settings; apex-based systems additionally depend on selecting an informative frame, and OFF-ApexNet~\cite{liong2019a} attributes weaker SMIC performance to imprecise apex selection. Using an onset--offset sequence avoids dependence on one selected apex but still relies on clip boundaries and sampling: repeated sampled frames yield near-zero-flow pairs, and large sampling intervals make the small-motion approximation less appropriate.

\subsection{Implications for this research}
\label{sec:flow-implications}

The central interpretive point is that optical flow is an \textit{analytical spatio-temporal feature extractor}: it estimates motion from paired frames before the network is trained. A learned spatial stem therefore acts on a representation that already makes displacement explicit, and with 156 clips, learning additional spatial filters may offer limited benefit or duplicate structure preprocessing has already exposed. This is a hypothesis about the stem's marginal value, not proof that learned spatial processing is unnecessary.

A temporal encoder addresses a different task: modelling relations among successive displacement fields, which pairwise flow estimation does not itself learn. Keeping those two roles distinct is necessary when attributing performance to the stem as against the transformer, and it is what the matched comparisons of this thesis are designed to separate. The design does not rank estimators, isolate alignment quality, or compare onset--apex against sequence representations; those remain limitations rather than gaps claimed as closed.


% =============================================================================
\section{Deformation Representation: The Optical Strain Tensor}
\label{sec:strain-review}
% =============================================================================

Optical strain supplies the third input channel. Its literature establishes a deformation descriptor derived from flow, several ways to use it, and limits relevant to interpreting the ablation. The tensor and its components are defined in \autoref{sec:optical-strain}; this section concerns the published evidence and one consequential divergence.

Flow describes displacement; strain describes its local spatial variation. A spatially uniform translation has zero displacement gradient, whereas neighbouring skin regions moving differently produce deformation. Shreve et al.~\cite{shreve2011} use this distinction to detect facial activity and describe robustness to moderate head translation, which does not imply immunity to arbitrary head movement or registration error. Their illumination- and make-up-robustness statements are attributed to earlier studies rather than measured anew, and should be read as reported motivation rather than as evidence about the present pipeline.

\subsection{The implemented magnitude and its divergence}
\label{sec:strain-divergence}

Liong et al.~\cite{liong2014a,liong2014b,liong2016} reduce the tensor to a scalar by the four-component Frobenius norm, which counts the two equal off-diagonal shear entries separately. \textbf{The implementation here counts shear once rather than twice}, giving the three-term magnitude of \autoref{sec:optical-strain}. Spatial derivatives come from central differences at one-pixel spacing, following Liong et al.~\cite{liong2014a}, but the scalar reduction differs from their convention.

STSTNet~\cite{liong2019b} cannot establish a precedent for the three-term form. Its printed magnitude equation repeats $\partial u$ in the mixed derivative where the symmetric tensor calls for $\partial v/\partial x$ alongside $\partial u/\partial y$, and it places the factor one-half \emph{outside} the squared derivative sum. Correcting the repeated derivative while retaining that placement yields the \textit{four}-term convention; moving the factor inside the square would instead yield the implemented form, but that is a further change to the printed expression. The intended implementation cannot be settled from the published equation alone, so this thesis declares a departure from the unambiguous four-term convention rather than claiming to reproduce STSTNet's magnitude.

The four-term form weights pure shear by $\sqrt{2}$ relative to the implemented one. Independent min--max normalisation removes a uniform scale factor but does not generally remove this difference, because normal and shear contributions vary across pixels and time. Neither convention was compared experimentally here, so their practical equivalence is not assumed.

\subsection{Recognition evidence and its protocol limits}
\label{sec:strain-evidence}

Strain has been used three ways: as a pooled descriptor (optical strain features, OSF), as a spatial weighting on LBP-TOP histograms (optical strain weighted features, OSW), and as a direct input channel. Liong et al.~\cite{liong2016} combine the first two, pooling spatially then temporally, and report the results of \autoref{tab:strain-performance}; the early work also applies Wiener and Gaussian filters explicitly to suppress noise~\cite{liong2014a}. STSTNet~\cite{liong2019b} supplies strain directly alongside onset--apex flow, which is the precedent for the three-channel input used here, independently of the magnitude ambiguity --- though this study retains all three channels across successive frame pairs rather than reducing each clip to one map.

\begin{table}[htbp]
	\centering
	\small
	\begin{tabular}{@{}lcc@{}}
		\hline
		\textbf{Method} & \textbf{SMIC accuracy (\%)} & \textbf{CASME II accuracy (\%)} \\
		\hline
		Baseline (texture) & 45.73 & 61.94 \\
		OSF                & 41.46 & 51.01 \\
		OSF + OSW          & 52.44 & \textbf{63.16} \\
		OFF + OFW (flow)   & 40.24 & \textbf{55.87} \\
		\hline
	\end{tabular}
	\caption[Strain-based descriptors against a matched flow control]{Strain-based descriptors against a matched flow-based control, after Liong et al.~\cite{liong2016} (Table 8). SMIC uses leave-one-subject-out; CASME II uses leave-one-\textit{video}-out, which is not subject-disjoint. The summary combines block settings, so equal entries elsewhere in the source are not matched-setting comparisons.}
	\label{tab:strain-performance}
\end{table}

The CASME II evidence depends on how strain is used and against what. OSF alone falls below the texture baseline; the combined descriptor improves on it by 0.41 to 1.22 points across block settings~\cite{liong2016}. Crucially, Liong et al. also test flow features and weights within the same framework, reaching 55.87\,\% against 63.16\,\% for the strain pair --- a \textbf{matched flow-versus-strain control}, so it would be incorrect to claim the literature never isolates strain against flow. What it concerns is hand-crafted descriptors under a protocol that is not subject-disjoint, and these accuracy gains are not comparable with pooled macro F1. It does not establish the gain from adding a strain channel to a learned sequence model.

\subsection{Implications for this research}
\label{sec:strain-implications}

Strain extends the analytical-feature argument of \autoref{sec:flow-implications}: the input contains a \textit{hand-specified differential operator}, exposing local deformation without learning that operator from 156 clips. A spatial stem may still learn useful combinations of those features, but it operates after both motion estimation and deformation extraction. This motivates testing its incremental contribution; it does not establish the cause of any observed architectural effect.

Three limits bound what this thesis can say about strain. Spatial differentiation can amplify flow-estimation noise, and \textbf{no dedicated strain post-filter is applied}, departing from the explicit filtering of Liong et al.~\cite{liong2014a}; Farneb\"ack's internal smoothing is not a separately evaluated strain filter. The shear weighting is the untested choice of \autoref{sec:strain-divergence}. Above all, \textbf{strain is present in all twelve evaluated configurations}, so its marginal contribution is not measured at all: a strain-on/strain-off comparison is the natural next ablation, and neither published gains from a different descriptor nor the performance of the complete system substitutes for it. A deterministic transformation of $u$ and $v$ adds no independent observation; its potential value lies in making deformation directly accessible to the classifier, and whether that helps must be measured rather than inferred from its construction.


% =============================================================================
\section{Temporal Normalisation: Fixed-Length Sampling}
\label{sec:tim-review}
% =============================================================================

Flow and strain fix \textit{what} is computed from each frame pair; temporal normalisation fixes \textit{which} pairs, and how many. The Temporal Interpolation Model (TIM) is credited in this corpus to Zhou et al. by Li et al.~\cite{liX2013,liX2018} and to Pfister et al. by Lu et al.~\cite{lu2015}, and is described here only through the reviewed papers that implement it.

Normalisation is necessary because any batched network requires a fixed input shape and the corpus supplies no constant, spanning a four-fold range of sequence length; and because duration is itself a nuisance variable --- Lu et al.~\cite{lu2015} normalise to remove ``temporal fluctuations of micro-expressions, which may introduce noise irrelevant to classification''. They also give the corpus's most compact description of TIM: ``a manifold-based interpolation method, which builds a low-dimensional embedding of an image sequence, and then interpolates a curve in the low-dimensional space. The interpolated frames are mapped back to the original high-dimensional space.'' The property that matters downstream is that TIM \textbf{synthesises frames that were never recorded}, having been proposed ``to generate sufficient frames from an image sequence'' --- an upsampler rather than only a length-standardiser.

\textbf{Ten frames is a convention, not a finding, and the studies that test it disagree.} TIM appears throughout the corpus as an unremarked preprocessing step, almost always at ten frames~\cite{liX2013,liX2018,lu2015}. Li et al.~\cite{liX2018} sweep interpolation length from 10 to 80 on the SMIC variants with the descriptor held fixed, and report that ``the best performance for all three datasets is achieved with interpolation to 10 frames'', because ``in longer sequences, changes along the time-dimension are diluted''. Ben et al.~\cite{ben2021} reach the opposite conclusion: holding alignment, descriptor and classifier fixed, they interpolate to 10 through 110 frames by both Newton interpolation and TIM on MMEW, finding Newton peaking at 33.3\,\% at 60 frames and TIM at \textbf{38.9\,\%} at 30 or 60, and adopt TIM at 60. Accuracy is non-monotonic in length in both. Xu et al.~\cite{xu2017} add the only multi-corpus sweep --- 10, 15 and 20 frames across CASME, CASME II and both SMIC subsets --- and depart from convention by using \textbf{linear} interpolation ``instead of Temporal Interpolation Model because in a high frame-rate scenario, linear interpolation suffices to characterize the motion pattern with less error''.

Two caveats limit what transfers. Ben et al.'s experiment uses MMEW with LBP-TOP, not CASME II and not a learned model, and runs largely in the \textit{upsampling} regime, so an optimum measured while synthesising frames need not hold when discarding them. And \textbf{frame synthesis is in tension with motion analysis}: where the downstream representation is appearance-based, synthesised frames are benign, but where it is motion, synthesised frames yield synthesised displacements. Xu et al.~\cite{xu2017} are the only reviewed paper to raise this, and they assert it rather than measure it. No reviewed study compares the two interpolators against a motion representation.

\subsection{Implications for this research}
\label{sec:tim-implications}

\textbf{A declaration about what this pipeline implements.} The module performing this step is named for the Temporal Interpolation Model, but what it computes is \textbf{uniform index sampling} of existing frames (\autoref{sec:three-channel-tensor}). No embedding is built, no curve is interpolated, no frame is synthesised. That is a materially different operation from TIM, and it is named accurately here rather than in the pipeline's own terms. It is defensible: because the downstream representation is flow between adjacent sampled frames, interpolation would manufacture displacement fields for motion that never occurred, whereas subsampling discards temporal resolution without fabricating it. What is given up is TIM's ability to \textit{lengthen} short clips.

The length is fixed at \textbf{33 frames sampled, giving 32 flow pairs}, bounded by the annotated onset and offset so the window spans the expression rather than surrounding neutral frames. That is close to the lower of Ben et al.'s two reported optima, but the correspondence should not be overstated: their 30 frames are \textit{interpolated} frames under a hand-crafted descriptor on a different corpus, which is not direct validation of 32 sampled flow fields under a learned model here. The length is not swept, and a sweep is declared as further work.

Three quantitative consequences follow, and they constrain the magnification conclusion rather than merely describing the sampler.

\textbf{Effective temporal resolution falls by a clip-dependent factor.} Sampling 33 frames from a clip of $N$ frames recorded at 200~fps gives an effective rate of $6600/N$ fps: \textbf{100 fps} at the median $N = 66$, which is SMIC's recording rate and half what CASME II was built to provide, and roughly \textbf{52 fps} for the longest clip, below the original CASME database's 60~fps~\cite{yan2014}. Eighty-one of the 156 clips are downsampled by a factor of two or more, which gives back part of what the 200~fps capture rate bought. No reviewed paper reports this cost.

\textbf{Magnification is applied after subsampling, which shifts the realised pass-band.} Magnification runs on the 33 sampled frames at the nominal 200~fps, but those frames span the clip's full duration, so the true inter-frame interval is $N/(200 \times 32)$ seconds rather than $1/200$. At the median, a band specified as 5--25 Hz corresponds to roughly 2.5--12.5 Hz physically, and the discrepancy grows with clip length. The realised band is therefore neither the 5--25 Hz that \autoref{sec:evm-implications} derives nor constant across clips. This is a declared divergence; magnifying before subsampling is declared as further work.

\textbf{Three clips are shorter than the sampling window.} Three of the 156 contain 31 frames against a target of 33, so sampling repeats indices and yields near-zero flow at those positions.

The unresolved question this leaves is what temporal normalisation costs a representation made of motion rather than appearance. Three studies sweep interpolation length and give three answers, only one of them touching CASME II and none using a learned temporal model. This thesis takes the position that follows from its own representation and quantifies what that costs in effective frame rate, but does not sweep the length, compare interpolators, or separate the effect of sequence length from that of the sampling regime.


% =============================================================================
\section{The Learned Spatial Backbone: Shallow Convolutional Networks}
\label{sec:cnn-review}
% =============================================================================

By 2019 the field had converged on optical flow as the input; the open question was what should consume it. A learned backbone promises what hand-crafted descriptors cannot~\cite{liong2019b}, but runs into the small-data constraint of \autoref{sec:corpus-threats}.

\textbf{The term ``3D CNN'' covers three materially different architectures in this literature.} STSTNet~\cite{liong2019b}, the canonical shallow model, calls itself a ``Shallow Triple Stream Three-dimensional CNN'', but its input is a single $28 \times 28 \times 3$ tensor whose third dimension holds horizontal flow, vertical flow and strain --- \textbf{not time}. Its three parallel single-layer streams, with 3, 5 and 8 kernels, supplement ``the small scale input data \ldots to avoid the problem of underfitting''; the clip was collapsed to one frame pair upstream. STRCN~\cite{xia2020a} by contrast extends convolutional connectivity genuinely ``across temporal domain'' with recurrent convolutional layers. Zhao et al.~\cite{zhaoS2021} apply a C3D-style network to a ten-frame flow sequence, so the temporal axis is present but short. Results reported across the corpus under the same label are therefore not like with like.

\textbf{The central controlled finding is that shrinking both model and input helps.} Xia et al.~\cite{xia2020b} treat depth and input resolution as controlled variables, building four models from a single convolutional layer to one convolutional plus three recurrent convolutional layers, and evaluating each across input resolutions from 20 to 300 pixels on optical flow. Deeper is worse: models 3 and 4 ``have worse recognition performances compared to shallower \ldots models'', even though all four are far shallower than ResNet-101. Higher resolution is worse with a threshold: performance ``will be degraded with larger input resolutions'' and for the deeper models ``begins to decrease dramatically when the resolution is bigger than \textbf{100 $\times$ 100}'' --- though shallow models are ``basically robust to the change of input resolution''. That qualification matters for this thesis, and so does the setting: Xia et al.'s study and STSTNet are both developed for the MEGC composite regime, which supplies substantially more training data per fold than a single corpus (\autoref{sec:megc-metrics}), so conclusions about affordable capacity understate the constraint here. Xia et al.~\cite{xia2020b} also take a third route to capacity control that this thesis follows --- developing ``three parameter-free modules \ldots integrate with RCN without increasing any learnable parameters'', so representational power grows while overfitting risk does not.

\textbf{No study isolates the backbone itself.} The corpus is not short of matched comparisons: it sets descriptor-only methods against learned ones on the same folds, and Xia et al.~\cite{xia2020b} vary individual architectural modules against a common baseline with every other parameter held constant. But those comparisons all \textit{add} to a backbone that stays in place. None removes the learned spatial stem entirely and passes the motion representation straight to a classifier or temporal model, so what the stem contributes against what it costs is nowhere measured.

\subsection{Implications for this research}
\label{sec:cnn-implications}

The component ablated here is a \textbf{three-stream shallow convolutional backbone}, one unshared two-layer stream per input channel, carrying \textbf{14,544 parameters} --- shallow by the corpus's standards, though roughly nine times STSTNet's. Where STSTNet's streams differ in filter count over the same input, these differ in input channel. Removing it leaves not an empty slot but spatial pooling followed by a learned linear projection, so the no-backbone route still contains learned input processing.

\textbf{A declaration about what this component is.} Every kernel in the backbone is spatial only: no filter spans two time steps, so despite the \texttt{Conv3d} implementation it is a \textbf{purely spatial} feature extractor (\autoref{sec:conv3d}). The ablation therefore does not test spatio-temporal convolution; it tests a learned spatial stem, and whatever it measures is evidence about \textit{this} stem --- spatial-only, at 224~$\times$~224, on flow and strain input --- and \textbf{not} evidence that convolution over time is unhelpful for micro-expression recognition, which this study never evaluated. In the field's own terms the choice is defensible: STSTNet, one of the strongest shallow baselines on the CASME II subset, likewise performs no temporal convolution.

\textbf{Input resolution is the clearest departure from Xia et al.'s guidance.} The backbone is fed at \textbf{224 $\times$ 224}, more than twice their threshold. This is not a departure from the corpus as a whole --- Bai et al.~\cite{bai2021} resize to 224 $\times$ 224 without alignment or cropping, as here; Li, Huang and Zhao~\cite{liY2021} resize faces to the same; and Xia et al.'s own earlier work~\cite{xia2020a} runs one stream at 300 $\times$ 245 --- but it departs from the specific threshold above, which is a composite-database result that bites hardest on models deeper than this one. Since this stem is shallow, the reviewed evidence predicts a smaller resolution effect than the deep-model figures suggest. The choice is also the expensive one: a 224 $\times$ 224 map carries sixty-four times the spatial elements of STSTNet's input, and the backbone convolves it at every one of the 32 time steps. Any effect measured for the stem is therefore \textbf{a resolution result as much as an architecture result}, and the reviewed evidence does not determine which mechanism dominates; re-running the backbone arm at a lower resolution is the first item of declared further work because it is cheap and would settle the question, not because the corpus implies a particular outcome.


% =============================================================================
\section{Parameter-Free Attention: SimAM}
\label{sec:simam-review}
% =============================================================================

SimAM re-weights convolutional feature maps without adding learned parameters. Yang et al.~\cite{yang2021} supply the mechanism and its image-recognition evidence; Xia et al.~\cite{xia2020b} motivate parameter-free modules in composite-database micro-expression recognition, supporting the design principle rather than validating SimAM on this task. The energy formulation and gate are given in \autoref{sec:simam}.

Yang et al.~\cite{yang2021} motivate attention through spatial suppression: neurons distinguishable from their surroundings are assigned greater importance, expressed as the linear separability of one activation from the others in its channel. This is a prior about activation statistics, not a detector of action units or minority classes --- a distinctive response may encode useful deformation, noise or an artefact --- so its suitability for flow-derived features requires evaluation rather than following from the neuroscience motivation. Two implementation details are owned here rather than by \autoref{ch:background}: the shared variance printed with the paper's equation~(5) divides by $M$ whereas its Figure~3 reference code divides by $M-1$, and this thesis follows the code; and the module's ``three-dimensional weights'' describe $C \times H \times W$ with $M = HW$ per frame, not a temporal volume, so extending the context across frames is an additional design decision rather than part of the original formulation.

Yang et al. evaluate SimAM with several network families on CIFAR and ImageNet, reporting competitive performance against learned attention modules, and their timing results show that unchanged parameter totals do not imply unchanged inference speed. For CIFAR they sweep $\lambda$ from $10^{-1}$ to $10^{-6}$ and select $10^{-4}$; they repeat the selection for ImageNet and choose $0.1$. Thus $10^{-4}$ is a dataset-selected operating point, not a universal setting justified by the formula.

\textbf{The limits of transfer are specific.} The source experiments concern natural images, different training-set sizes and learned feature distributions. Here the input is flow and strain, so distinctiveness need not bear the same relationship to useful evidence. The image experiments also do not decide whether video statistics should be computed per frame or over a clip. And the absence of learned attention weights does not establish that runtime, memory or overfitting concerns disappear. These are empirical limits on transferring the result, not reasons to assume success or failure.

\subsection{Implications for this research}
\label{sec:simam-implications}

SimAM is applied separately to each convolutional stream before concatenation, introducing no trainable weights or biases anywhere. Two choices diverge from the source. Mean and variance span $(T,H,W)$ within each channel with divisor $THW - 1$, rather than one $(H,W)$ plane, so an activation is judged distinctive relative to its whole clip rather than its own frame; this could emphasise a brief contraction, but does not guarantee low weights at onset or high weights at the apex, since the score depends on squared deviation and not on expression labels. And $\lambda = 10^{-4}$ is retained without re-tuning on CASME II. Neither context choice is compared experimentally.

SimAM is also enabled only alongside the convolutional stem in this model, which describes this implementation rather than a mathematical prohibition. Four of the sixteen candidate configurations are therefore excluded, leaving twelve evaluated and \textbf{four SimAM matched pairs}, and the measured effect is conditional on the stem being present. Retention should be judged against both recognition performance and computational cost: parameter count alone cannot settle it, because the module is cheap rather than free.

\subsection{Relevance and scope}
\label{sec:simam-gap}

The study evaluates SimAM in a motion-input, small-corpus setting through four matched comparisons under complete leave-one-subject-out, testing the implemented module as a whole rather than establishing the benefit of each design choice. Three follow-ups remain: whether the distinctiveness prior suits flow-derived features, re-tuning $\lambda$, and comparing per-frame with whole-clip statistics. Each requires new evaluations despite the module introducing no learned weights.


% =============================================================================
\section{The Temporal Encoder: Transformers}
\label{sec:transformer-review}
% =============================================================================

The temporal encoder is motivated by Zhang et al.'s Short and Long Range Relation Based Spatio-Temporal Transformer (SLSTT)~\cite{zhang2022} and its Vision Transformer foundation~\cite{dosovitskiy2021}. The component evaluated here differs substantially from published SLSTT, particularly in the axis over which self-attention operates, and distinguishing the two is the purpose of this section.

A micro-expression evolves through onset, peak and relaxation, so retaining a motion sequence permits a model to relate separated stages rather than compressing them immediately into one descriptor; Zhang et al.~\cite{zhang2022} motivate their design through the need to capture both local and global patterns. This thesis compares a temporal transformer against mean pooling over the same sequence: mean pooling discards order, whereas positional information and self-attention allow relations among frame-pair features to be modelled. That motivates the comparison without establishing in advance which generalises better. ViT itself~\cite{dosovitskiy2021} splits an image into patches, projects them into tokens, adds learned positional embeddings and relates them by attention; Dosovitskiy et al. also show that without the locality and translation-equivariance biases of convolution, image transformers can perform poorly when trained with insufficient data, while large-scale pre-training improves them. That is a caution for a 156-clip study, not a universal requirement that every transformer be pre-trained.

\textbf{Published SLSTT combines three stages.} First, \textbf{onset-referenced optical flow} describes sampled frames relative to the onset rather than to their immediate predecessors, which Zhang et al. argue gives a more structured trajectory towards and away from the apex; their experiments use it throughout, without a matched short-term-flow ablation establishing its independent gain. Second, a \textbf{spatial ViT encoder} processes each flow field as patches, using ImageNet-pre-trained ViT-B/16 with twelve layers at $384 \times 384$. Third, an \textbf{LSTM temporal aggregator} combines the per-frame features, with a mean aggregator supplying the comparison variant. Self-attention is spatial in this architecture; recurrence handles time.

\subsection{The aggregation comparison, and what it can support}
\label{sec:slstt-results}

The comparison relevant here is internal to SLSTT and holds the spatial transformer fixed in both arms: CASME II UF1 rises from \textbf{0.844} with the mean aggregator to \textbf{0.901} with the LSTM, a difference of 0.057~\cite{zhang2022}. That supports evaluating aggregation choices within that framework. It is \textbf{not} an estimate of what an LSTM would add to this thesis, where mean pooling follows \emph{temporal} self-attention rather than independent spatial encoding, and it isolates an aggregator rather than the presence of a transformer, since both variants contain one.

Protocol labels need care. Zhang et al. describe three-class leave-one-subject-out evaluations on the databases separately and together under a Composite Database Evaluation heading, so that heading alone does not establish that every CASME II entry uses pooled cross-database training; they also report a distinct Sole Database Evaluation using original labels, including five CASME II classes~\cite{zhang2022}. Neither setting should be equated automatically with this thesis's 156-clip three-class pool. Published SLSTT also retains a pre-trained spatial encoder, so its results do not isolate the benefit of a transformer trained from scratch on this study's motion representation.

\subsection{Implementation and six declared divergences}
\label{sec:transformer-implications}

The evaluated component receives 32 tokens of width 96, from either the convolutional stem or spatial pooling followed by a learned linear projection. The encoder uses four pre-norm layers, eight heads, feed-forward width 256 and dropout 0.1; fixed sinusoidal positions precede attention and mean pooling follows it. With the encoder disabled the input tokens are averaged directly.

\begin{table}[htbp]
	\centering
	\small
	\begin{tabular}{@{}p{0.23\linewidth} p{0.34\linewidth} p{0.35\linewidth}@{}}
		\hline
		\textbf{Property} & \textbf{Published SLSTT~\cite{zhang2022}} & \textbf{This thesis} \\
		\hline
		Attention axis   & Spatial patches within a frame              & Temporal tokens within a clip \\
		Aggregation      & LSTM over spatially encoded frames          & Mean after temporal self-attention \\
		Initialisation   & ImageNet-pre-trained ViT-B/16               & Trained from scratch \\
		Capacity         & Twelve layers, width 768; $384\times384$ in & Four layers, width 96; 32 tokens \\
		Positions        & Learned spatial embeddings                  & Fixed sinusoidal temporal encoding \\
		Input flow       & Relative to onset                           & Relative to the preceding sampled frame \\
		\hline
	\end{tabular}
	\caption[Divergences from published SLSTT]{Six material divergences from published SLSTT. The component evaluated here is a small temporal transformer encoder, not a reproduction of the published architecture.}
	\label{tab:slstt-divergences}
\end{table}

Six matched pairs evaluate the encoder with other component choices fixed within each pair. They test the whole encoder --- positional encoding, feed-forward layers and normalisation included --- against direct mean pooling, and do not isolate self-attention from those accompanying operations or compare it against another learned temporal model. Three follow-ups follow from \autoref{tab:slstt-divergences}: testing an LSTM aggregator, comparing onset-referenced with consecutive-frame flow, and replacing fixed positions with learned embeddings. All three remain untested, and changing the input representation requires re-training before any benefit can be claimed. What this study can establish is a matched comparison of the implemented encoder against an order-blind baseline within this pipeline; it neither validates published SLSTT nor establishes a general ranking of attention, recurrence and convolution.


% =============================================================================
\section{Training Under Severe Class Imbalance}
\label{sec:imbalance-review}
% =============================================================================

The 99 : 32 : 25 distribution of \autoref{sec:label-taxonomy} is a property of what can be elicited in a laboratory, not a sampling accident, so it cannot be corrected by collecting more data from the same paradigm. The corpus corrects for it at three points --- the data, the loss and the metric --- and this thesis uses all three. The mechanisms are given in \autoref{sec:focal-loss} and \autoref{sec:balanced-sampling}.

\textbf{At the loss.} Zhao et al.~\cite{zhaoS2021} bring focal loss into micro-expression recognition ``to alleviate the inefficient model training caused by the class imbalance in the MEs Datasets'', adopting the formulation developed for one-stage object detection. The division of labour inside it is what matters: $\gamma$ re-weights by \textit{difficulty} and is treated as a constant, ``set as 2 in practice'', while $\alpha$ re-weights by \textit{class frequency} and is ``treated as a hyper-parameter to set by cross validation''. Conflating the two is the error discussed below. Zhao et al. compare focal loss against cross-entropy across a sweep of $\alpha$, but hold $\gamma$ fixed, so the focusing exponent remains unvalidated on this task.

\textbf{At the data.} Xia et al.~\cite{xia2020a} pair ``temporal data augmentation strategies as well as a balanced loss \ldots jointly'' to overcome ``limited and imbalanced training samples''. Their own ablation removes each correction in turn and reports the balanced loss improving performance only ``slightly'' --- at most about two points, and in one cell costing accuracy. This is the corpus's only controlled look at composing two corrections, and no reviewed paper reports the pairing failing outright.

\textbf{At the metric.} \autoref{sec:megc-metrics} defines the pooled macro-averaged measure. Liong et al.~\cite{liong2016} add the argument from the subject side, reporting F1, precision and recall because these ``provide a more meaningful perspective than accuracy rates when the datasets used are naturally imbalanced since each subject has a different number of video samples.'' Both arguments compound here, where classes are skewed \textbf{and} per-subject clip counts range from 1 to 33.

One further gap: \textbf{label smoothing is named in the corpus but never examined on this task.} Dosovitskiy et al.~\cite{dosovitskiy2021} list it among the regularisers they tune, but no reviewed micro-expression paper uses it or reports what it does under class skew.

\subsection{Implications for this research}
\label{sec:imbalance-double-correction}

Focal loss is adopted with \textbf{$\gamma = 2.0$}, matching Zhao et al.'s reported setting. The class-frequency term is implemented as an optional per-class $\alpha$ vector but \textbf{switched off at run time}, because frequency correction is assigned to a weighted sampler active in every configuration. Label smoothing of 0.05 is applied inside the focal loss, declared as a choice this corpus supports neither way. The primary metric is pooled macro F1, with accuracy quoted only alongside the always-Negative floor.

\textbf{What this project contributes is a negative result about composing corrections.} The configuration is resolved at run time by a single rule: class weighting in the loss is enabled only when the balanced sampler is \textit{not} active. That rule exists because of two failures in opposite directions, both recoverable from this project's own run history. In an earlier code era the loss carried inverse-frequency class weights and no sampler existed; under that single correction the proposed configuration abandoned the \textbf{Positive} class outright, a failure initially misdiagnosed as over-capacity overfitting. A later era added the sampler while leaving the loss weights unconditional, so both corrections ran at once; the collapse then inverted, and configurations abandoned the \textbf{Negative} majority class instead, one predicting a single class for every clip. With the guard in place, no configuration in the reported study abandons any class.

This is worth reporting because no reviewed paper warns that combining corrections can invert the intended effect rather than merely blunt it. It must be stated carefully, however: the two collapses occurred in code eras differing from the reported run in more than the correction rule, and the final sweep contains no matched arm with the guard disabled. \textbf{The guard's causal benefit is therefore not isolated}, and the evidence does not establish that correcting at exactly one point is universally best --- only that correcting twice, in this pipeline, produced a class collapse that correcting once did not.

A second declaration follows. Because every configuration shares the same sampler, loss and $\gamma$, this thesis measures \textbf{no} effect attributable to the imbalance treatment: it is a constant, not a factor. The claim supported is that the four architectural components were compared under a training regime that does not abandon minority classes, not that this regime is optimal.

\subsection{Relevance and scope}
\label{sec:imbalance-gap}

The corrections themselves are well established. What the corpus does not settle is how two corrections compose, what $\gamma$ should be on this task, and what over-correction looks like when it fails. This thesis contributes to the first and third by documenting two opposite class-collapse failures and the run-time rule that resolved them. On the second it offers no improvement, since $\gamma = 2.0$ and label smoothing of 0.05 are held constant across all twelve configurations; sweeping either is declared as further work.


% =============================================================================
\section{Synthesis and Research Gap}
\label{sec:synthesis-and-gap}
% =============================================================================

This section consolidates the preceding relevance statements into one research gap, and collects in one place every declared divergence between the system built here and the literature it draws on.

Read as a whole, the reviewed literature settles a good deal. The corpora are small, demographically narrow, laboratory-bound and irreducibly imbalanced by what can be elicited, and CASME II remains the primary single-corpus benchmark under its native leave-one-subject-out protocol. Motion rather than appearance is the appropriate input; a strain field adds a further cue, though on CASME II only weakly and under a protocol that is not subject-disjoint; and both the magnification factor and the sequence length have interior optima located only by sweep. At this data size, capability bought without parameters outperforms capability bought with them, and self-attention is the strongest-performing mechanism reported. Accuracy alone is inadequate under skew and must give way to macro-averaged measures accumulated across folds before averaging.

\subsection{Four limitations that recur across every topic}
\label{sec:four-limitations}

The per-section statements above are not independent observations. Four structural problems generate most of them.

\textbf{(a) Components are isolated one at a time, inside one proposed design, and never jointly.} Matched-pair measurement is not absent from the corpus: \autoref{tab:evm-effect} isolates magnification, \autoref{tab:strain-performance} isolates strain against a flow control, Xia et al.~\cite{xia2020b} compare a recurrent convolutional network against three variants each adding one module with every other parameter held constant, and SLSTT's Mean-versus-LSTM comparison isolates its aggregator (\autoref{sec:slstt-results}). What no reviewed study does is vary \textit{several} components together, so that their interactions can be read. Every ablation is internal to one proposed design and moves one factor from that design's own defaults; across papers, systems are compared whole, so a difference in the reported number cannot be attributed to any one stage. There is therefore no basis for deciding which learned stages are load-bearing, nor whether one component's value depends on another's presence.

\textbf{(b) The evaluation protocol is unstable and under-reported.} Sample counts, label-set size and the choice among leave-one-subject-out, leave-one-video-out and leave-one-sample-out vary inconsistently, and the field's most-quoted baseline is described three ways by its own sources (\autoref{sec:original-baseline}). Given this corpus's fold composition these are not bookkeeping differences: they change what the numbers mean.

\textbf{(c) The strongest results depend on regimes a single-corpus study cannot access.} The best CASME II figures come from composite-database training (\autoref{sec:megc-metrics}) and, for SLSTT, ImageNet pre-training --- neither a property of the architecture being credited. Almost no evidence exists about what these components do when trained from scratch on one small corpus, which is the situation most practitioners are in.

\textbf{(d) Preprocessing parameters have interior optima and are transferred without re-validation.} The magnification factor, the sequence length and the input resolution all rise then fall, and none has a principled selection rule. Values are carried between pipelines whose downstream models differ entirely: $\alpha = 30$ moves from a single-frame appearance encoder to other settings, ten-frame interpolation propagates from one baseline paper to the whole field, and $\lambda = 10^{-4}$ transfers from CIFAR though its own authors re-tuned it for ImageNet.

\subsection{The research gap}
\label{sec:research-gap}

Taking (a) to (d) together:

\textbf{The reviewed literature has established a standard pipeline for micro-expression recognition --- magnify, compute motion, normalise duration, extract spatially, model temporally --- and contains matched ablations of individual stages, but no study varies several of those stages jointly so that their relative contributions and interactions can be read. Because components are varied one at a time inside a single proposed system, under protocols that differ in unstated ways, and predominantly in training regimes richer than a single small corpus provides, it is not currently possible to say from this literature which parts of that pipeline earn their place, which are inherited convention, and which depend on the presence of another.}

This matters practically as well as scientifically. A practitioner assembling this pipeline today has no published basis for allocating effort: for judging whether the most computationally expensive stage is the most valuable, whether a component with a negligible average effect may still matter in a specific failure, or whether a ranking read off an incomplete evaluation of a corpus this size is stable at all.

\subsection{How this thesis addresses the gap, and what it does not}
\label{sec:addressing-gap}

\textbf{What it does.} This thesis treats the pipeline as a factorial experiment rather than a proposal: four components varied independently across every architecturally valid cell of a 2\textsuperscript{4} matrix, twelve configurations, each evaluated under \textbf{complete 25-fold leave-one-subject-out on CASME II} with every non-varied factor held identical. That supplies the joint measurement (a) identifies as missing, reporting the protocol in full addresses (b) for this study at least, and training from scratch is the regime (c) identifies as unevidenced. Each component's marginal contribution is read from pairs of runs differing in one factor alone, and the distribution of those four effects is the substantive answer to the gap. Two further contributions follow from the review rather than the experiment: a \textbf{repaired measurement}, where an early data-routing defect left the magnification switch inert and the review's evidence of a clear, repeatable magnification effect is what made bit-identical results detectable as a defect rather than a null (\autoref{sec:evm-implications}); and \textbf{two opposite class-collapse failures} that no reviewed paper reports (\autoref{sec:imbalance-double-correction}).

\textbf{What it does not do.} Three of the four limitations are only partially addressed. On (b), this study reports its own protocol completely but cannot repair the field's inconsistency, and its own N = 156 differs from MEGC's 145. On (c), it establishes what these components do from scratch on one corpus, which is a different question from --- not an answer to --- how they behave under composite training. On (d), the thesis inherits parameter values rather than sweeping them: $\alpha = 10$, T = 32, $\lambda = 10^{-4}$ and $\gamma = 2.0$ are fixed throughout, so the ablation measures architecture at one point in preprocessing space rather than across it.

\subsection{A consolidated ledger of declared divergences}
\label{sec:divergence-ledger}

Reviewing the literature closely enough to write this chapter surfaced twelve points at which the implemented system departs from the work it draws on. Several were discovered only through the review. They are collected in \autoref{tab:divergence-ledger} so that no result rests on an unstated assumption.

\begin{table}[htbp]
	\centering
	\footnotesize
	\begin{tabular}{@{}p{0.85cm} p{6.0cm} p{6.6cm}@{}}
		\hline
		\textbf{\S} & \textbf{Divergence} & \textbf{Consequence} \\
		\hline
		\ref{sec:label-taxonomy} & Working set is N = 156, including sadness and fear in Negative; MEGC's CASME II subset is N = 145 & Comparisons against challenge figures are not exactly like-for-like \\
		\ref{sec:casme2-construction} & The corpus's own registered release is not used; the pipeline reads the original recorded frames, unregistered and uncropped & No landmarking or alignment of any kind precedes the network; tolerance to pose and scale must be learned from 156 clips \\
		\ref{sec:evm-implications} & Temporal band is the unnarrowed 5--25 Hz duration rule, where Bai et al. narrowed it to 15--25 Hz as noisy & The pass-band admits the low-frequency range they rejected \\
		\ref{sec:evm-implications} & Band-pass is realised by an FFT coefficient mask, not the Butterworth filter of the reviewed description & A different filter shape, with different transition behaviour at the band edges \\
		\ref{sec:strain-divergence} & Strain magnitude counts shear once, not twice as in the published Frobenius norm & Shear weighted lower by $\sqrt{2}$ relative to normal strain; untested \\
		\ref{sec:strain-implications} & No dedicated strain filtering, where the literature applies Wiener and Gaussian filters & Spatial differentiation can amplify flow-estimation noise, and nothing downstream suppresses it \\
		\ref{sec:tim-implications} & Temporal normalisation is uniform index sampling, \textbf{not} manifold-based TIM; no frames are synthesised & Avoids fabricating motion; forfeits the ability to lengthen short clips \\
		\ref{sec:tim-implications} & Magnification is applied \textbf{after} subsampling, at a nominal 200 fps & Realised pass-band is clip-dependent and below the intended 5--25 Hz \\
		\ref{sec:cnn-implications} & Backbone kernels are spatial only --- \textbf{no temporal mixing} & The ablation tests a learned \textit{spatial} stem, not spatio-temporal convolution \\
		\ref{sec:cnn-implications} & Backbone input is 224 $\times$ 224, against Xia et al.'s guidance of $\leq$ 100 $\times$ 100 --- though Bai et al. and Li, Huang and Zhao both use 224 $\times$ 224 & Whatever the stem is measured to contribute may be a resolution effect as much as an architecture effect \\
		\ref{sec:simam-implications} & SimAM statistics pooled over the whole clip rather than per frame as in the source; $\lambda$ not re-tuned & Neurons judged against the clip, not the frame; defensible but untested \\
		\ref{sec:transformer-implications} & Six divergences from published SLSTT (temporal not spatial attention, mean not LSTM aggregation, from scratch, smaller, sinusoidal encoding, short-term flow) & The temporal encoder is \textbf{not} a reproduction of SLSTT \\
		\ref{sec:imbalance-double-correction} & Label smoothing of 0.05, named in the corpus only by Dosovitskiy et al. and never examined on this task & A choice the micro-expression corpus supports neither way, held constant and unmeasured \\
		\hline
	\end{tabular}
	\caption[Every declared departure from the reviewed literature]{Every declared departure from the reviewed literature, with its consequence. The first column resolves to the subsection that owns each divergence.}
	\label{tab:divergence-ledger}
\end{table}

Three of these --- the strain formula, the magnification ordering and the input resolution --- were identified by the review itself rather than being known design decisions, and each is declared as proposed further work. None invalidates the ablation, since every configuration shares the same choices. What they constrain is the \textit{wording} of the conclusions: what is measured is the contribution of a spatial stem at 224 $\times$ 224 and of temporal self-attention in this pipeline, not a verdict on convolution over time or on SLSTT. The limits of what can be claimed are therefore fixed before any number is reported --- a small, single-seed, demographically narrow corpus measured against richer published regimes, with these divergences standing throughout.
 IN PLACE OF chapter3.tex and
% chapter3_shortened.tex. Never include two of them: they share labels.
%
% WHAT CHANGED
%   * Repeated mechanism/derivation passages removed: they are owned by
%     Chapter 2 (magnification, brightness constancy, strain tensor, SimAM
%     energy, focal loss, F1). This chapter now carries published evidence and
%     transfer limits only, and contains NO display mathematics.
%   * The nine per-section "research gap" endings are replaced by short
%     relevance statements, with one qualified synthesis at the end.
%   * Local decision tables (EVM/flow/strain/TIM/CNN/SimAM/imbalance) removed:
%     they restated Chapter 2 or 4 and were then restated again in the ledger.
%   * 17 tables -> 6. 2 figures -> 1. The dataset-composition figure is gone
%     (its content is now one prose sentence); the fold-composition figure is
%     retained because the pooled-versus-fold argument depends on it.
%
% THE LEDGER'S FIRST COLUMN IS NOW \ref{}, NOT HARD-CODED NUMBERS.
% The twelve divergence rows previously cited 3.1.2 ... 3.9.7 literally, which
% any restructuring silently falsifies. They now resolve automatically.
%
% LABELS: every label referenced from introduction.tex, chapter4_shortened.tex,
% chapter5_shortened.tex or chapter6.tex is preserved (24 of them), plus those
% used by the unshortened chapters, so either manuscript state resolves.
% Verified by scratchpad scan; re-run it after any further edit.
%
% NOTE: the old header warned that Chapter 2 cites Chapter 3 subsections BY
% NUMBER. That is no longer true -- the shortened Chapter 2 has none -- so
% subsection positions here are free. The only by-number dependency was the
% ledger's own first column, now removed.
% =============================================================================

\chapter{Literature Review}
\label{ch:literature}

This chapter reviews the literature bearing on each component of the recognition pipeline evaluated in this thesis: the corpus every experiment runs on, the four techniques placed under ablation, the input transformations they operate on, and the training regime held constant across all of them. Each section states what the relevant studies found, what their evidence does not establish, and what follows for this experiment. Mechanisms already derived in \autoref{ch:background} are not derived again. \autoref{sec:synthesis-and-gap} consolidates the closing statements into one research gap and collects every declared departure from the reviewed methods.


% =============================================================================
\section{The Evaluation Corpus: CASME II and the Spontaneous Benchmark Family}
\label{sec:corpus}
% =============================================================================

The corpus is treated first because on a database of this size several decisions usually presented as modelling choices are in fact forced by properties of the data.

\subsection{Why a dedicated spontaneous corpus was necessary}
\label{sec:why-spontaneous}

Automatic expression recognition matured on large corpora of \textit{posed} expressions, reporting accuracies above 90\,\% on the six basic emotions by the early 2010s~\cite{yan2014}. Micro-expression research could not follow, because the phenomenon is involuntary and cannot be performed on instruction. Two problems motivated CASME II directly.

The first is \textbf{construct validity}: Yan et al.~\cite{yan2014} and Li et al.~\cite{liX2013} both observe that deliberately produced ``micro-expressions'' differ systematically from spontaneous ones, so a system trained on posed data may be a recogniser of fast deliberate movement rather than of micro-expressions. The second is \textbf{acquisition quality}. Defined jointly by very short duration and very low intensity, the phenomenon sits close to the noise floor of ordinary video; Yan et al.~\cite{yan2014} argue that CASME's 60~fps coarsely samples a movement lasting a fraction of a second --- while marking the benefit of higher rates as unresolved --- and that mains-frequency flicker contaminated earlier recordings. Elicitation was a third difficulty: deception paradigms produced samples contaminated by conversational movement, and the workable paradigm, which CASME II adopts, has participants watch high-arousal video while attempting to keep the face neutral~\cite{liX2013}.

\begin{table}[htbp]
	\centering
	\footnotesize
	\setlength{\tabcolsep}{4pt}
	\begin{tabular}{l c c c c c c}
		\hline
		\textbf{Corpus} & \textbf{Subj.} & \textbf{Micro samples} & \textbf{fps} & \textbf{Spont.} & \textbf{FACS} & \textbf{Ethn. / age} \\
		\hline
		SMIC~\cite{liX2013}              & 16 valid          & 164                & 100          & Yes          & No           & 3 / 26.7 \\
		\textbf{CASME II}~\cite{yan2014} & \textbf{26 valid} & \textbf{247 / 255} & \textbf{200} & \textbf{Yes} & \textbf{Yes} & \textbf{1 / 22.03} \\
		SAMM~\cite{davison2018}          & 32                & 159                & 200          & Yes          & Yes          & 13 / 33.24 \\
		CAS(ME)$^2$~\cite{qu2016}        & 22                & 53 (+250 macro)    & 30           & Yes          & Yes          & 1 / 22.59 \\
		\hline
	\end{tabular}
	\caption[The spontaneous micro-expression corpora]{The spontaneous corpora that define the current context, compiled from Yan et al.~\cite{yan2014}, Davison et al.~\cite{davison2018}, Li et al.~\cite{liX2013} and Qu et al.~\cite{qu2016}. CASME II carries two sample counts: 247 in the paper, 255 in the released coding.}
	\label{tab:spontaneous-corpora}
\end{table}

\subsection{CASME II: construction and annotation}
\label{sec:casme2-construction}

CASME II~\cite{yan2014} recruited thirty-five participants, mean age 22.03 (SD 1.60), recording at \textbf{200 fps} and 640~$\times$~480 without inter-frame compression under flicker-free illumination, giving a facial region of roughly 280~$\times$~340 pixels. Each retained sample carries onset and offset indices, FACS action-unit labels from two independent coders with disagreements arbitrated by discussion, and an emotion label assigned from the coded action units, participant self-report and the eliciting episode's content --- Yan et al. arguing against forcing spontaneous samples into the six basic categories, since a single action unit is frequently ambiguous.

Two specifics propagate into this thesis. The original release applies a normalisation from which this work departs: 68 Active Shape Model landmarks are detected on the \textit{first frame only} and the resulting transform applied unchanged to every subsequent frame; this pipeline reads the original recorded frames instead, unregistered and uncropped. And the published paper reports \textbf{247} micro-expressions while the distributed coding file contains \textbf{255}, a post-publication revision rarely acknowledged and one reason sample counts for ``CASME II'' vary. The present work uses the released 255-sample coding.

\subsection{The label taxonomy and the sparsity it creates}
\label{sec:label-taxonomy}

Yan et al.~\cite{yan2014} provide five emotion categories with action-unit criteria and are candid about the result: ``some types of facial expressions are difficult to elicit in laboratory situations, thus the samples in different categories distributed unequally.'' The released coding is finer, exposing seven raw labels. Counting directly from the label table used here gives happiness 32, disgust 63, surprise 25, repression 27, sadness 7, fear 2, and others 99.

\textbf{Seven-class classification is not viable on this corpus.} Fear is represented by two clips in the entire database and sadness by seven. Under any subject-disjoint protocol --- and subject-disjointness has been this benchmark's norm since its introduction (\autoref{sec:original-baseline}) --- a two-clip class cannot be meaningfully trained or tested, and per-class metrics degenerate.

The field's response is to group by affective valence. The convention originates with SMIC, released with exactly three classes~\cite{liX2013}, and was carried into cross-corpus work by MEGC 2019~\cite{see2019}. Discarding ``others'' and grouping the remainder gives the working set used throughout: \textbf{Negative 99, Positive 32, Surprise 25, N = 156}, a ratio of roughly 4 : 1.3 : 1.

\textbf{A declared divergence from the MEGC convention.} MEGC 2019 reports the CASME II portion of its composite database as \textbf{145} samples, footnoting that CASME II's Negative class covers ``Disgust and Repression''~\cite{see2019} and building its subset from the paper's five-category release, in which sadness and fear are absorbed into ``others''. This thesis uses the seven-label coding and assigns sadness and fear to Negative, following the mapping MEGC gives for \textbf{SAMM} rather than for CASME II. The consequence is an \textbf{eleven-clip difference}: N = 156 against 145, with Negative at 99 rather than 88. Nine of the eleven are the sadness and fear clips; the remaining two are unaccounted for, since disgust and repression together number 90 in the released coding against the 88 the challenge reports --- a residual noted rather than resolved. The direction is mild: the nine added clips all fall in the majority class and come from the categories coders found hardest to assign, so if anything they make this task marginally harder than the challenge's.

\subsection{The original baseline, and what it does not establish}
\label{sec:original-baseline}

Yan et al.~\cite{yan2014} baseline the corpus with LBP-TOP~\cite{zhaoG2007} and a support vector machine, sweeping spatial and temporal radii and excluding $R_T = 1$ deliberately, because at 200~fps the change between adjacent frames is negligible --- an early acknowledgement that a high sampling rate is not automatically an advantage. Five-class accuracies under \textbf{leave-one-subject-out} cluster in the high fifties to low sixties, the best configuration reaching \textbf{63.41\,\%}. Three observations matter.

\textbf{Leave-one-subject-out is native to this benchmark}, not a stricter protocol adopted later, so any evaluation reporting a single random split departs from the benchmark's own convention. \textbf{The baseline is reported as accuracy alone}, and on a five-class problem where ``others'' is 102 of 247 samples the majority-constant classifier scores 41\,\% without modelling anything, so a headline near 63.41\,\% is not by itself evidence of a working system. \textbf{And the sources disagree on the protocol behind that figure}: Yan et al.~\cite{yan2014} state leave-one-subject-out in abstract and results prose, the caption of the table reporting it reads ``leave-one-out cross-validation'' without the subject qualifier, and Li et al.~\cite{liX2018} reproduce the figure marked as leave-one-\textit{sample}-out, a materially easier protocol. The most widely quoted number on this benchmark thus rests on a protocol its own sources describe three ways; this thesis follows the originating paper's prose and records the discrepancy rather than resolving it.

Accuracy also moves by several points across a purely descriptive radius sweep, with no change to data, split or classifier. Where effect sizes of that magnitude arise from descriptor settings, similar differences between architectures cannot be interpreted without a variance estimate --- a limitation this thesis inherits and declares.

\subsection{Sibling corpora, and why this study remains single-corpus}
\label{sec:sibling-corpora}

Three further databases define the context, each isolating a specific limitation of CASME II. \textbf{SMIC}~\cite{liX2013} is the origin of the three-class taxonomy and of fixed-length temporal interpolation, both inherited here; it is not used as a corpus because it is neither FACS-coded nor apex-annotated. \textbf{SAMM}~\cite{davison2018} addresses the CASME family's demographic narrowness with 13 ethnicities against CASME II's one, Davison et al. observing that a corpus drawn from one ethnicity and one narrow age band restricts analysis to ``similar looking participants''; it is stronger on that count but smaller in usable samples and spread across seven categories. \textbf{CAS(ME)$^2$}~\cite{qu2016} pairs macro- and micro-expression samples from the same participants, making it the natural resource for domain transfer --- the strategy behind the strongest MEGC 2019 entry~\cite{liu2019,see2019} --- but records at 30~fps, supplies only 53 micro-expression samples, and is evaluated leave-one-\textit{video}-out, which does not enforce subject-disjointness.

CASME II is used alone because it has the highest temporal resolution of any spontaneous corpus of comparable size, a precondition for magnification in a defined band and for dense frame-to-frame flow; because its onset, apex and offset annotation makes temporal normalisation well-defined rather than heuristic; and because it remains the most widely reported single-corpus benchmark. The cost --- that no claim here extends to other populations, capture conditions or corpora --- is declared rather than argued away.

\subsection{How CASME II is evaluated: MEGC 2019 and the metric convention}
\label{sec:megc-metrics}

The most consequential development in how CASME II is \textit{used} is the Composite Database Evaluation protocol of the Second Facial Micro-Expression Grand Challenge~\cite{see2019}. It merges SMIC, CASME II and SAMM into one composite database of \textbf{442 samples} --- Negative 250, Positive 109, Surprise 83 --- of which CASME II contributes 145, and runs leave-one-subject-out across the composite. Plain accuracy is rejected explicitly on grounds of class imbalance, and two balanced metrics mandated: \textbf{Unweighted F1 (UF1)}, the macro-averaged F1 defined in \autoref{sec:confusion-matrix}, and \textbf{Unweighted Average Recall}, the balanced accuracy.

What matters is not the algebra but the \textbf{construction}: UF1 accumulates true positives, false positives and false negatives per class \textbf{over all folds} \textit{before} any per-class F1 is computed. As \autoref{sec:pooled-vs-fold} establishes, that is a different estimator from the average of per-fold macro F1 scores, and the distinction is not cosmetic here. Because clips are distributed very unequally across subjects and many single-subject folds contain only one of the three classes, a per-fold macro F1 is bounded well below one by fold composition alone (\autoref{fig:fold-composition}). The primary metric of this thesis, \textit{pooled macro F1}, is therefore identical in construction to UF1; the per-fold average appearing in the raw logs is not a metric the field uses.

\begin{figure}[htbp]
	\centering
	\includegraphics[width=\textwidth]{fig3_2_fold_composition.png}
	\caption[Class composition of the 25 leave-one-subject-out folds]{Class composition of the 25 leave-one-subject-out folds, and the per-fold macro-F1 ceiling it implies.}
	\label{fig:fold-composition}
\end{figure}

\textbf{One structural caveat governs the challenge's reported figures}, which range up to 0.7885 composite UF1 and 0.8764 on the CASME II subset~\cite{see2019}. Every one comes from a model trained on the \textit{composite} database --- roughly 442 samples across three corpora and 68 subjects --- then scored on the CASME II portion of the pooled predictions. A CASME II column is therefore not a CASME II-only result: it reports how a cross-corpus model performs on CASME II clips, with approximately \textbf{2.8 times} the training data available in any single-corpus fold, where the experiment conducted here trains each fold on roughly 150 clips. The two regimes are not interchangeable, and single-corpus results should not be compared against those figures without stating the difference.

Two further observations from the challenge bear on this design. The top three submissions all chose \textbf{optical flow} as their representation, a convergence that supports the motion-based input used here; and one submission substituted mid-position frames for annotated apex frames with little loss, suggesting that precise apex localisation may be less critical than assumed.

\subsection{Threats to validity carried by the corpus}
\label{sec:corpus-threats}

The following are not criticisms of CASME II, which remains the best available resource for this task, but constraints any study built on it inherits. Fold and duration statistics are computed from the 156-clip label table used here.

\textbf{Scale and imbalance.} 156 usable clips across three classes, distributed 99 : 32 : 25, small enough that parameter count is itself a risk factor. This motivates parameter-free or low-parameter components wherever a choice exists, and it makes accuracy unusable alone: the always-Negative classifier attains \textbf{0.6346 accuracy} at a macro F1 of only \textbf{0.2588}, so that floor is quoted alongside every result.

\textbf{Subject-level non-uniformity.} \textbf{Subject 17 alone contributes 33 clips, 21\,\% of the working set}, while three subjects contribute one each, so fold sizes vary by a factor of 33; \textbf{ten of the 25 folds contain a single class}, eight contain two, and only seven all three. This forces full 25-fold leave-one-subject-out and is why the pooled metric rather than the fold average is primary. It is also why a single held-out split is not a reliable instrument here: which subjects fall in the test set can move the result by more than any architectural change under study. Subject 18's three clips are all ``others'' and are removed by the three-class filter, so the run has \textbf{25 folds, not 26} --- stated explicitly because fold counts differ between papers using this corpus.

\textbf{Demographic homogeneity and laboratory ecology.} All participants are Chinese with mean age 22.03, named by Davison et al.~\cite{davison2018} as a principal motivation for SAMM; the MEGC organisers saw the effect in reverse, reporting that adaptation from CK+ helped on SMIC and SAMM but ``not so for CASME II, which contain predominantly Chinese subjects''~\cite{see2019}. Recording is frontal under fixed lamps with participants instructed not to move or speak, so almost none of the head motion, illumination variation, occlusion or speech-related movement a deployed system would face is present. Together these motivate a \textbf{motion} representation rather than raw pixels --- appearance being exactly what a 26-subject corpus would let a network memorise --- and they are why no cross-dataset or in-the-wild claim is made anywhere in this thesis.

\textbf{Duration, annotation and version ambiguity.} Onset-to-offset lengths range from \textbf{31 to 126 frames}, median 66 --- at 200~fps, 0.155\,s to 0.63\,s --- so every clip is resampled to a fixed length before training. Labels derive from coded action units, self-report and stimulus content: a construct produced by trained coders under a rubric, taken here as given and recorded as a declared limitation rather than as measurement error. And three sample counts circulate under the name ``CASME II'' --- 247, 255 and 145 --- with both five- and seven-class label sets in use and rarely stated, so N is reported explicitly alongside every figure here.

\subsection{Implications for this research}
\label{sec:corpus-implications}

Two further properties determine settings rather than threatening validity: the \textbf{200 fps} capture rate fixes the magnification pass-band at a matched \textbf{5--25 Hz}, and the baseline's sensitivity to a descriptive hyper-parameter sets the standard of evidence, so effect sizes below the resolution of a 156-clip corpus are reported as unresolved rather than as findings. No confidence interval is available, since the single-seed design and unretained per-clip predictions preclude one.

The consequence that shapes interpretation rather than construction is that \textbf{the small-sample regime predicts which components should help}. 156 clips cannot supply the data needed to learn a general spatio-temporal feature extractor from scratch, particularly once much of that extraction has been performed analytically before the network sees the data, whereas a component modelling the temporal arc of an expression exploits exactly the structure the corpus's onset, apex and offset annotation guarantees is present. This is a hypothesis the experiment tests, not a result it assumes. What the corpus literature does not supply is a stable basis for comparing methods: only Li et al.~\cite{liX2018} quantify how much of a reported difference is attributable to the protocol rather than the method, and then only for a single matched row of their own method.


% =============================================================================
\section{Motion Magnification: Eulerian Video Magnification}
\label{sec:evm-review}
% =============================================================================

Magnification is reviewed through the four corpus papers that apply it: Bai et al.~\cite{bai2021}, Li et al.~\cite{liX2018}, and Li, Huang and Zhao~\cite{liY2018,liY2021}. The original amplitude- and phase-based formulations lie outside the corpus and are described only as those papers attribute them. The mechanism itself is given in \autoref{sec:motion-magnification}.

The corpus design of \autoref{sec:why-spontaneous} answers the \textit{short-duration} half of the micro-expression definition but not the \textit{low-intensity} half. A movement spread across dozens of frames at 200~fps may still occupy only a few grey levels across a small patch of skin, and Yan et al.~\cite{yan2014} report discarding samples outright because the movement was too faint to place an onset. Li et al.~\cite{liX2018} state the remedy directly --- amplify the motion before any descriptor is computed, so that the difference between categories is enlarged rather than the classifier being made more sensitive. Being entirely upstream of the model is what makes magnification a clean ablation factor here.

Three parameters are free in any application: the temporal pass-band, which decides \textit{which} motions are amplified; the factor $\alpha$, which decides \textit{how much}; and the pyramid depth. On the first, Bai et al.~\cite{bai2021} report the rule given by the original authors --- the band should match the duration of the phenomenon, for micro-expressions 1/5 to 1/25\,s, equivalently \textbf{5--25 Hz} --- while noting that this full band produced ``noticeable noise'' in their setting, which led them to narrow it to 15--25 Hz. This is the only explicit guidance on band selection anywhere in the corpus.

\subsection{Evidence that magnification improves recognition}
\label{sec:evm-evidence}

The strongest case is Li et al.~\cite{liX2018}, who sweep the magnification factor rather than fixing it, across $\alpha \in \{1, 2, 4, 8, 12, 16, 20, 24, 30\}$ with $\alpha = 1$ denoting no magnification, on both SMIC and CASME II under leave-one-subject-out. Two findings emerge.

\textbf{Magnification helps consistently}, ``for all three kinds of features on all four testing datasets'', and on CASME II substantially (\autoref{tab:evm-effect}); their best full configuration reaches 67.21\,\% on CASME II, competitive with the contemporaneous state of the art.

\begin{table}[htbp]
	\centering
	\small
	\begin{tabular}{l c c c}
		\hline
		\textbf{Descriptor} & \textbf{No magnification} & \textbf{With magnification} & \textbf{$\Delta$} \\
		\hline
		LBP  & 55.87\,\% & 60.73\,\% & \textbf{+4.86} \\
		HOG  & 57.49\,\% & 63.97\,\% & \textbf{+6.48} \\
		HIGO & 57.09\,\% & 67.21\,\% & \textbf{+10.12} \\
		\hline
	\end{tabular}
	\caption[Effect of Eulerian motion magnification on CASME II]{Effect of magnification on CASME II under leave-one-subject-out, after Li et al.~\cite{liX2018} (Table 6). This is the corpus's one controlled magnification-on/off comparison.}
	\label{tab:evm-effect}
\end{table}

\textbf{The relationship between $\alpha$ and accuracy is non-monotonic.} Li et al.~\cite{liX2018} describe the curves as ``rainbow-shaped'', with best performance ``generally achieved when the motion is magnified in the range of $[8, 16]$'', and give a mechanistic account of both failure modes: too little magnification ``might not be enough to reveal the ME motion progress'', while too much ``degrades the performance because too many artifacts are induced''. Magnification is therefore a parameter with an interior optimum, not a monotone improvement to be maximised.

Bai et al.~\cite{bai2021} provide a second demonstration with a fundamentally different downstream model --- a pre-trained VGGFace2 encoder feeding a bidirectional LSTM --- reaching \textbf{75.76\,\% with phase-based magnification against a 60.60\,\% baseline} on SMIC, with the phase-based variant outperforming the amplitude-based one. The comparison is a distinct one and worth retaining, but the baseline is drawn from an earlier publication rather than from a magnification-off arm of their own, so it is not a matched control in the sense of \autoref{tab:evm-effect}. Because the pipeline shares nothing with hand-crafted descriptors, the two results together indicate the benefit is not an artefact of one feature family --- at least when the downstream representation is appearance-based.

Two further papers apply magnification in a different regime. Li, Huang and Zhao~\cite{liY2018,liY2021} magnify a single spotted apex frame at \textbf{$\alpha = 30$}, at the extreme end of Li et al.'s sweep, and augment by retaining the apex plus two frames either side. That $\alpha$ is asserted rather than derived, and the regimes differ enough --- a narrow window consumed as largely static appearance, against full magnified sequences described by spatio-temporal descriptors --- that the discrepancy need not be a contradiction. The safe reading is the weaker one: $\alpha$ is a parameter of the \textit{combination} of magnification and downstream representation, and cannot be transferred between pipelines without re-validation.

\subsection{Known limitations}
\label{sec:evm-limitations}

Four failure modes are documented, and the corpus is unusually candid about them.

\textbf{Magnification amplifies everything in band, not only expressions.} This is the most important and the most easily overlooked. Li et al.~\cite{liX2018} decline to use magnification for micro-expression \textit{spotting} at all, ``because it magnifies unwanted motions (e.g. head movements) at the same time''. A rigid head translation whose temporal frequency falls inside the pass-band is amplified exactly as a brow raise is. Prior alignment and cropping mitigate this; they do not eliminate it.

\textbf{Artefacts grow with $\alpha$}, reported by all four papers. \textbf{There is no principled procedure for choosing $\alpha$}: one value in the corpus is obtained by exhaustive sweep~\cite{liX2018}, and the apex-frame papers assert theirs. And \textbf{evaluation is confined to appearance-based downstream representations} --- LBP, HOG and HIGO in Li et al.~\cite{liX2018}, VGGFace2 features in Bai et al.~\cite{bai2021}, VGG-Face in both Li, Huang and Zhao papers. Every one is computed on magnified pixel intensities. \textbf{No study in the corpus measures what magnification contributes when the representation handed to the classifier is itself an explicit motion field.}

\subsection{Implications for this research}
\label{sec:evm-implications}
\label{sec:evm-motion-hypothesis}

The review fixes two settings and leaves one hypothesis. The temporal band is set to the \textbf{unnarrowed 5--25 Hz} duration rule, so Bai et al.'s narrowing is not followed and the pass-band admits the low-frequency range they rejected as noisy; this is a declared divergence. The factor is set to \textbf{$\alpha = 10$}, inside the reported $[8,16]$ optimum, and the apex papers' $\alpha = 30$ is explicitly not adopted because this thesis consumes full sequences. Amplitude-based magnification is used over a four-level Laplacian pyramid, and the band-pass is realised by an \textbf{FFT coefficient mask} rather than the Butterworth filter the reviewed description uses (\autoref{sec:temporal-filtering}) --- a further declared difference. Stage 1 runs twice, producing magnified and non-magnified tensor sets, so the switch selects a directory and leaves the network untouched.

Behind those settings lies a hypothesis the corpus does not test, stated because it frames the interpretation. Magnification amplifies the \textit{amplitude} of intensity variation in a band, and an appearance descriptor on magnified frames sees a directly stronger signal. An optical flow field is a different object: flow estimates \textit{displacement}, and amplifying the underlying intensity change by $\alpha$ does not simply scale the recovered displacement by $\alpha$. It changes the conditioning of the estimation problem, potentially making faint motion recoverable where it was not, while introducing displacement artefacts of the kind all four papers warn about; and any residual global scaling is further attenuated by the per-channel normalisation applied before training. The plausible outcome is that magnification in front of a flow-and-strain representation retains the part of its benefit consisting of \textit{making otherwise-unrecoverable motion recoverable} while losing the part consisting of \textit{increasing contrast for an appearance descriptor} --- so its measured contribution should be smaller on average than the +5 to +10 points of \autoref{tab:evm-effect}. Six matched pairs differing in nothing but the tensor directory are what allow that to be tested.

\textbf{A defect the review made detectable.} In this project's earliest evaluations every magnification-enabled configuration produced results \textit{bit-identical} to its disabled counterpart. Given that all four reviewed papers report clear, repeatable effects, an effect of exactly zero across six independent pairs is not a plausible null but a symptom: both arms were reading the same tensor directory, so the hypothesis had never been exercised. The defect was repaired before the reported run, and wiring the existing verification helper in as a start-up assertion is declared further work.

What the literature leaves open is threefold: every evaluation places magnification in front of an appearance representation, conflating two mechanisms that no experiment separates; $\alpha$ is transferred between studies without re-validation; and magnification is always evaluated as a two-cell addition to an otherwise fixed pipeline, so its interaction with other components is never measured. This thesis places it in front of an explicit motion representation, selects $\alpha$ from the full-sequence optimum, and measures it as one factor in a factorial ablation whose spread across architectures is itself the finding.


% =============================================================================
\section{Motion Representation: Dense Optical Flow}
\label{sec:flow-review}
% =============================================================================

Dense optical flow supplies the horizontal and vertical displacement channels. The relevant literature concerns three choices: the estimator, the frame pairs it processes, and whether flow serves as a recognition feature or an alignment tool. The estimation assumptions themselves are given in \autoref{sec:optical-flow}.

Xu et al.~\cite{xu2017} motivate facial-dynamics descriptors by the difficulty of interpreting micro-expressions through texture alone: flow describes apparent movement directly, whereas image intensity also records appearance and illumination. This does not guarantee independence from identity or recording conditions, since facial geometry, alignment and brightness changes still affect estimation. The benchmark evidence supports the choice --- the MEGC 2019 organisers report that their three highest-ranked submissions used optical flow~\cite{see2019}, OFF-ApexNet~\cite{liong2019a} uses horizontal and vertical flow as network inputs, STSTNet~\cite{liong2019b} adds strain, and Zhao et al.~\cite{zhaoS2021} retain a sequence of flow fields --- but this convergence is not a controlled demonstration that motion always outperforms appearance.

\textbf{The estimator is not settled by the corpus.} Zhao et al.~\cite{zhaoS2021} adopt Farneb\"ack through OpenCV, citing established availability; OFF-ApexNet~\cite{liong2019a} uses TV-L1, citing preservation of motion discontinuities; Xu et al.~\cite{xu2017} use the method of Sun, Roth and Black. STSTNet~\cite{liong2019b} specifies onset--apex flow without naming its estimator, so its similarity to OFF-ApexNet is insufficient to attribute TV-L1 to it. The advantages claimed for TV-L1 should not be read as a measured recognition gain over Farneb\"ack on CASME II: no controlled three-estimator comparison exists in the corpus. This thesis fixes Farneb\"ack as a preprocessing choice and leaves the comparison open.

\textbf{Onset--apex flow or a sequence is the consequential choice.} OFF-ApexNet~\cite{liong2019a} and STSTNet~\cite{liong2019b} represent each clip by flow from onset to apex, resized to $28\times28$ --- an economical summary of the displacement from neutral to peak, suitable for a small training set, which omits the intermediate trajectory and the relaxation after the apex. Zhao et al.~\cite{zhaoS2021} instead build an eleven-frame key-frame sequence and compute ten adjacent-pair fields, retaining changes in motion across the expression. The distinction is between one displacement summary and a sequence whose order a temporal model can exploit; it is not merely a difference in input size. This study follows the sequence approach but samples existing frames uniformly rather than constructing key frames, which preserves a temporal input for the transformer ablation. It does not measure the recognition cost of replacing a sequence with an onset--apex pair: that comparison is not an experimental arm.

Flow is also used for alignment rather than recognition. Xu et al.~\cite{xu2017} coarsely align by landmarks and then refine in-sequence using pixel-level movement, motivated specifically by micro-expressions, where a small registration error can be comparable to the displacement being recognised. This pipeline performs no flow-based fine alignment, so corpus preprocessing must not be treated as evidence that this error source has been eliminated. It remains the most promising unexploited improvement to preprocessing.

\textbf{Limitations of the representation.} Flow remains sensitive to brightness-constancy violations, alignment error and estimator settings; apex-based systems additionally depend on selecting an informative frame, and OFF-ApexNet~\cite{liong2019a} attributes weaker SMIC performance to imprecise apex selection. Using an onset--offset sequence avoids dependence on one selected apex but still relies on clip boundaries and sampling: repeated sampled frames yield near-zero-flow pairs, and large sampling intervals make the small-motion approximation less appropriate.

\subsection{Implications for this research}
\label{sec:flow-implications}

The central interpretive point is that optical flow is an \textit{analytical spatio-temporal feature extractor}: it estimates motion from paired frames before the network is trained. A learned spatial stem therefore acts on a representation that already makes displacement explicit, and with 156 clips, learning additional spatial filters may offer limited benefit or duplicate structure preprocessing has already exposed. This is a hypothesis about the stem's marginal value, not proof that learned spatial processing is unnecessary.

A temporal encoder addresses a different task: modelling relations among successive displacement fields, which pairwise flow estimation does not itself learn. Keeping those two roles distinct is necessary when attributing performance to the stem as against the transformer, and it is what the matched comparisons of this thesis are designed to separate. The design does not rank estimators, isolate alignment quality, or compare onset--apex against sequence representations; those remain limitations rather than gaps claimed as closed.


% =============================================================================
\section{Deformation Representation: The Optical Strain Tensor}
\label{sec:strain-review}
% =============================================================================

Optical strain supplies the third input channel. Its literature establishes a deformation descriptor derived from flow, several ways to use it, and limits relevant to interpreting the ablation. The tensor and its components are defined in \autoref{sec:optical-strain}; this section concerns the published evidence and one consequential divergence.

Flow describes displacement; strain describes its local spatial variation. A spatially uniform translation has zero displacement gradient, whereas neighbouring skin regions moving differently produce deformation. Shreve et al.~\cite{shreve2011} use this distinction to detect facial activity and describe robustness to moderate head translation, which does not imply immunity to arbitrary head movement or registration error. Their illumination- and make-up-robustness statements are attributed to earlier studies rather than measured anew, and should be read as reported motivation rather than as evidence about the present pipeline.

\subsection{The implemented magnitude and its divergence}
\label{sec:strain-divergence}

Liong et al.~\cite{liong2014a,liong2014b,liong2016} reduce the tensor to a scalar by the four-component Frobenius norm, which counts the two equal off-diagonal shear entries separately. \textbf{The implementation here counts shear once rather than twice}, giving the three-term magnitude of \autoref{sec:optical-strain}. Spatial derivatives come from central differences at one-pixel spacing, following Liong et al.~\cite{liong2014a}, but the scalar reduction differs from their convention.

STSTNet~\cite{liong2019b} cannot establish a precedent for the three-term form. Its printed magnitude equation repeats $\partial u$ in the mixed derivative where the symmetric tensor calls for $\partial v/\partial x$ alongside $\partial u/\partial y$, and it places the factor one-half \emph{outside} the squared derivative sum. Correcting the repeated derivative while retaining that placement yields the \textit{four}-term convention; moving the factor inside the square would instead yield the implemented form, but that is a further change to the printed expression. The intended implementation cannot be settled from the published equation alone, so this thesis declares a departure from the unambiguous four-term convention rather than claiming to reproduce STSTNet's magnitude.

The four-term form weights pure shear by $\sqrt{2}$ relative to the implemented one. Independent min--max normalisation removes a uniform scale factor but does not generally remove this difference, because normal and shear contributions vary across pixels and time. Neither convention was compared experimentally here, so their practical equivalence is not assumed.

\subsection{Recognition evidence and its protocol limits}
\label{sec:strain-evidence}

Strain has been used three ways: as a pooled descriptor (optical strain features, OSF), as a spatial weighting on LBP-TOP histograms (optical strain weighted features, OSW), and as a direct input channel. Liong et al.~\cite{liong2016} combine the first two, pooling spatially then temporally, and report the results of \autoref{tab:strain-performance}; the early work also applies Wiener and Gaussian filters explicitly to suppress noise~\cite{liong2014a}. STSTNet~\cite{liong2019b} supplies strain directly alongside onset--apex flow, which is the precedent for the three-channel input used here, independently of the magnitude ambiguity --- though this study retains all three channels across successive frame pairs rather than reducing each clip to one map.

\begin{table}[htbp]
	\centering
	\small
	\begin{tabular}{@{}lcc@{}}
		\hline
		\textbf{Method} & \textbf{SMIC accuracy (\%)} & \textbf{CASME II accuracy (\%)} \\
		\hline
		Baseline (texture) & 45.73 & 61.94 \\
		OSF                & 41.46 & 51.01 \\
		OSF + OSW          & 52.44 & \textbf{63.16} \\
		OFF + OFW (flow)   & 40.24 & \textbf{55.87} \\
		\hline
	\end{tabular}
	\caption[Strain-based descriptors against a matched flow control]{Strain-based descriptors against a matched flow-based control, after Liong et al.~\cite{liong2016} (Table 8). SMIC uses leave-one-subject-out; CASME II uses leave-one-\textit{video}-out, which is not subject-disjoint. The summary combines block settings, so equal entries elsewhere in the source are not matched-setting comparisons.}
	\label{tab:strain-performance}
\end{table}

The CASME II evidence depends on how strain is used and against what. OSF alone falls below the texture baseline; the combined descriptor improves on it by 0.41 to 1.22 points across block settings~\cite{liong2016}. Crucially, Liong et al. also test flow features and weights within the same framework, reaching 55.87\,\% against 63.16\,\% for the strain pair --- a \textbf{matched flow-versus-strain control}, so it would be incorrect to claim the literature never isolates strain against flow. What it concerns is hand-crafted descriptors under a protocol that is not subject-disjoint, and these accuracy gains are not comparable with pooled macro F1. It does not establish the gain from adding a strain channel to a learned sequence model.

\subsection{Implications for this research}
\label{sec:strain-implications}

Strain extends the analytical-feature argument of \autoref{sec:flow-implications}: the input contains a \textit{hand-specified differential operator}, exposing local deformation without learning that operator from 156 clips. A spatial stem may still learn useful combinations of those features, but it operates after both motion estimation and deformation extraction. This motivates testing its incremental contribution; it does not establish the cause of any observed architectural effect.

Three limits bound what this thesis can say about strain. Spatial differentiation can amplify flow-estimation noise, and \textbf{no dedicated strain post-filter is applied}, departing from the explicit filtering of Liong et al.~\cite{liong2014a}; Farneb\"ack's internal smoothing is not a separately evaluated strain filter. The shear weighting is the untested choice of \autoref{sec:strain-divergence}. Above all, \textbf{strain is present in all twelve evaluated configurations}, so its marginal contribution is not measured at all: a strain-on/strain-off comparison is the natural next ablation, and neither published gains from a different descriptor nor the performance of the complete system substitutes for it. A deterministic transformation of $u$ and $v$ adds no independent observation; its potential value lies in making deformation directly accessible to the classifier, and whether that helps must be measured rather than inferred from its construction.


% =============================================================================
\section{Temporal Normalisation: Fixed-Length Sampling}
\label{sec:tim-review}
% =============================================================================

Flow and strain fix \textit{what} is computed from each frame pair; temporal normalisation fixes \textit{which} pairs, and how many. The Temporal Interpolation Model (TIM) is credited in this corpus to Zhou et al. by Li et al.~\cite{liX2013,liX2018} and to Pfister et al. by Lu et al.~\cite{lu2015}, and is described here only through the reviewed papers that implement it.

Normalisation is necessary because any batched network requires a fixed input shape and the corpus supplies no constant, spanning a four-fold range of sequence length; and because duration is itself a nuisance variable --- Lu et al.~\cite{lu2015} normalise to remove ``temporal fluctuations of micro-expressions, which may introduce noise irrelevant to classification''. They also give the corpus's most compact description of TIM: ``a manifold-based interpolation method, which builds a low-dimensional embedding of an image sequence, and then interpolates a curve in the low-dimensional space. The interpolated frames are mapped back to the original high-dimensional space.'' The property that matters downstream is that TIM \textbf{synthesises frames that were never recorded}, having been proposed ``to generate sufficient frames from an image sequence'' --- an upsampler rather than only a length-standardiser.

\textbf{Ten frames is a convention, not a finding, and the studies that test it disagree.} TIM appears throughout the corpus as an unremarked preprocessing step, almost always at ten frames~\cite{liX2013,liX2018,lu2015}. Li et al.~\cite{liX2018} sweep interpolation length from 10 to 80 on the SMIC variants with the descriptor held fixed, and report that ``the best performance for all three datasets is achieved with interpolation to 10 frames'', because ``in longer sequences, changes along the time-dimension are diluted''. Ben et al.~\cite{ben2021} reach the opposite conclusion: holding alignment, descriptor and classifier fixed, they interpolate to 10 through 110 frames by both Newton interpolation and TIM on MMEW, finding Newton peaking at 33.3\,\% at 60 frames and TIM at \textbf{38.9\,\%} at 30 or 60, and adopt TIM at 60. Accuracy is non-monotonic in length in both. Xu et al.~\cite{xu2017} add the only multi-corpus sweep --- 10, 15 and 20 frames across CASME, CASME II and both SMIC subsets --- and depart from convention by using \textbf{linear} interpolation ``instead of Temporal Interpolation Model because in a high frame-rate scenario, linear interpolation suffices to characterize the motion pattern with less error''.

Two caveats limit what transfers. Ben et al.'s experiment uses MMEW with LBP-TOP, not CASME II and not a learned model, and runs largely in the \textit{upsampling} regime, so an optimum measured while synthesising frames need not hold when discarding them. And \textbf{frame synthesis is in tension with motion analysis}: where the downstream representation is appearance-based, synthesised frames are benign, but where it is motion, synthesised frames yield synthesised displacements. Xu et al.~\cite{xu2017} are the only reviewed paper to raise this, and they assert it rather than measure it. No reviewed study compares the two interpolators against a motion representation.

\subsection{Implications for this research}
\label{sec:tim-implications}

\textbf{A declaration about what this pipeline implements.} The module performing this step is named for the Temporal Interpolation Model, but what it computes is \textbf{uniform index sampling} of existing frames (\autoref{sec:three-channel-tensor}). No embedding is built, no curve is interpolated, no frame is synthesised. That is a materially different operation from TIM, and it is named accurately here rather than in the pipeline's own terms. It is defensible: because the downstream representation is flow between adjacent sampled frames, interpolation would manufacture displacement fields for motion that never occurred, whereas subsampling discards temporal resolution without fabricating it. What is given up is TIM's ability to \textit{lengthen} short clips.

The length is fixed at \textbf{33 frames sampled, giving 32 flow pairs}, bounded by the annotated onset and offset so the window spans the expression rather than surrounding neutral frames. That is close to the lower of Ben et al.'s two reported optima, but the correspondence should not be overstated: their 30 frames are \textit{interpolated} frames under a hand-crafted descriptor on a different corpus, which is not direct validation of 32 sampled flow fields under a learned model here. The length is not swept, and a sweep is declared as further work.

Three quantitative consequences follow, and they constrain the magnification conclusion rather than merely describing the sampler.

\textbf{Effective temporal resolution falls by a clip-dependent factor.} Sampling 33 frames from a clip of $N$ frames recorded at 200~fps gives an effective rate of $6600/N$ fps: \textbf{100 fps} at the median $N = 66$, which is SMIC's recording rate and half what CASME II was built to provide, and roughly \textbf{52 fps} for the longest clip, below the original CASME database's 60~fps~\cite{yan2014}. Eighty-one of the 156 clips are downsampled by a factor of two or more, which gives back part of what the 200~fps capture rate bought. No reviewed paper reports this cost.

\textbf{Magnification is applied after subsampling, which shifts the realised pass-band.} Magnification runs on the 33 sampled frames at the nominal 200~fps, but those frames span the clip's full duration, so the true inter-frame interval is $N/(200 \times 32)$ seconds rather than $1/200$. At the median, a band specified as 5--25 Hz corresponds to roughly 2.5--12.5 Hz physically, and the discrepancy grows with clip length. The realised band is therefore neither the 5--25 Hz that \autoref{sec:evm-implications} derives nor constant across clips. This is a declared divergence; magnifying before subsampling is declared as further work.

\textbf{Three clips are shorter than the sampling window.} Three of the 156 contain 31 frames against a target of 33, so sampling repeats indices and yields near-zero flow at those positions.

The unresolved question this leaves is what temporal normalisation costs a representation made of motion rather than appearance. Three studies sweep interpolation length and give three answers, only one of them touching CASME II and none using a learned temporal model. This thesis takes the position that follows from its own representation and quantifies what that costs in effective frame rate, but does not sweep the length, compare interpolators, or separate the effect of sequence length from that of the sampling regime.


% =============================================================================
\section{The Learned Spatial Backbone: Shallow Convolutional Networks}
\label{sec:cnn-review}
% =============================================================================

By 2019 the field had converged on optical flow as the input; the open question was what should consume it. A learned backbone promises what hand-crafted descriptors cannot~\cite{liong2019b}, but runs into the small-data constraint of \autoref{sec:corpus-threats}.

\textbf{The term ``3D CNN'' covers three materially different architectures in this literature.} STSTNet~\cite{liong2019b}, the canonical shallow model, calls itself a ``Shallow Triple Stream Three-dimensional CNN'', but its input is a single $28 \times 28 \times 3$ tensor whose third dimension holds horizontal flow, vertical flow and strain --- \textbf{not time}. Its three parallel single-layer streams, with 3, 5 and 8 kernels, supplement ``the small scale input data \ldots to avoid the problem of underfitting''; the clip was collapsed to one frame pair upstream. STRCN~\cite{xia2020a} by contrast extends convolutional connectivity genuinely ``across temporal domain'' with recurrent convolutional layers. Zhao et al.~\cite{zhaoS2021} apply a C3D-style network to a ten-frame flow sequence, so the temporal axis is present but short. Results reported across the corpus under the same label are therefore not like with like.

\textbf{The central controlled finding is that shrinking both model and input helps.} Xia et al.~\cite{xia2020b} treat depth and input resolution as controlled variables, building four models from a single convolutional layer to one convolutional plus three recurrent convolutional layers, and evaluating each across input resolutions from 20 to 300 pixels on optical flow. Deeper is worse: models 3 and 4 ``have worse recognition performances compared to shallower \ldots models'', even though all four are far shallower than ResNet-101. Higher resolution is worse with a threshold: performance ``will be degraded with larger input resolutions'' and for the deeper models ``begins to decrease dramatically when the resolution is bigger than \textbf{100 $\times$ 100}'' --- though shallow models are ``basically robust to the change of input resolution''. That qualification matters for this thesis, and so does the setting: Xia et al.'s study and STSTNet are both developed for the MEGC composite regime, which supplies substantially more training data per fold than a single corpus (\autoref{sec:megc-metrics}), so conclusions about affordable capacity understate the constraint here. Xia et al.~\cite{xia2020b} also take a third route to capacity control that this thesis follows --- developing ``three parameter-free modules \ldots integrate with RCN without increasing any learnable parameters'', so representational power grows while overfitting risk does not.

\textbf{No study isolates the backbone itself.} The corpus is not short of matched comparisons: it sets descriptor-only methods against learned ones on the same folds, and Xia et al.~\cite{xia2020b} vary individual architectural modules against a common baseline with every other parameter held constant. But those comparisons all \textit{add} to a backbone that stays in place. None removes the learned spatial stem entirely and passes the motion representation straight to a classifier or temporal model, so what the stem contributes against what it costs is nowhere measured.

\subsection{Implications for this research}
\label{sec:cnn-implications}

The component ablated here is a \textbf{three-stream shallow convolutional backbone}, one unshared two-layer stream per input channel, carrying \textbf{14,544 parameters} --- shallow by the corpus's standards, though roughly nine times STSTNet's. Where STSTNet's streams differ in filter count over the same input, these differ in input channel. Removing it leaves not an empty slot but spatial pooling followed by a learned linear projection, so the no-backbone route still contains learned input processing.

\textbf{A declaration about what this component is.} Every kernel in the backbone is spatial only: no filter spans two time steps, so despite the \texttt{Conv3d} implementation it is a \textbf{purely spatial} feature extractor (\autoref{sec:conv3d}). The ablation therefore does not test spatio-temporal convolution; it tests a learned spatial stem, and whatever it measures is evidence about \textit{this} stem --- spatial-only, at 224~$\times$~224, on flow and strain input --- and \textbf{not} evidence that convolution over time is unhelpful for micro-expression recognition, which this study never evaluated. In the field's own terms the choice is defensible: STSTNet, one of the strongest shallow baselines on the CASME II subset, likewise performs no temporal convolution.

\textbf{Input resolution is the clearest departure from Xia et al.'s guidance.} The backbone is fed at \textbf{224 $\times$ 224}, more than twice their threshold. This is not a departure from the corpus as a whole --- Bai et al.~\cite{bai2021} resize to 224 $\times$ 224 without alignment or cropping, as here; Li, Huang and Zhao~\cite{liY2021} resize faces to the same; and Xia et al.'s own earlier work~\cite{xia2020a} runs one stream at 300 $\times$ 245 --- but it departs from the specific threshold above, which is a composite-database result that bites hardest on models deeper than this one. Since this stem is shallow, the reviewed evidence predicts a smaller resolution effect than the deep-model figures suggest. The choice is also the expensive one: a 224 $\times$ 224 map carries sixty-four times the spatial elements of STSTNet's input, and the backbone convolves it at every one of the 32 time steps. Any effect measured for the stem is therefore \textbf{a resolution result as much as an architecture result}, and the reviewed evidence does not determine which mechanism dominates; re-running the backbone arm at a lower resolution is the first item of declared further work because it is cheap and would settle the question, not because the corpus implies a particular outcome.


% =============================================================================
\section{Parameter-Free Attention: SimAM}
\label{sec:simam-review}
% =============================================================================

SimAM re-weights convolutional feature maps without adding learned parameters. Yang et al.~\cite{yang2021} supply the mechanism and its image-recognition evidence; Xia et al.~\cite{xia2020b} motivate parameter-free modules in composite-database micro-expression recognition, supporting the design principle rather than validating SimAM on this task. The energy formulation and gate are given in \autoref{sec:simam}.

Yang et al.~\cite{yang2021} motivate attention through spatial suppression: neurons distinguishable from their surroundings are assigned greater importance, expressed as the linear separability of one activation from the others in its channel. This is a prior about activation statistics, not a detector of action units or minority classes --- a distinctive response may encode useful deformation, noise or an artefact --- so its suitability for flow-derived features requires evaluation rather than following from the neuroscience motivation. Two implementation details are owned here rather than by \autoref{ch:background}: the shared variance printed with the paper's equation~(5) divides by $M$ whereas its Figure~3 reference code divides by $M-1$, and this thesis follows the code; and the module's ``three-dimensional weights'' describe $C \times H \times W$ with $M = HW$ per frame, not a temporal volume, so extending the context across frames is an additional design decision rather than part of the original formulation.

Yang et al. evaluate SimAM with several network families on CIFAR and ImageNet, reporting competitive performance against learned attention modules, and their timing results show that unchanged parameter totals do not imply unchanged inference speed. For CIFAR they sweep $\lambda$ from $10^{-1}$ to $10^{-6}$ and select $10^{-4}$; they repeat the selection for ImageNet and choose $0.1$. Thus $10^{-4}$ is a dataset-selected operating point, not a universal setting justified by the formula.

\textbf{The limits of transfer are specific.} The source experiments concern natural images, different training-set sizes and learned feature distributions. Here the input is flow and strain, so distinctiveness need not bear the same relationship to useful evidence. The image experiments also do not decide whether video statistics should be computed per frame or over a clip. And the absence of learned attention weights does not establish that runtime, memory or overfitting concerns disappear. These are empirical limits on transferring the result, not reasons to assume success or failure.

\subsection{Implications for this research}
\label{sec:simam-implications}

SimAM is applied separately to each convolutional stream before concatenation, introducing no trainable weights or biases anywhere. Two choices diverge from the source. Mean and variance span $(T,H,W)$ within each channel with divisor $THW - 1$, rather than one $(H,W)$ plane, so an activation is judged distinctive relative to its whole clip rather than its own frame; this could emphasise a brief contraction, but does not guarantee low weights at onset or high weights at the apex, since the score depends on squared deviation and not on expression labels. And $\lambda = 10^{-4}$ is retained without re-tuning on CASME II. Neither context choice is compared experimentally.

SimAM is also enabled only alongside the convolutional stem in this model, which describes this implementation rather than a mathematical prohibition. Four of the sixteen candidate configurations are therefore excluded, leaving twelve evaluated and \textbf{four SimAM matched pairs}, and the measured effect is conditional on the stem being present. Retention should be judged against both recognition performance and computational cost: parameter count alone cannot settle it, because the module is cheap rather than free.

\subsection{Relevance and scope}
\label{sec:simam-gap}

The study evaluates SimAM in a motion-input, small-corpus setting through four matched comparisons under complete leave-one-subject-out, testing the implemented module as a whole rather than establishing the benefit of each design choice. Three follow-ups remain: whether the distinctiveness prior suits flow-derived features, re-tuning $\lambda$, and comparing per-frame with whole-clip statistics. Each requires new evaluations despite the module introducing no learned weights.


% =============================================================================
\section{The Temporal Encoder: Transformers}
\label{sec:transformer-review}
% =============================================================================

The temporal encoder is motivated by Zhang et al.'s Short and Long Range Relation Based Spatio-Temporal Transformer (SLSTT)~\cite{zhang2022} and its Vision Transformer foundation~\cite{dosovitskiy2021}. The component evaluated here differs substantially from published SLSTT, particularly in the axis over which self-attention operates, and distinguishing the two is the purpose of this section.

A micro-expression evolves through onset, peak and relaxation, so retaining a motion sequence permits a model to relate separated stages rather than compressing them immediately into one descriptor; Zhang et al.~\cite{zhang2022} motivate their design through the need to capture both local and global patterns. This thesis compares a temporal transformer against mean pooling over the same sequence: mean pooling discards order, whereas positional information and self-attention allow relations among frame-pair features to be modelled. That motivates the comparison without establishing in advance which generalises better. ViT itself~\cite{dosovitskiy2021} splits an image into patches, projects them into tokens, adds learned positional embeddings and relates them by attention; Dosovitskiy et al. also show that without the locality and translation-equivariance biases of convolution, image transformers can perform poorly when trained with insufficient data, while large-scale pre-training improves them. That is a caution for a 156-clip study, not a universal requirement that every transformer be pre-trained.

\textbf{Published SLSTT combines three stages.} First, \textbf{onset-referenced optical flow} describes sampled frames relative to the onset rather than to their immediate predecessors, which Zhang et al. argue gives a more structured trajectory towards and away from the apex; their experiments use it throughout, without a matched short-term-flow ablation establishing its independent gain. Second, a \textbf{spatial ViT encoder} processes each flow field as patches, using ImageNet-pre-trained ViT-B/16 with twelve layers at $384 \times 384$. Third, an \textbf{LSTM temporal aggregator} combines the per-frame features, with a mean aggregator supplying the comparison variant. Self-attention is spatial in this architecture; recurrence handles time.

\subsection{The aggregation comparison, and what it can support}
\label{sec:slstt-results}

The comparison relevant here is internal to SLSTT and holds the spatial transformer fixed in both arms: CASME II UF1 rises from \textbf{0.844} with the mean aggregator to \textbf{0.901} with the LSTM, a difference of 0.057~\cite{zhang2022}. That supports evaluating aggregation choices within that framework. It is \textbf{not} an estimate of what an LSTM would add to this thesis, where mean pooling follows \emph{temporal} self-attention rather than independent spatial encoding, and it isolates an aggregator rather than the presence of a transformer, since both variants contain one.

Protocol labels need care. Zhang et al. describe three-class leave-one-subject-out evaluations on the databases separately and together under a Composite Database Evaluation heading, so that heading alone does not establish that every CASME II entry uses pooled cross-database training; they also report a distinct Sole Database Evaluation using original labels, including five CASME II classes~\cite{zhang2022}. Neither setting should be equated automatically with this thesis's 156-clip three-class pool. Published SLSTT also retains a pre-trained spatial encoder, so its results do not isolate the benefit of a transformer trained from scratch on this study's motion representation.

\subsection{Implementation and six declared divergences}
\label{sec:transformer-implications}

The evaluated component receives 32 tokens of width 96, from either the convolutional stem or spatial pooling followed by a learned linear projection. The encoder uses four pre-norm layers, eight heads, feed-forward width 256 and dropout 0.1; fixed sinusoidal positions precede attention and mean pooling follows it. With the encoder disabled the input tokens are averaged directly.

\begin{table}[htbp]
	\centering
	\small
	\begin{tabular}{@{}p{0.23\linewidth} p{0.34\linewidth} p{0.35\linewidth}@{}}
		\hline
		\textbf{Property} & \textbf{Published SLSTT~\cite{zhang2022}} & \textbf{This thesis} \\
		\hline
		Attention axis   & Spatial patches within a frame              & Temporal tokens within a clip \\
		Aggregation      & LSTM over spatially encoded frames          & Mean after temporal self-attention \\
		Initialisation   & ImageNet-pre-trained ViT-B/16               & Trained from scratch \\
		Capacity         & Twelve layers, width 768; $384\times384$ in & Four layers, width 96; 32 tokens \\
		Positions        & Learned spatial embeddings                  & Fixed sinusoidal temporal encoding \\
		Input flow       & Relative to onset                           & Relative to the preceding sampled frame \\
		\hline
	\end{tabular}
	\caption[Divergences from published SLSTT]{Six material divergences from published SLSTT. The component evaluated here is a small temporal transformer encoder, not a reproduction of the published architecture.}
	\label{tab:slstt-divergences}
\end{table}

Six matched pairs evaluate the encoder with other component choices fixed within each pair. They test the whole encoder --- positional encoding, feed-forward layers and normalisation included --- against direct mean pooling, and do not isolate self-attention from those accompanying operations or compare it against another learned temporal model. Three follow-ups follow from \autoref{tab:slstt-divergences}: testing an LSTM aggregator, comparing onset-referenced with consecutive-frame flow, and replacing fixed positions with learned embeddings. All three remain untested, and changing the input representation requires re-training before any benefit can be claimed. What this study can establish is a matched comparison of the implemented encoder against an order-blind baseline within this pipeline; it neither validates published SLSTT nor establishes a general ranking of attention, recurrence and convolution.


% =============================================================================
\section{Training Under Severe Class Imbalance}
\label{sec:imbalance-review}
% =============================================================================

The 99 : 32 : 25 distribution of \autoref{sec:label-taxonomy} is a property of what can be elicited in a laboratory, not a sampling accident, so it cannot be corrected by collecting more data from the same paradigm. The corpus corrects for it at three points --- the data, the loss and the metric --- and this thesis uses all three. The mechanisms are given in \autoref{sec:focal-loss} and \autoref{sec:balanced-sampling}.

\textbf{At the loss.} Zhao et al.~\cite{zhaoS2021} bring focal loss into micro-expression recognition ``to alleviate the inefficient model training caused by the class imbalance in the MEs Datasets'', adopting the formulation developed for one-stage object detection. The division of labour inside it is what matters: $\gamma$ re-weights by \textit{difficulty} and is treated as a constant, ``set as 2 in practice'', while $\alpha$ re-weights by \textit{class frequency} and is ``treated as a hyper-parameter to set by cross validation''. Conflating the two is the error discussed below. Zhao et al. compare focal loss against cross-entropy across a sweep of $\alpha$, but hold $\gamma$ fixed, so the focusing exponent remains unvalidated on this task.

\textbf{At the data.} Xia et al.~\cite{xia2020a} pair ``temporal data augmentation strategies as well as a balanced loss \ldots jointly'' to overcome ``limited and imbalanced training samples''. Their own ablation removes each correction in turn and reports the balanced loss improving performance only ``slightly'' --- at most about two points, and in one cell costing accuracy. This is the corpus's only controlled look at composing two corrections, and no reviewed paper reports the pairing failing outright.

\textbf{At the metric.} \autoref{sec:megc-metrics} defines the pooled macro-averaged measure. Liong et al.~\cite{liong2016} add the argument from the subject side, reporting F1, precision and recall because these ``provide a more meaningful perspective than accuracy rates when the datasets used are naturally imbalanced since each subject has a different number of video samples.'' Both arguments compound here, where classes are skewed \textbf{and} per-subject clip counts range from 1 to 33.

One further gap: \textbf{label smoothing is named in the corpus but never examined on this task.} Dosovitskiy et al.~\cite{dosovitskiy2021} list it among the regularisers they tune, but no reviewed micro-expression paper uses it or reports what it does under class skew.

\subsection{Implications for this research}
\label{sec:imbalance-double-correction}

Focal loss is adopted with \textbf{$\gamma = 2.0$}, matching Zhao et al.'s reported setting. The class-frequency term is implemented as an optional per-class $\alpha$ vector but \textbf{switched off at run time}, because frequency correction is assigned to a weighted sampler active in every configuration. Label smoothing of 0.05 is applied inside the focal loss, declared as a choice this corpus supports neither way. The primary metric is pooled macro F1, with accuracy quoted only alongside the always-Negative floor.

\textbf{What this project contributes is a negative result about composing corrections.} The configuration is resolved at run time by a single rule: class weighting in the loss is enabled only when the balanced sampler is \textit{not} active. That rule exists because of two failures in opposite directions, both recoverable from this project's own run history. In an earlier code era the loss carried inverse-frequency class weights and no sampler existed; under that single correction the proposed configuration abandoned the \textbf{Positive} class outright, a failure initially misdiagnosed as over-capacity overfitting. A later era added the sampler while leaving the loss weights unconditional, so both corrections ran at once; the collapse then inverted, and configurations abandoned the \textbf{Negative} majority class instead, one predicting a single class for every clip. With the guard in place, no configuration in the reported study abandons any class.

This is worth reporting because no reviewed paper warns that combining corrections can invert the intended effect rather than merely blunt it. It must be stated carefully, however: the two collapses occurred in code eras differing from the reported run in more than the correction rule, and the final sweep contains no matched arm with the guard disabled. \textbf{The guard's causal benefit is therefore not isolated}, and the evidence does not establish that correcting at exactly one point is universally best --- only that correcting twice, in this pipeline, produced a class collapse that correcting once did not.

A second declaration follows. Because every configuration shares the same sampler, loss and $\gamma$, this thesis measures \textbf{no} effect attributable to the imbalance treatment: it is a constant, not a factor. The claim supported is that the four architectural components were compared under a training regime that does not abandon minority classes, not that this regime is optimal.

\subsection{Relevance and scope}
\label{sec:imbalance-gap}

The corrections themselves are well established. What the corpus does not settle is how two corrections compose, what $\gamma$ should be on this task, and what over-correction looks like when it fails. This thesis contributes to the first and third by documenting two opposite class-collapse failures and the run-time rule that resolved them. On the second it offers no improvement, since $\gamma = 2.0$ and label smoothing of 0.05 are held constant across all twelve configurations; sweeping either is declared as further work.


% =============================================================================
\section{Synthesis and Research Gap}
\label{sec:synthesis-and-gap}
% =============================================================================

This section consolidates the preceding relevance statements into one research gap, and collects in one place every declared divergence between the system built here and the literature it draws on.

Read as a whole, the reviewed literature settles a good deal. The corpora are small, demographically narrow, laboratory-bound and irreducibly imbalanced by what can be elicited, and CASME II remains the primary single-corpus benchmark under its native leave-one-subject-out protocol. Motion rather than appearance is the appropriate input; a strain field adds a further cue, though on CASME II only weakly and under a protocol that is not subject-disjoint; and both the magnification factor and the sequence length have interior optima located only by sweep. At this data size, capability bought without parameters outperforms capability bought with them, and self-attention is the strongest-performing mechanism reported. Accuracy alone is inadequate under skew and must give way to macro-averaged measures accumulated across folds before averaging.

\subsection{Four limitations that recur across every topic}
\label{sec:four-limitations}

The per-section statements above are not independent observations. Four structural problems generate most of them.

\textbf{(a) Components are isolated one at a time, inside one proposed design, and never jointly.} Matched-pair measurement is not absent from the corpus: \autoref{tab:evm-effect} isolates magnification, \autoref{tab:strain-performance} isolates strain against a flow control, Xia et al.~\cite{xia2020b} compare a recurrent convolutional network against three variants each adding one module with every other parameter held constant, and SLSTT's Mean-versus-LSTM comparison isolates its aggregator (\autoref{sec:slstt-results}). What no reviewed study does is vary \textit{several} components together, so that their interactions can be read. Every ablation is internal to one proposed design and moves one factor from that design's own defaults; across papers, systems are compared whole, so a difference in the reported number cannot be attributed to any one stage. There is therefore no basis for deciding which learned stages are load-bearing, nor whether one component's value depends on another's presence.

\textbf{(b) The evaluation protocol is unstable and under-reported.} Sample counts, label-set size and the choice among leave-one-subject-out, leave-one-video-out and leave-one-sample-out vary inconsistently, and the field's most-quoted baseline is described three ways by its own sources (\autoref{sec:original-baseline}). Given this corpus's fold composition these are not bookkeeping differences: they change what the numbers mean.

\textbf{(c) The strongest results depend on regimes a single-corpus study cannot access.} The best CASME II figures come from composite-database training (\autoref{sec:megc-metrics}) and, for SLSTT, ImageNet pre-training --- neither a property of the architecture being credited. Almost no evidence exists about what these components do when trained from scratch on one small corpus, which is the situation most practitioners are in.

\textbf{(d) Preprocessing parameters have interior optima and are transferred without re-validation.} The magnification factor, the sequence length and the input resolution all rise then fall, and none has a principled selection rule. Values are carried between pipelines whose downstream models differ entirely: $\alpha = 30$ moves from a single-frame appearance encoder to other settings, ten-frame interpolation propagates from one baseline paper to the whole field, and $\lambda = 10^{-4}$ transfers from CIFAR though its own authors re-tuned it for ImageNet.

\subsection{The research gap}
\label{sec:research-gap}

Taking (a) to (d) together:

\textbf{The reviewed literature has established a standard pipeline for micro-expression recognition --- magnify, compute motion, normalise duration, extract spatially, model temporally --- and contains matched ablations of individual stages, but no study varies several of those stages jointly so that their relative contributions and interactions can be read. Because components are varied one at a time inside a single proposed system, under protocols that differ in unstated ways, and predominantly in training regimes richer than a single small corpus provides, it is not currently possible to say from this literature which parts of that pipeline earn their place, which are inherited convention, and which depend on the presence of another.}

This matters practically as well as scientifically. A practitioner assembling this pipeline today has no published basis for allocating effort: for judging whether the most computationally expensive stage is the most valuable, whether a component with a negligible average effect may still matter in a specific failure, or whether a ranking read off an incomplete evaluation of a corpus this size is stable at all.

\subsection{How this thesis addresses the gap, and what it does not}
\label{sec:addressing-gap}

\textbf{What it does.} This thesis treats the pipeline as a factorial experiment rather than a proposal: four components varied independently across every architecturally valid cell of a 2\textsuperscript{4} matrix, twelve configurations, each evaluated under \textbf{complete 25-fold leave-one-subject-out on CASME II} with every non-varied factor held identical. That supplies the joint measurement (a) identifies as missing, reporting the protocol in full addresses (b) for this study at least, and training from scratch is the regime (c) identifies as unevidenced. Each component's marginal contribution is read from pairs of runs differing in one factor alone, and the distribution of those four effects is the substantive answer to the gap. Two further contributions follow from the review rather than the experiment: a \textbf{repaired measurement}, where an early data-routing defect left the magnification switch inert and the review's evidence of a clear, repeatable magnification effect is what made bit-identical results detectable as a defect rather than a null (\autoref{sec:evm-implications}); and \textbf{two opposite class-collapse failures} that no reviewed paper reports (\autoref{sec:imbalance-double-correction}).

\textbf{What it does not do.} Three of the four limitations are only partially addressed. On (b), this study reports its own protocol completely but cannot repair the field's inconsistency, and its own N = 156 differs from MEGC's 145. On (c), it establishes what these components do from scratch on one corpus, which is a different question from --- not an answer to --- how they behave under composite training. On (d), the thesis inherits parameter values rather than sweeping them: $\alpha = 10$, T = 32, $\lambda = 10^{-4}$ and $\gamma = 2.0$ are fixed throughout, so the ablation measures architecture at one point in preprocessing space rather than across it.

\subsection{A consolidated ledger of declared divergences}
\label{sec:divergence-ledger}

Reviewing the literature closely enough to write this chapter surfaced twelve points at which the implemented system departs from the work it draws on. Several were discovered only through the review. They are collected in \autoref{tab:divergence-ledger} so that no result rests on an unstated assumption.

\begin{table}[htbp]
	\centering
	\footnotesize
	\begin{tabular}{@{}p{0.85cm} p{6.0cm} p{6.6cm}@{}}
		\hline
		\textbf{\S} & \textbf{Divergence} & \textbf{Consequence} \\
		\hline
		\ref{sec:label-taxonomy} & Working set is N = 156, including sadness and fear in Negative; MEGC's CASME II subset is N = 145 & Comparisons against challenge figures are not exactly like-for-like \\
		\ref{sec:casme2-construction} & The corpus's own registered release is not used; the pipeline reads the original recorded frames, unregistered and uncropped & No landmarking or alignment of any kind precedes the network; tolerance to pose and scale must be learned from 156 clips \\
		\ref{sec:evm-implications} & Temporal band is the unnarrowed 5--25 Hz duration rule, where Bai et al. narrowed it to 15--25 Hz as noisy & The pass-band admits the low-frequency range they rejected \\
		\ref{sec:evm-implications} & Band-pass is realised by an FFT coefficient mask, not the Butterworth filter of the reviewed description & A different filter shape, with different transition behaviour at the band edges \\
		\ref{sec:strain-divergence} & Strain magnitude counts shear once, not twice as in the published Frobenius norm & Shear weighted lower by $\sqrt{2}$ relative to normal strain; untested \\
		\ref{sec:strain-implications} & No dedicated strain filtering, where the literature applies Wiener and Gaussian filters & Spatial differentiation can amplify flow-estimation noise, and nothing downstream suppresses it \\
		\ref{sec:tim-implications} & Temporal normalisation is uniform index sampling, \textbf{not} manifold-based TIM; no frames are synthesised & Avoids fabricating motion; forfeits the ability to lengthen short clips \\
		\ref{sec:tim-implications} & Magnification is applied \textbf{after} subsampling, at a nominal 200 fps & Realised pass-band is clip-dependent and below the intended 5--25 Hz \\
		\ref{sec:cnn-implications} & Backbone kernels are spatial only --- \textbf{no temporal mixing} & The ablation tests a learned \textit{spatial} stem, not spatio-temporal convolution \\
		\ref{sec:cnn-implications} & Backbone input is 224 $\times$ 224, against Xia et al.'s guidance of $\leq$ 100 $\times$ 100 --- though Bai et al. and Li, Huang and Zhao both use 224 $\times$ 224 & Whatever the stem is measured to contribute may be a resolution effect as much as an architecture effect \\
		\ref{sec:simam-implications} & SimAM statistics pooled over the whole clip rather than per frame as in the source; $\lambda$ not re-tuned & Neurons judged against the clip, not the frame; defensible but untested \\
		\ref{sec:transformer-implications} & Six divergences from published SLSTT (temporal not spatial attention, mean not LSTM aggregation, from scratch, smaller, sinusoidal encoding, short-term flow) & The temporal encoder is \textbf{not} a reproduction of SLSTT \\
		\ref{sec:imbalance-double-correction} & Label smoothing of 0.05, named in the corpus only by Dosovitskiy et al. and never examined on this task & A choice the micro-expression corpus supports neither way, held constant and unmeasured \\
		\hline
	\end{tabular}
	\caption[Every declared departure from the reviewed literature]{Every declared departure from the reviewed literature, with its consequence. The first column resolves to the subsection that owns each divergence.}
	\label{tab:divergence-ledger}
\end{table}

Three of these --- the strain formula, the magnification ordering and the input resolution --- were identified by the review itself rather than being known design decisions, and each is declared as proposed further work. None invalidates the ablation, since every configuration shares the same choices. What they constrain is the \textit{wording} of the conclusions: what is measured is the contribution of a spatial stem at 224 $\times$ 224 and of temporal self-attention in this pipeline, not a verdict on convolution over time or on SLSTT. The limits of what can be claimed are therefore fixed before any number is reported --- a small, single-seed, demographically narrow corpus measured against richer published regimes, with these divergences standing throughout.
